\section{Derivations for Section \ref{sec:simple}}

\subsection{The bond demand view}

Define the normalized bond demand $b_{j} \equiv \frac{R_f^{-1} B_j}{W_j} = 1 -c_{j}(r,\pi) - \alpha_j (1-c_{j}(r,\pi))$ as the bond-to-wealth ratio:
\begin{equation}
    b_j = (1-\alpha_j) (1-c_j(r,\pi)) = 1-c_{j}(r,\pi) - \frac{P Q_j}{W_j}.
\end{equation}

Notice that the bond demand depends not only on the interest rate, but also on the price of the risky asset, that is, the cross-elasticity will generically be non-zero.

The market clearing condition for bonds can be expressed as follows
\begin{equation}
    \underbrace{x_a b_a}_{\textbf{active bond demand}} = \underbrace{- x_p b_p}_{\textbf{net bond supply}},
\end{equation}
where $x_j \equiv \frac{\omega_j W_{j}}{\omega_a W_a + \omega_p W_p} = \omega_j Q_{j,-1} $ denotes the wealth share of investor $j$.

Notice that $\frac{P Q_j}{W_j} = \frac{P}{Y+P} \frac{\omega_j Q_j}{x_j}$. 
We can then write the bond market-clearing condition as follows:
\begin{equation}
    1 - \overline{c}(\mu-p) -\frac{e^p}{1+e^p} \left[\omega_a Q_a + \omega_p Q_p\right] = 0 \iff
    \underbrace{\frac{1}{1+e^p} - \overline{c}(\mu-p)}_{\text{excess supply of goods}} -\frac{e^p}{1+e^p} \left[\omega_a Q_a + \omega_p Q_p-1\right] = 0
\end{equation}

The left-hand side of the two expressions above gives you total bond demand. 
In the partial equilibrium case, the risk premium adjusts to clear the market for risky assets, so $\omega_a Q_a + \omega_p Q_p=1$, keeping the interest rate fixed.
But, due to the cross-elasticity in bond demand, a change in the price of the risky asset --- keeping the interest rate fixed --- would create a disequilibrium in the bond market.
For instance, in the case $\chi_p \equiv - \overline{c}'(\mu-p) > 0$, a decrease in the price of the risky asset would create an insufficient demand for goods and an excess demand for bonds.
The excess demand for bonds would create an incentive to reduce the interest rate, 

% The $\hat{b}_j \equiv b_j - b_j^*$. Then, the linearized bond demand for the passive investor is given by
% \begin{equation}
%     \hat{b}_p = -(1-\alpha_p^*) \left[\frac{\partial c_p^*}{\partial r} \hat{r} + \frac{\partial c_p^*}{\partial \pi} \hat{\pi} \right] - (1-c_p^*) \overline{\alpha}_p^* \hat{\alpha}_p. 
% \end{equation}

% The linearized bond demand for the active investor is given by
% \begin{equation}
%     \hat{b}_a = -(1-\alpha_a^*) \left[\frac{\partial c_a^*}{\partial r} \hat{r} + \frac{\partial c_a^*}{\partial \pi} \hat{\pi} \right] - (1-c_a^*) \frac{\partial \alpha_a^*}{\partial \pi} \hat{\pi}.
% \end{equation}

% The risk premium is given by
% \begin{equation}
%     \hat{\pi} = - \hat{p} - \hat{r}.
% \end{equation}

% We can then write the bond demand as follows:
% \begin{equation}
%     \hat{b}_j = - \zeta_{}
% \end{equation}

% The linearized bond demand for a passive investor is given by
% \begin{equation}
%     b_j - b_j^* = \left[ r -r^* +\frac{c_j'(\mu-p^*)}{1-c_j(\mu-p^*)}   (p-p^*)\right] b_j^* - \alpha_j^* e^{r^*} (1-c_j(\mu-p^*))  \hat{\alpha}_j
% \end{equation}

% For simplicity, focus on the case $\overline{\alpha}_p = 1$, so $b^*_p = b_a^* = 0$.
% The passive bond demand is then given by
% \begin{equation}
%     x_p b_p=  f^b,
% \end{equation}
% where $f^b \equiv  -x_p [1-c_p(\mu-p^*)] \hat{\alpha}_p$.
% The active bond demand is given by
% \begin{equation}
%     x_a b_a  = -\zeta^b_p (p-p^*) - \zeta^b_r (r-r^*),
% \end{equation}
% where $\zeta^b_p = \zeta^b_r = -x_a g_a'(\mu-p^*-r^*)[1-c_a(\mu-p^*)]$.

% The demand system can be written as follows:
% \begin{equation}
%     \left[
%     \begin{matrix}
%         \zeta_p^q & \zeta^q_r\\
%         \zeta_p^b & \zeta^b_r 
%     \end{matrix}
%     \right] \left[ 
%     \begin{matrix}
%         p-p^*\\
%         r-r^*
%     \end{matrix}
%     \right] = \left[ 
%     \begin{matrix}
%         f^q\\
%     f^b
%     \end{matrix}
%     \right],
% \end{equation}
% where we denote here $f^q \equiv f$ for symmetry.

% Inverting the system above, we obtain
% \begin{equation}
%     \left[ 
%     \begin{matrix}
%         p-p^*\\
%         r-r^*
%     \end{matrix}
%     \right] = 
%     \frac{1}{\zeta_{p}^q \zeta^b_r - \zeta_r^q \zeta ^b_p } \left[
%     \begin{matrix}
%         \zeta^b_r  & -\zeta^q_r\\
%         -\zeta_p^b & \zeta_p^q
%     \end{matrix}
%     \right] 
%     \left[ 
%     \begin{matrix}
%         f^q\\
%     f^b
%     \end{matrix}
%     \right].
% \end{equation}

% The price is given by
% \begin{equation}
%     p-p^* = \frac{\zeta_r^b f^q - \zeta_r^q f^b}{\zeta_{p}^q \zeta^b_r - \zeta_r^q \zeta ^b_p } = 
%      x_p \hat{\alpha}_p\frac{\zeta_r^b  + \zeta_r^q (1-c_p(\mu-p^*))}{\zeta_{p}^q \zeta^b_r - \zeta_r^q \zeta ^b_p } = 0,
% \end{equation}
% using the fact that $\zeta^q_r = x_a g_a'(\mu-p^*-r^*)$ and $c_a(\mu-p^*) = c_p(\mu-p^*)$.

\section{Derivations for Section \ref{sec:model}}

\subsection{Investors' problem}

The Hamilton-Jacobi-Bellman (HJB) equation for investor $j$ can be written as
\begin{align}
    0 &= \max_{C_j, \alpha_j} f_j(C_j,V_j) + V_{j,W}\left[ r W_j + \pi^\top \alpha_j W_j  - C_j\right] +V_{j,X} \mu_X\nonumber\\
    &\qquad + \frac{1}{2} V_{j,WW} W_j^2 \|\sigma_j\|^2 + V_{j,WX}W_j \sigma_X \sigma_j^\top    + \frac{1}{2} \sum_{k=1}^3\ \sigma_{X,k}^\top V_{j,XX} \sigma_{X,k},
\end{align}
subject to 
\begin{equation}
    \alpha_j = \overline{\alpha}_p[ \alpha_{D,t}^m, \alpha_{B,t}^m, 0]^\top \qquad \text{for } j = 0, \qquad \qquad 
    \sigma_j \sigma_j^\top \leq \overline{\sigma}^2 \qquad \text{for } j = 2.
\end{equation}
given $\sigma_{j} = \alpha_j^\top \Sigma$, where $\alpha_{k,t}^m \equiv \frac{P_{k,t}}{P_{Y,t}}$ for $k \in \{D,B\}$.
For ease of notation, we dropped time subscripts.
Note that  $V_{j,X}$ and $V_{j,WX}$ are $1\times N$ vectors, $V_{j,XX}$ is a $N\times N$ matrix, and both $V_{j,W}$ and $V_{j,WW}$ are scalars.
The drift $\mu_X$ is a $N\times 1$ vector, the diffusion $\sigma_X$ is a $N\times 3$ matrix, while $\sigma_R$ is a $1\times 3$ vector.
The notation $\sigma_{X,k}$ denotes the $k-$th column of $\sigma_{X}$, that is, the exposure to the $k-$th Brownian motion.



The optimality condition for consumption and portfolio share are given by
\begin{align}
    f_{j,C}(C_j,V_j) &= V_{j,W}\\
    V_{j,W} \pi &= - V_{j,WW} W_j \Sigma \Sigma^\top \alpha_j -  \Sigma (V_{j,WX} \sigma_X)^\top + \tilde{\nu}_{j,t} \Sigma \Sigma^\top \alpha_{j}.
\end{align}
where the second optimality condition holds for $j = 1,2$, and $\tilde{\nu}_{j,t}$ is the Lagrange multiplier on the leverage constraint. 

The optimal consumption is given by
\begin{equation}
 C_j = \rho^\psi ((1-\gamma_j) V_j)^{\frac{1-\gamma_j\psi }{1-\gamma_j}}V_{j,W}^{-\psi}.
\end{equation}

The optimal portfolio share for an active investor is given by
\begin{equation}
\sigma_j =  -\frac{V_{j,W}}{V_{j,WW} W_j + \tilde{\nu}_{j,t}}  \eta_j   - \frac{V_{j,WX}}{V_{j,WW} W_j + \tilde{\nu}_{j,t}} \sigma_X.
\end{equation}
using the fact that $\pi = \Sigma \eta^\top$.

Given the homotheticity of preferences, the value function for investor $j$ can be written as
\begin{equation}\label{app:value_function}
    V_{j,t} = \left(\frac{c_{j,t}}{\rho^\psi}\right)^{\frac{1-\gamma_j}{1-\psi}} \frac{W_{j,t}^{1-\gamma_j}}{1-\gamma_j}.
\end{equation}

This particular parametrization of the value function implies that the consumption-wealth ratio is given by
\begin{equation}
    \frac{C_{j,t}}{W_{j,t}} = c_{j,t}.
\end{equation}

The optimal risk exposure for active investors is given by
\begin{equation}
    \sigma_{j,t}
    =\frac{1}{1+\nu_{j,t}} \left[ \frac{\eta_t}{\gamma_j}  + \frac{1-\gamma_j^{-1}}{\psi-1} \sigma_{c_j,t} \right],
\end{equation}
where $\nu_{j,t} \equiv \tilde{\nu}_{j,t} / (V_{j,WW,t} W_{j,t})$ is the normalized Lagrange multiplier on the leverage constraint.

Plugging the consumption-wealth ratio into the HJB equation and using Equation \eqref{app:value_function}, we obtain
\begin{align}
        0 &=  \frac{\rho   }{1-\psi^{-1}} \left[  \rho^{-1} c_j-1 \right] + r  + \eta \sigma_j^ \top   - c_j + \frac{1}{1-\psi} \left[ \frac{\xi_{j,X}}{\xi_j} \mu_X + \frac{1}{2} \sum_{k=1}^d \sigma_{X,k}^\top \frac{\xi_{j,XX}}{\xi_j} \sigma_{X,k}\right]\nonumber\\
    &\qquad - \frac{\gamma_j}{2} \| \sigma_j\|^2  + \frac{1-\gamma_j}{1-\psi} \frac{c_{j,X}}{c_j} \sigma_X  \sigma_j^ \top    + \frac{1}{2} \frac{\psi-\gamma_j}{(1-\psi)^2} \sum_{k=1}^d \sigma_{X,k}^\top \frac{c_{j,X}^\top}{c_j} \frac{c_{j,X}}{c_j} \sigma_{X,k}.
\end{align}

Rearranging the expression above, we obtain
\begin{small}
\begin{equation}
    c_{j,t} = \psi \rho + (1-\psi) \left[ r_t + \eta_t \sigma_{j,t}^\top - \frac{\gamma_j}{2} \| \sigma_{j,t}\|^2 \right] + \mu_{c_j,t} + (1-\gamma_j)  \sigma_{c_j,t} \sigma_{j,t}^\top + \frac{\psi - \gamma_j}{1-\psi} \frac{\|\sigma_{c_j,t} \|^2}{2},
\end{equation}
\end{small}
where the law of motion of $c_{j,t}$ is given by
\begin{equation}
    \frac{d c_{j,t}}{c_{j,t}} = \mu_{c_j,t} dt + \sigma_{c_j,t} d Z_t,
\end{equation}
and the drift and diffusion of $c_{j,t}$ are given by Ito's lemma:
\begin{align}
    \mu_{c_j,t} &= \frac{c_{j,X}}{c_j} \mu_{X,t} + \frac{1}{2} \sum_{k=1}^d \sigma_{X,k,t}^\top \frac{c_{j,XX,t}}{c_{j,t}} \sigma_{X,k,t}, \qquad \qquad 
    \sigma_{c_j,t} = \frac{c_{j,X}}{c_j} \sigma_{X,t}.
\end{align}


\subsection{Pricing condition}

Let $p_{D,t} = \frac{P_{D,t}}{Y_t}$ and $p_{B,t} = \frac{P_{B,t}}{Y_t}$ denote the normalized prices of the dividend and bond claims, respectively.
The normalized price of the aggregate claim is $p_{Y,t} = \frac{P_{Y,t}}{Y_t}$.
The pricing condition for asset $k \in\{D,B\}$ is given by
\begin{equation}
    r_t + \sigma_{R_k,t} \eta_t^\top = \frac{1}{p_{k,t}} + \mu + \mu_{p_k,t} + \sigma \sigma_{p_k,t}^\top,
\end{equation}
where $\sigma_{R,t} = \sigma + \sigma_{p_k,t}$ and $(\mu_{p_k,t}, \sigma_{p_k,t})$ are given by Ito's lemma:
\begin{equation}
    \mu_{p_k,t} = \frac{p_{k,X,t}}{p_{k,t}} \mu_{X,t} + \frac{1}{2}\sum_{k=1}^d \sigma_{X,k,t}^\top\,\frac{p_{k,XX,t}}{p_{k,t}}\,\sigma_{X,k,t}, \qquad \qquad 
    \sigma_{p_k,t} = \frac{p_{k,X,t}}{p_{k,t}} \sigma_{X,t}.
\end{equation}

The expected return and risk exposure on the aggregate claim are given by
 $\mu_{R_Y,t} = \alpha_{D,t}^m \mu_{R_D,t} + \alpha_{B,t}^m \mu_{R_B,t}$ and $\sigma_{R_Y,t} = \alpha_{D,t}^m \sigma_{R_D,t} + \alpha_{B,t}^m \sigma_{R_B,t}$, respectively.

% Let $y_t \equiv Y_t/P_t$ denote the dividend yield on the risky asset.
% From Equation~\eqref{eq:risky-asset-ret}, we can write the expected return on the risky asset as:
% \begin{equation}
%     r_t + \theta_t \| \sigma_{R,t}\| = y_t + \frac{1}{dt} \frac{d(Y_t/y_t)}{(Y_t/y_t)} = y_t + \mu - \mu_{y,t} + \|\sigma_{y,t}\|^2 - \sigma \sigma_{y,t}^\top.
% \end{equation}

% Rearranging the expression above, we obtain
% \begin{equation}\label{app:pricing_cond}
%     y_t = r_t + \theta_t \|\sigma_{R,t}\| -\mu + \mu_{y,t} - \|\sigma_{R,t}\|^2+ \sigma \sigma_{R,t}^\top,
% \end{equation}
% where $\sigma_{R,t} = \sigma - \sigma_{y,t}$ and $(\mu_{y,t}, \sigma_{y,t})$ are given by Ito's lemma:
% \begin{equation}
%     \mu_{y,t} = y_{X,t} \mu_{X,t} + \frac{1}{2}\sum_{k=1}^d \sigma_{X,k,t}^\top\,y_{XX,t}\,\sigma_{X,k,t}, \qquad \qquad 
%     \sigma_{y,t} = y_{X,t} \sigma_{X,t}.
% \end{equation}

\subsection{Aggregate state variable}

Define the share of wealth of type-$j$ investors as follows
\begin{equation}
    x_{j,t} \equiv \frac{\omega_j W_{j,t}}{\sum_{i=0}^2 \omega_i W_{i,t}}.
\end{equation}


We define the aggregate state variable as $X_t = (x_{1,t}, x_{2,t}, \overline{\alpha}_{p,t}, s_t)$.

\paragraph{Market clearing conditions.}
We can write the market clearing condition for the equity claim, corporate bond claim, and short-term bonds as follows:
\begin{equation}
    \sum_{j=0}^2 \omega_j W_{j,t} \alpha_{j,k,t} = P_{k,t}, \qquad 
    \sum_{j=0}^2 \omega_j W_{j,t} (1-\alpha_{j,D,t}-\alpha_{j,B,t}) = 0,
\end{equation}
for $k \in \{D,B\}$.

Using the fact that $P_{Y,t} = P_{D,t} + P_{B,t}$, we obtain:
\begin{equation}
    \sum_{j=0}^2 \omega_j W_{j,t} = P_{Y,t}.
\end{equation}

The market clearing condition for goods is given by
\begin{equation}
    \sum_{j=0}^2 x_{j,t} c_{j,t} = \frac{1}{p_{Y,t}}.
\end{equation}

We can then write the market clearing condition for the equity claim, corporate bond claim, and the derivative as follows:
\begin{equation}
    \sum_{j=0}^2 x_{j,t} \alpha_{j,D,t} = \alpha_{D,t}^m, \qquad 
    \sum_{j=0}^2 x_{j,t} \alpha_{j,B,t} = \alpha_{B,t}^m, \qquad 
    \sum_{j=0}^2 x_{j,t} \alpha_{j,F,t} = 0.
\end{equation}  

Let $\alpha_t^m = [\alpha_{D,t}^m, \alpha_{B,t}^m,0]^\top$ denote the market weights of the dividend, bond, and derivative claims.
We can then write the market clearing conditions above in a compact way:
\begin{equation}
    \sum_{j=0}^2 x_{j,t} \alpha_{j,t} = \alpha_t^m.
\end{equation}
Multiplying both sides by $\Sigma$ and using the fact that $\sigma_{R_Y,t} = \omega_{D,t} \sigma_{R_D,t} + \omega_{B,t} \sigma_{R_B,t}$, we obtain
\begin{equation}
    \sum_{j=0}^2 x_{j,t} \sigma_{j,t} = \sigma_{R_Y,t}.
\end{equation}

\paragraph{Law of motion of the aggregate state variable.}
The law of motion of $\overline{\alpha}_{p,t}$ is given by \eqref{eq:alpha_p_process}.
The law of motion of $s_t$ is given by \eqref{eq:s_process}.
To compute the law of motion of $x_{j,t}$, first, note that the law of motion of wealth for a type-$j$ investor can be written as
\begin{equation}
    \frac{d W_{j,t}}{W_{j,t}} = \left[ r_t + \eta_t \sigma_{j,t}^\top - c_{j,t} \right]dt + \sigma_{j,t} d Z_t.
\end{equation}

From Ito's lemma, the law of motion of $x_{j,t}$ is given by
\begin{align*}
    \frac{dx_{j,t}}{x_{j,t}} 
    &= \Big[ r_t + \pi_t^\top \alpha_{j,t} - c_{j,t} - \mu_{P_Y,t} 
        - (\sigma_{j,t} - \sigma_{R_Y,t}) \, \sigma_{R_Y,t}^\top 
        + \kappa \frac{\omega_j - x_{j,t}}{x_{j,t}} \Big] dt  
        + (\sigma_{j,t} - \sigma_{R_Y,t}) \, dZ_t.
\end{align*}
using $\mu_{P_Y,t} = \mu_{R_Y,t} - \mu_{R_D,t} \alpha_{D,t}^m - \mu_{R_B,t} \alpha_{B,t}^m$.


\subsection{Block recursivity and the mutual fund theorem}

To compute the equilibrium, one needs to solve a system of $5$ partial differential equations (PDEs), involving the three consumption-wealth ratios $\{c_{j,t}(X_t)\}_{j=0}^2$ and the two price-dividend ratios $\{p_{k,t}(X_t)\}_{k=D,B}$. %$\xi_{j}(x,\overline{\alpha}_p)$
These functions depend on $4$ state variables, the $2-$dimensional vector $x_{t}$, the portfolio share of passive investors $\overline{\alpha}_{p,t}$, and the dividend share $s_t$.

The system of PDEs satisfies an important \textit{block-recursive} structure.
We can first solve for a reduced system of PDEs involving $(c_j,p_Y)$ in the set of states $\hat{X}_t = (x_{1,t}, x_{2,t}, \overline{\alpha}_{p,t})$.
Given the solution of the reduced system of PDEs, we can then solve for $p_{D,t}$ and $p_{B,t}$ in the full set of states $X_t = (x_{1,t}, x_{2,t}, \overline{\alpha}_{p,t}, s_t)$.

\subsubsection{Block recursivity}

We next derive the system of PDEs. 
We will show that one can solve for a reduced system of PDEs involving $(c_j,p_Y)$.
% This requires expressing all equilibrium quantities as functions of the aggregate state variables $X_t$ and the functions $(c_j,p_Y)$ and their derivatives.
This is the crucial in the obtaining the mutual fund theorem in our setting.

\begin{lemma}\label{lemma:block_recursivity}
    The consumption-wealth ratios, $\{c_{j,t}(X_t)\}_{j=0}^2$, and the price-dividend ratio for the aggregate claim, $p_Y(X_t)$, satisfy the following system of PDEs:
    \begin{align}
        c_{j,t} &= \psi \rho 
        + (1-\psi) \left[ r_t + \eta_{t}\,\sigma_{j,t}^\top - \frac{\gamma_j}{2} \|\sigma_{j,t}\|^2 \right] 
        +  \mu_{c_j,t} 
        + (1-\gamma_j) \sigma_{c_j,t} \sigma_{j,t}^\top 
        + \frac{\psi - \gamma_j}{1-\psi} \frac{\|\sigma_{c_j,t}\|^2}{2}\\
        \frac{1}{p_{Y,t}} &= r_t + \eta_{t}\,(\sigma + \sigma_{p_Y,t})^\top - \big(\mu + \mu_{p_Y,t} + \sigma \, \sigma_{p_Y,t}^\top\big),        
\end{align}
where $(r_t, \eta_t, \sigma_{j,t}, \mu_{c_j,t}, \sigma_{c_j,t}, \mu_{p_Y,t}, \sigma_{p_Y,t})$ are all functions of $X_t$, $(c_j,p_Y)$ and their derivatives.    
\end{lemma}

\begin{proof}
    We proceed in four steps. 
    First, we show how to solve for the diffusion of the endogenous state variables, $\sigma_{x_1,t}$ and $\sigma_{x_2,t}$, in terms of $X_t$, $(c_j,p_Y)$ and their derivatives.
    Second, we show how to solve for the drift of the endogenous state variables, $\mu_{x_1,t}$ and $\mu_{x_2,t}$, again in terms of $X_t$, $(c_j,p_Y)$ and their derivatives.
    Third, we show how to solve for the interest rate, $r_t$.
    Finally, we show how to obtain the PDEs for $(c_j,p_Y)$.

    In the derivations below, we drop the explicit dependence on $X$ for brevity.

    \paragraph{Step 1: Risk exposure of endogenous state variables.}
    Consider the expression for the risk exposure of the aggregate claim and the consumption-wealth ratio:
    \begin{align}
        \sigma_{R_Y} &= \sigma + \sum_{k=1}^2 \frac{p_{Y,x_k}}{p_{Y}} \sigma_{x_k} + \frac{p_{Y,\overline{\alpha}_p}}{p_{Y}} \sigma_{\overline{\alpha}_p} + \frac{p_{Y,s}}{p_{Y}} \sigma_{s},\\
        \sigma_{c_j} &= \sum_{k=1}^2 \frac{c_{j,x_k}}{c_{j}} \sigma_{x_k} + \frac{c_{j,\overline{\alpha}_p}}{c_{j}} \sigma_{\overline{\alpha}_p} + \frac{c_{j,s}}{c_{j}} \sigma_{s},
    \end{align}
    
    Notice that we can write $\sigma_{R_Y}$ as follows:
    \begin{equation}
        \sigma_{R_Y} = \chi_{\sigma_{R_Y},0} + \chi_{\sigma_{R_Y},1} \sigma_{x_1} + \chi_{\sigma_{R_Y},2} \sigma_{x_2},
    \end{equation}
    where the coefficients $\chi_{\sigma_{R_Y},k}$, $k=0,1,2$, are functions of $X$ given by 
    \begin{align}
        \chi_{\sigma_{R_Y},0} &= \sigma  +\frac{p_{Y,\overline{\alpha}_p}}{p_{Y}} \sigma_{\overline{\alpha}_p} + \frac{p_{Y,s}}{p_{Y}} \sigma_{s} ,\\
        \chi_{\sigma_{R_Y},k} &= \frac{p_{Y,x_k}}{p_{Y}}, k=1,2.
    \end{align}

    Similarly, we can write $\sigma_{c_j}$ as follows:
    \begin{equation}
        \sigma_{c_j} = \chi_{\sigma_{c_j},0} + \chi_{\sigma_{c_j},1} \sigma_{x_1} + \chi_{\sigma_{c_j},2} \sigma_{x_2},
    \end{equation}
    where the coefficients $\chi_{\sigma_{c_j},k}$, $k=0,1,2$, are functions of $X$ given by 
    \begin{align}
        \chi_{\sigma_{c_j},0} &= \frac{c_{j,\overline{\alpha}_p}}{c_{j}} \sigma_{\overline{\alpha}_p} + \frac{c_{j,s}}{c_{j}} \sigma_{s}\\
        \chi_{\sigma_{c_j},k} &= \frac{c_{j,x_k}}{c_{j}}, k=1,2.
    \end{align}

    The market price of risk 
    \begin{align}
        \eta &= \gamma_{a} \left[\chi_{a}\sigma_{R_Y} - \sum_{j=0}^2 x_{j} \frac{1-\gamma_j^{-1}}{\psi-1} \, \frac{\sigma_{c_j}}{1+\nu_{j}}\right],
    \end{align}
    where $\chi_{a} =\frac{1 - (1-x_{1}-x_{2}) \overline{\alpha}_{p}}{x_{1}+x_{2}}$ is a known function of the aggregate state variables $X$.
    
    We can then write $\eta$ as follows:
    \begin{equation}
        \eta = \chi_{\eta,0} + \chi_{\eta,1} \sigma_{x_1} + \chi_{\eta,2} \sigma_{x_2},
    \end{equation}
    where the coefficients $\chi_{\eta,k}$, $k=0,1,2$, are functions of $X$ given by 
    \begin{align}
        \chi_{\eta,k} &= \gamma_{a} \left[ \chi_{a}\chi_{\sigma_{R_Y},k} - \sum_{j=0}^2 x_{j} \frac{1-\gamma_j^{-1}}{\psi-1} \, \frac{\chi_{\sigma_{c_j},k}}{1+\nu_{j}}\right].
    \end{align}


    The risk exposure of investor $j = 0,1,2$ is given by
    \begin{align}
        \sigma_{0} &= \overline{\alpha}_{p} \sigma_{R_Y}\\
        \sigma_{j} &= \frac{1}{1+\nu_{j}} \left[ \frac{\eta}{\gamma_j} + \frac{1-\gamma_j^{-1}}{\psi-1} \, \sigma_{c_j} \right], \qquad j=1,2.
    \end{align}

    
    We can then solve for $\sigma_{j}$ in terms of the risk exposure of the endogenous state variables:
    \begin{equation}
        \sigma_{j} = \chi_{\sigma_j,0} + \chi_{\sigma_j,1} \sigma_{x_1} + \chi_{\sigma_j,2} \sigma_{x_2},
    \end{equation}
    where the coefficients $\chi_{\sigma_j,k}$, $k=0,1,2$, are functions of $X$ given by 
    \begin{align}
        \chi_{\sigma_0,k} &= \overline{\alpha}_{p} \chi_{\sigma_{R_Y},k}\\
        \chi_{\sigma_j,k} &= \frac{1}{1+\nu_{j}} \left[ \frac{\chi_{\eta,k}}{\gamma_j} + \frac{1-\gamma_j^{-1}}{\psi-1} \chi_{\sigma_{c_j},k} \right], \qquad j=1,2.
    \end{align}

    The risk exposure of the endogenous state variables is given by
    \begin{equation}
        \sigma_{x_j} = x_{j} \left[ \sigma_{j} -  \sum_{k=0}^2 x_{k} \sigma_{k}  \right],
    \end{equation}
    using the fact that $\sigma_{R_Y} = \sum_{k=0}^2 x_{k} \sigma_{k}$.
    
    This gives us a system of equations determining $\sigma_{x_j,t}$:\small
    \begin{equation}
            \begin{bmatrix}
                1 - x_{1} \sum_{k=0}^2 \left( \chi_{\sigma_1,1} - \chi_{\sigma_k,1} \right) x_{k} & - x_{1} \sum_{k=0}^2 \left( \chi_{\sigma_1,2} - \chi_{\sigma_k,2} \right) x_{k}\\
                - x_{2} \sum_{k=0}^2 \left( \chi_{\sigma_2,1} - \chi_{\sigma_k,1} \right) x_{k} & 1 - x_{2} \sum_{k=0}^2 \left( \chi_{\sigma_2,2} - \chi_{\sigma_k,2} \right) x_{k}
            \end{bmatrix}
            \begin{bmatrix}
                \sigma_{x_1}\\
                \sigma_{x_2}
            \end{bmatrix}
            =
            \begin{bmatrix}
              x_{1} \sum_{k=0}^2 \left( \chi_{\sigma_1,0} - \chi_{\sigma_k,0} \right) x_{k} \\
              x_{2} \sum_{k=0}^2 \left( \chi_{\sigma_2,0} - \chi_{\sigma_k,0} \right) x_{k}
            \end{bmatrix}.
    \end{equation}\normalsize
    
    Solving the system above, we obtain $\sigma_{x_j}$ as a function of $X$ and the functions $c_j$ and $p_{Y}$ and their derivatives.
    Given $\sigma_{x_j}$, we can solve for $\sigma_{j},$ $\eta$, $\sigma_{c_j}$, and $\sigma_{R_Y}$ using the expressions above. 
    The risk premium on the aggregate claim is given by $\pi_{Y} = \sigma_{R_Y} \eta^\top$.

    \paragraph{Step 2: Drift of endogenous state variables.} 
Given $(\sigma_j, \eta, \sigma_{R_Y})$, we can solve for the drift of the endogenous state variables:
\begin{align}
    \mu_{x_j} &= x_{j} \left[ 
        (\eta-\sigma_{R_Y}) \left( \sigma_{j} -\sigma_{R_Y} \right)^\top - \left( c_{j} - \frac{1}{p_{Y}}  \right)
     \right] + \kappa (\omega_j - x_{j}),
\end{align}

This gives us $\mu_{X}$ and $\sigma_{X}$ in terms of $X$, $(c_j,p_Y)$ and their derivatives.
We can then solve for the drift of $c_j$ and $p_{Y}$ using Ito's lemma:
\begin{align}
    \mu_{c_j} &= \frac{c_{j,X}}{c_{j}} \mu_{X} + \frac{1}{2}\sum_{k=1}^3 \sigma_{X,k}^\top\,\frac{c_{j,XX}}{c_{j}}\,\sigma_{X,k},\\
    \mu_{p_Y} &= \frac{p_{Y,X}}{p_{Y}} \mu_{X} + \frac{1}{2}\sum_{k=1}^3 \sigma_{X,k}^\top\,\frac{p_{Y,XX}}{p_{Y}}\,\sigma_{X,k}.
\end{align}

\paragraph{Step 3: Interest rate.}
Given the expressions derived in the previous steps, we can solve for the shifter of the consumption-wealth ratio $\xi_{j,t}$:
\begin{equation}
    \xi_{j} = \mu_{c_j} 
    + (1-\gamma_j) \sigma_{c_j} \sigma_{j}^\top 
    + \frac{\psi - \gamma_j}{1-\psi} \frac{\|\sigma_{c_j}\|^2}{2}.
\end{equation}

We can then solve for $\xi = \sum_{j=0}^2 x_{j} \xi_{j}$ and the interest rate:
\begin{align}
    r &= \rho 
    + \psi^{-1} \left(\mu_{p_Y} + \mu + \sigma \sigma_{p_Y}^\top + \xi \right)
    + (1-\psi^{-1}) \sum_{j=0}^2 x_{j} 
        \frac{\gamma_j }{2} \|\sigma_{j}\|^2
    - \pi_{Y}.
\end{align}

\paragraph{Step 4: PDEs for the consumption-wealth ratio and price-dividend ratio.}
From the expression for the consumption-wealth ratio and the pricing condition for the aggregate claim, we can write the PDE for $c_j$ and $p_Y$ as follows:
\begin{align}
        c_{j} &= \psi \rho 
        + (1-\psi) \left[ r + \eta\,\sigma_{j}^\top - \frac{\gamma_j}{2} \|\sigma_{j}\|^2 \right] 
        +  \mu_{c_j} 
        + (1-\gamma_j) \sigma_{c_j} \sigma_{j}^\top 
        + \frac{\psi - \gamma_j}{1-\psi} \frac{\|\sigma_{c_j}\|^2}{2}\\
        \frac{1}{p_{Y}} &= r + \eta\,\sigma_{R_Y}^\top - \big(\mu + \mu_{p_Y} + \sigma \, \sigma_{p_Y}^\top\big),        
\end{align}
Notice that $(r, \eta, \sigma_{j}, \mu_{c_j}, \sigma_{c_j}, \mu_{p_Y}, \sigma_{p_Y})$ are all functions of $X$, $(c_j,p_Y)$ and their derivatives.
Hence, these equations form a system of PDEs for $(c_j,p_Y)$.
\end{proof}

A corollary of the above proposition is that the consumption-wealth ratio and price-dividend ratio are independent of the dividend share $s$.
\begin{cor}
    The consumption-wealth ratio $c_j$ and price-dividend ratio $p_Y$ are a function of the reduced state vector $\hat{X} = (x_1,x_2,\overline{\alpha}_p)$.
\end{cor}
\begin{proof}
    The dividend share $s$ does not appear in the expressions for PDE system for $(c_j,p_Y)$.
    They only appear indirectly through $\mu_{s}$ and $\sigma_{s}$.
    If $(c_j,p_Y)$ are a function of $\hat{X}$, then their derivatives with respect to $s$ are zero, and $s$ completely drops out of the PDE system for $(c_j,p_Y)$.
\end{proof}

\subsubsection{The mutual fund theorem}

Lemma~\ref{lemma:block_recursivity} and its corollary show that the consumption-wealth ratio and price-dividend ratio of the aggregate claim are independent of the dividend share $s$.
This implies that investor's consumption or wealth choices are not exposed to idiosyncratic dividend-share shocks:
\begin{equation}
    \sigma_{j,t,3} = \sigma_{R_Y,t,3} = \sigma_{R_F,t,3} = 0,
\end{equation}
where the third component of the risk exposure vector is the dividend-share shock.

We can implement the risk exposure $\sigma_{j,t}$ by holding the market index fund and a derivative.
Let $\hat{\sigma}_{j,t} = [\sigma_{j,t,1}, \sigma_{j,t,2}]$, $\hat{\Sigma}_{t} = \left[\begin{matrix}
 \hat{\sigma}_{R_Y,t} \\
 \hat{\sigma}_{R_F,t}
\end{matrix}\right]$, and $\hat{\alpha}_{j,t} = [\alpha_{j,Y,t}, \alpha_{j,F,t}]^\top$, then we can define the portfolio shares of the market index fund and the derivative as:
\begin{equation}
    \hat{\sigma}_{j,t} = \hat{\alpha}_{j,t}^\top \hat{\Sigma}_{t} \Rightarrow \hat{\alpha}_{j,t}^\top = \hat{\sigma}_{j,t} \hat{\Sigma}_{t}^{-1}.
\end{equation}

We can then write the risk exposure $\sigma_{j,t}$ as:
\begin{equation}
    \sigma_{j,t} = \alpha_{j,Y,t} \sigma_{R_Y,t} + \alpha_{j,F,t} \sigma_{R_F,t}.
\end{equation}
Using the fact that $\sigma_{R_Y,t} = \frac{P_{D,t}}{P_{Y,t}} \sigma_{R_D,t} + \frac{P_{B,t}}{P_{Y,t}} \sigma_{R_B,t}$, we obtain:
\begin{equation}
    \sigma_{j,t} = \alpha_{j,Y,t} \frac{P_{D,t}}{P_{Y,t}} \sigma_{R_D,t} + \alpha_{j,Y,t} \frac{P_{B,t}}{P_{Y,t}} \sigma_{R_B,t} + \alpha_{j,F,t} \left( \frac{P_{B,t}}{P_{Y,t}} \sigma_{R_Y,t} + \frac{P_{D,t}}{P_{Y,t}} \sigma_{R_D,t} \right).
\end{equation}
Hence, this implies that the portfolio shares of dividend and bond claims are given by:
\begin{equation}
    \alpha_{j,D,t} = \alpha_{j,Y,t} \frac{P_{D,t}}{P_{Y,t}},
    \quad
    \alpha_{j,B,t} = \alpha_{j,Y,t} \frac{P_{B,t}}{P_{Y,t}},
\end{equation}
where $\alpha_{j,Y,t}$ is the portfolio share of the market index fund defined.

\paragraph{Pricing the individual claims.} The normalized price of the dividend claim satisfies the pricing condition:
\begin{equation}
    \frac{s_t}{p_{D,t}} = r_t + \sigma_{R_D,t} \eta_t^\top - \big(\mu + \mu_{p_D,t} + \sigma \, \sigma_{p_D,t}^\top\big),
\end{equation}
where
\begin{equation}
    \mu_{p_{D,t}} = \frac{p_{D,X,t}}{p_{D,t}} \mu_{X,t} + \frac{1}{2}\sum_{k=1}^3 \sigma_{X,k,t}^\top\,\frac{p_{D,XX,t}}{p_{D,t}}\,\sigma_{X,k,t}, \qquad \qquad 
    \sigma_{p_{D,t}} = \frac{p_{D,X,t}}{p_{D,t}} \sigma_{X,t}.
\end{equation}

% \paragraph{Mapping to \citet{MenzlySantosVeronesi2004}.} We can map the dividend share $s$ to notation in \citet{MenzlySantosVeronesi2004} as follows:
% \begin{equation}
%     \sigma_s = v^1- v^0, \qquad (1-\overline{s}) v^0 + \overline{s} v^1 = 0,
% \end{equation}
% which implies that
% \begin{equation}
%     v^0 + \overline{s} (v^1- v^0) = 0 \Rightarrow v^0 = - \overline{s} \sigma_s, \qquad v^1 = (1-\overline{s}) \sigma_s,
% \end{equation}
% \begin{equation}
%     v^1 - s^ 1 v^ 1 - (1-s^1)v^0 = v^1- v^0 - s^ 1 (v^1- v^0) = (1-s^1)(v^1- v^0)
% \end{equation}
% \begin{equation}
%     v^0 - s^ 1 v^ 1 - (1-s^1)v^0 = v^0- v^0 - s^ 1 (v^1- v^0)= -s^ 1 (v^1- v^0)
% \end{equation}

\section{Derivations for Section \ref{sec:elasticity}: Perturbation Solution}\label{app:derivations-perturbation-solution}

In this section, we discuss the derivations of the perturbation solution in detail. 
Consistent with the block-recursivity property in Proposition~\ref{prop:mft}, the consumption-wealth ratios and the price-dividend ratio of the aggregate claim are independent of the dividend share $s$.

For these variables, we can focus on the reduced set of states $\hat{X} = (x, \overline{\alpha}_{p})$, where $x = (x_1, x_2)$ is the vector of wealth shares of the active investors. 
To price the individual claims, we need to consider the full set of states $X = (x, \overline{\alpha}_{p}, s)$.

The perturbation is defined by the following parameter restrictions:
\begin{equation}
    \gamma_j = \gamma (1+\hat{\gamma}_j \epsilon), \qquad 
    \overline{\sigma} = \|\sigma\| (1+\hat{\sigma} \epsilon), \qquad 
    \kappa = \hat{\kappa} \epsilon, 
\end{equation}
where $\sum_{j=0}^2 \omega_j \hat{\gamma}_j = 0$.

We also make assumptions about the dynamics of the exogenous state variables.
First, the portfolio share of passive investors is written as $\overline{\alpha}_{p} = 1 + \hat{\alpha}_{p} \epsilon$, where $\hat{\alpha}_{p}$ follows the process
\begin{equation}
d \overline{\alpha}_{p} = \kappa_p (\overline{\alpha} - \overline{\alpha}_{p}) dt + \sigma_p \epsilon dZ_t,
\end{equation}
where $\overline{\alpha} = 1+\hat{\alpha} \epsilon$.
Equivalently, $d\hat{\alpha}_{p} = \kappa_p(\hat{\alpha}-\hat{\alpha}_{p})dt + \sigma_p dZ_t$.

Second, the dividend share is written as $s = \overline{s} + \hat{s} \epsilon$, where $\hat{s}$ follows the process induced by
\begin{equation}
    d s = \kappa_s ( \overline{s} - s )dt + s (1-s) \sigma_s \epsilon d Z_t.
\end{equation} 

The case $\epsilon = 0$ corresponds to the benchmark economy, in which there is no preference heterogeneity and the exogenous state variables are constant.
In particular, passive investors are fully invested in the risky asset and the dividend share is fixed at $s = \overline{s}$.


We expand the consumption-wealth ratio of investor $j\in\{0,1,2\}$ and the price-dividend ratio of the aggregate claim as follows:
\begin{align}
    c_{j}(x, \overline{\alpha}_p;\epsilon) &= c_{j,0}(x, \hat{\alpha}_p) + c_{j,1}(x, \hat{\alpha}_p) \epsilon+ c_{j,2}(x, \hat{\alpha}_p) \epsilon^2 + \mathcal{O}(\epsilon^3),\\
    p_{Y}(x, \overline{\alpha}_p;\epsilon) &= p_{Y,0}(x, \hat{\alpha}_p) + p_{Y,1}(x, \hat{\alpha}_p) \epsilon+ p_{Y,2}(x, \hat{\alpha}_p) \epsilon^2 + \mathcal{O}(\epsilon^3).
\end{align}

For each individual claim $k\in\{D,B\}$, we use the expansion
\begin{align}
    p_{k}(x, \overline{\alpha}_p, s;\epsilon) &= p_{k,0}(x, \hat{\alpha}_p, \hat{s}) + p_{k,1}(x, \hat{\alpha}_p, \hat{s}) \epsilon+ p_{k,2}(x, \hat{\alpha}_p, \hat{s}) \epsilon^2 + \mathcal{O}(\epsilon^3).
\end{align}

The following lemma summarizes some basic properties of the perturbation solution that will be useful for the derivations below.
\begin{lemma}[Basic properties of the perturbation solution]\label{lemma:basic_properties}
    The perturbation solution satisfies the following properties:
    \begin{enumerate}   
        \item The zeroth-order terms are constant:
        \begin{align}
            c_{j,0}(x, \hat{\alpha}_p) = c_{j}^*, \qquad 
            p_{Y,0}(x, \hat{\alpha}_p) = p_{Y}^*, \qquad 
            p_{k,0}(x, \hat{\alpha}_p, \hat{s}) = p_{k}^*, \qquad k \in \{D,B\},
        \end{align}
        for some constants $c_{j}^*, p_{Y}^*, p_{k}^*$ independent of $(x, \hat{\alpha}_p, \hat{s})$.
        \item The drift and diffusion of the state variables are first-order in $\epsilon$:
        \begin{align}
            \mu_{X} = \mathcal{O}(\epsilon), \qquad 
            \sigma_{X} = \mathcal{O}(\epsilon).
        \end{align}
        \item The first-order and second-order corrections to the consumption-wealth ratio and the price-dividend ratio for the aggregate claim can be written as:\small
        \begin{align}
            c_{j,1}(x, \hat{\alpha}_p) &= \psi_{j,0}(x) + \psi_{j,\hat{\alpha}_p}(x) \hat{\alpha}_p, \qquad 
            c_{j,2}(x, \hat{\alpha}_p) = \Psi_{j,0}(x) + \Psi_{j,\hat{\alpha}_p}(x) \hat{\alpha}_p + \Psi_{j,\hat{\alpha}_p^2}(x) \hat{\alpha}_p^2, \\
            p_{Y,1}(x, \hat{\alpha}_p) &= \psi_{Y,0}(x) + \psi_{Y,\hat{\alpha}_p}(x) \hat{\alpha}_p, \qquad 
            p_{Y,2}(x, \hat{\alpha}_p) = \Psi_{Y,0}(x) + \Psi_{Y,\hat{\alpha}_p}(x) \hat{\alpha}_p + \Psi_{Y,\hat{\alpha}_p^2}(x) \hat{\alpha}_p^2,
        \end{align}\normalsize
        where $\psi_{j,l}(x)$ and $\Psi_{j,l}(x)$, for $j\in\{0,1,2,Y\}$ and $l\in\{0,\hat{\alpha}_p,\hat{\alpha}_p^2\}$, are functions of $x$.


        \item The first-order and second-order corrections to the price-dividend ratio for the individual claim can be written as:\small
        \begin{align}
            p_{k,1}(x, \hat{\alpha}_p, \hat{s}) &= \psi_{k,0}(x) + \psi_{k,\hat{\alpha}_p}(x) \hat{\alpha}_p + \psi_{k,\hat{s}}(x) \hat{s}, \\
            p_{k,2}(x, \hat{\alpha}_p, \hat{s}) &= \Psi_{k,0}(x) + \Psi_{k,\hat{\alpha}_p}(x) \hat{\alpha}_p + \Psi_{k,\hat{s}}(x) \hat{s} + \Psi_{k,\hat{\alpha}_p^2}(x) \hat{\alpha}_p^2 + \Psi_{k,\hat{\alpha}_p \hat{s}}(x) \hat{\alpha}_p \hat{s} + \Psi_{k,\hat{s}^2}(x) \hat{s}^2,
        \end{align}\normalsize
        where $\psi_{k,l}(x)$ and $\Psi_{k,l}(x)$, for $k\in\{D,B\}$ and $l\in\{0,\hat{\alpha}_p,\hat{\alpha}_p^2,\hat{s},\hat{\alpha}_p \hat{s},\hat{s}^2\}$, are functions of $x$.

        % \item The derivatives with respect to the state variables evaluated at $\epsilon = 0$ are given by:
        % \begin{align}
        %     \frac{\partial c_{j}}{\partial x} = 0_{1\times 2}, \qquad 
        %     \frac{\partial c_{j}}{\partial \overline{\alpha}_p} = \psi_{j,\hat{\alpha}_p}(x), \qquad 
        %     \frac{\partial p_{Y}}{\partial x} = 0_{1\times 2}, \qquad 
        %     \frac{\partial p_{Y}}{\partial \overline{\alpha}_p} = \psi_{Y,\hat{\alpha}_p}(x).
        % \end{align}
    \end{enumerate}
\end{lemma}

These properties follow from the definition of the perturbation solution.
When $\epsilon = 0$, the exogenous state variables are constant and investors have identical preferences.
Hence, the zeroth-order terms are independent of $(x,\hat{\alpha}_p,\hat{s})$.
The fact that the drift and diffusion of the state variables are first-order in $\epsilon$ comes from the assumptions on the dynamics of the exogenous state variables, and the fact that the zeroth-order terms are constant.
Properties 3 and 4 describe a solution that is global in the endogenous wealth shares and local in the exogenous state variables.

Our goal is to solve for the coefficients in this perturbation expansion.
We proceed in steps.
We first characterize the benchmark economy.
We then derive the first- and second-order corrections to the consumption-wealth ratios and the price-dividend ratio of the aggregate claim.
Finally, we derive the corresponding corrections to the price-dividend ratios of the individual claims.



\subsection{Proof of Lemma~\ref{lemma:benchmark}}\label{app:lemma1}

\begin{proof}
The assumption $\epsilon = 0$ implies that there is no preference heterogeneity and that passive investors are fully invested in the aggregate risky claim.
We guess and verify that, in this benchmark economy, expected returns are constant and the leverage constraint is slack.
In particular, the wealth distribution plays no role in equilibrium.
This implies that $\mu_{c_j,0}(X) = \mu_{p_Y,0}(X) = 0$ and $\sigma_{c_j,0}(X) = \sigma_{p_Y,0}(X) = 0_{1\times 3}$, where $0_{1\times 3}$ is a $1\times 3$ vector of zeros.

\paragraph{Risk exposure and market price of risk.}
The risk exposure vector for the aggregate claim is given by
\begin{equation}
    \sigma_{R_Y,0}(X) = \sigma.
\end{equation}
 
The market price of risk is therefore
\begin{equation}
    \eta_0(X) = \gamma \sigma,
\end{equation}
using the fact that $h_{a,t} = 0$ and $\chi_{a,t} = 1$.

The risk exposure of investor $j$ is therefore
\begin{equation}
    \sigma_{j,0}(X) = \sigma.
\end{equation}
Because $\overline{\sigma} = \|\sigma\|$, the leverage constraint is satisfied in the benchmark economy.

\paragraph{Risk premium and interest rate.}
The risk premium on the aggregate claim is given by
\begin{equation}
    \pi_{Y,0}(X) = \eta_0(X) \sigma_{R_Y,0}(X)^\top = \gamma \|\sigma\|^2,
\end{equation}
where $\|\sigma\|^2 = \sigma \sigma^\top$.

The interest rate is given by
\begin{align}
    r_0(X) &= \rho 
    + \psi^{-1} \mu 
    -(1+\psi^{-1}) 
        \frac{\gamma }{2} \|\sigma\|^2,
\end{align}
using the fact that $\xi_t = 0$, $\sigma_{j,t} = \sigma$, $\mu_{p_Y,t} = 0$, and $\sigma_{p_Y,t} = 0_{1\times 3}$ in the benchmark economy.

\paragraph{Consumption-wealth ratio and price-dividend ratio.}
The consumption-wealth ratio $c_{j,0}(X)$ is given by
\begin{align}
    c_{j,0}(X) &= \psi \rho 
    + (1-\psi) \left[ r_0(X) + \eta_0(X)\,\sigma^\top - \frac{\gamma}{2} \|\sigma\|^2 \right].
\end{align}

Plugging in the expression for $r_0(X)$ and $\eta_0(X)$, we obtain
\begin{equation}
    c_{j,0}(X) = \rho  - \left(1-\psi^{-1}\right)  \left( \mu - \frac{\gamma\|\sigma\|^2}{2} \right),
\end{equation}
We assume $\rho  > \left(1-\psi^{-1}\right)  \left( \mu - \frac{\gamma\|\sigma\|^2}{2} \right)$, so the benchmark consumption-wealth ratio is positive.

The goods-market-clearing condition then implies
\begin{equation}
    \frac{1}{p_{Y,0}(X)} = \rho  - \left(1-\psi^{-1}\right)  \left( \mu - \frac{\gamma\|\sigma\|^2}{2} \right).
\end{equation}

\paragraph{Drift and diffusion of wealth shares.}
The drift and diffusion of the wealth shares are given by
\begin{align}
    \mu_{x_j,0}(X) &= 0, \qquad
    \sigma_{x_j,0}(X) = 0_{1\times 3}, \qquad
    \forall j \in \{1,2\}.
\end{align}
using the fact that $\kappa = 0$, $\sigma_{j,0}(X) = \sigma_{R_Y,0}(X)$, and $c_{j,0}(X) = \frac{1}{p_{Y,0}(X)}$.

\paragraph{Derivative claim.}
The risk exposure of the derivative claim is given by
\begin{equation}
    \sigma_{R_F,0} = [\sigma_{R_F,0}^1, \sigma_{R_F,0}^2, \sigma_{R_F,0}^3].
\end{equation}
These three quantities are defined by the normalization conditions
\begin{equation}
    \sigma_{R_Y,0} \sigma_{R_F,0}^\top = 0, \qquad 
    \sigma_{R_F,0} \sigma_{R_F,0}^\top = \sigma_{R_Y,0} \sigma_{R_Y,0}^\top,  \qquad 
    \sigma_{R_F,0}^3 = 0.
\end{equation}

Using $\sigma_{R_Y,0} = [\sigma^1,0,0]$, we obtain
\begin{equation}
    \sigma_{R_F,0}^1 = 0, \qquad 
    \sigma_{R_F,0}^2 = \sigma^1, \qquad 
    \sigma_{R_F,0}^3 = 0.
\end{equation}

\paragraph{Portfolio shares.}
The portfolio shares on the aggregate claim and the derivative satisfy
\begin{equation}
    \sigma_{j,0} = \alpha_{j,Y} \sigma_{R_Y,0} + \alpha_{j,F} \sigma_{R_F,0}.
\end{equation}
Using $\sigma_{j,0} = \sigma_{R_Y,0}$, we obtain $\alpha_{j,Y} = 1$ and $\alpha_{j,F} = 0$ for $j = 0,1,2$.

\end{proof}

\subsection{Closed-form solution to individual claims}

I consider next the characterization of the price of the individual claims when the interest rate and the price of risk are constant. 
In contrast to the benchmark economy, we do not assume that $s$ is close to $\overline{s}$, instead we show that one can obtain closed-form solutions, analogous to the ones in \citet{MenzlySantosVeronesi2004}, in the special case $\gamma = 1$. 
The closed-form solutions provide a useful reference to compare with the perturbation solution derived below.

\paragraph{Pricing the individual claims.}
The price of the dividend claim satisfies the pricing condition:
\begin{equation}
    P_{D,t} = \mathbb{E}_t\left[ \int_t^\infty\frac{\Lambda_\tau}{\Lambda_t} s_\tau Y_\tau d\tau \right],
\end{equation}

The SDF and the aggregate endowment follow the processes:
\begin{align}
    \frac{d \Lambda_t}{\Lambda_t} = -r^* dt - \eta^* dZ_t, \qquad \qquad 
    \frac{d Y_t}{Y_t} = \mu dt + \sigma dZ_t.
\end{align}

We can solve for the price of the dividend claim in closed form in the case $\gamma = 1$, so $\eta^ * = \sigma$.
In this case, we can then write the SDF and the aggregate endowment as:
\begin{align}
    \Lambda_t = \exp\left( -\int_0^t \left( r^* + \frac{\sigma \sigma^\top}{2} \right) ds - \sigma Z_t \right) \Lambda_0, \qquad 
    Y_t = \exp\left( \int_0^t \left( \mu - \frac{\sigma \sigma^\top}{2} \right) ds + \sigma Z_t \right) Y_0.
\end{align}

Combining the expressions for $\Lambda_t$ and $Y_t$, we obtain:
\begin{equation}
    P_{D,t} = \int_t^\infty e^{-\delta (\tau-t)} \mathbb{E}_t\left[ s_\tau \right] d\tau Y_t,
\end{equation}
where $\delta = r^* +\sigma \sigma^\top- \mu $, using the fact that $Z_t$ cancels out in the product $\Lambda_t Y_t$.

The conditional expectation of the dividend share is given by
\begin{equation}
    \mathbb{E}_t\left[ s_\tau \right] = \overline{s} + (s_t - \overline{s}) e^{-\kappa_s (\tau-t)}.
\end{equation}

The price of the dividend claim is given by
\begin{equation}
    P_{D,t} = \left[ \frac{\overline{s}}{\delta} + \frac{s_t - \overline{s}}{\delta + \kappa_s} \right] Y_t
\end{equation}
where $Y_t = \exp\left( \int_0^t \left( \mu - \frac{\sigma \sigma^\top}{2} \right) ds + \sigma Z_t \right) Y_0$.

Hence, the normalized price of the dividend and corporate bond claims are given by
\begin{equation}
    p_{D,0}(X) = \frac{1}{\delta}\left[\frac{\delta}{\delta + \kappa_s} s_t + \frac{\kappa_s}{\delta+\kappa_s}  \overline{s}  \right], \qquad \qquad 
    p_{B,0}(X) = \frac{1}{\delta}\left[\frac{\delta}{\delta + \kappa_s} (1-s_t) + \frac{\kappa_s}{\delta+\kappa_s}  (1-\overline{s})  \right].
\end{equation}

\paragraph{Portfolio share of individual claims.}
The portfolio share of investor $j$ on the dividend and bond claims are given by
\begin{equation}
    \alpha_{j,D,t} = \frac{p_{D,t}}{p_{Y,t}}, \qquad \qquad 
    \alpha_{j,B,t} = \frac{p_{B,t}}{p_{Y,t}}.
\end{equation}

In the benchmark economy with $\gamma = 1$, we have that 
\begin{equation}
    \alpha_{j,D,0}(X) = \frac{\delta}{\delta + \kappa_s} s + \frac{\kappa_s}{\delta + \kappa_s} \overline{s}, \qquad \qquad 
    \alpha_{j,B,0}(X) = \frac{\delta}{\delta + \kappa_s} (1-s) + \frac{\kappa_s}{\delta + \kappa_s} (1-\overline{s}),
\end{equation}
using the fact that $\delta = 1/p_{Y,0}(X)$ when $\gamma = 1$.

% \paragraph{Convenient reparametrization.} 
% It is instructive to consider a convenient reparametrization of the price-dividend ratios.
% Consider the change of variables: $\tilde{p}_{D,t} = p_{D,t} s_t$ and $\tilde{p}_{B,t} = p_{B,t} (1-s_t)$, both being  affine functions of the state $s_t$. 
% Notice this implies that $P_{k,t} = \tilde{p}_{k,t} Y_t$ for $k = D,B$.
% In terms of this reparametrization, the pricing condition for the dividend and bond claims is given by
% \begin{equation}
%     r + \sigma_{R_k} \eta^\top = \frac{s_t}{\tilde{p}_{k}} + \mu + \mu_{\tilde{p}_k} + \sigma \sigma_{\tilde{p}_k}^\top,
% \end{equation}
% where $\sigma_{R_k} = \sigma + \sigma_{p_k}$ and $(\mu_{p_k}, \sigma_{p_k})$ are given by Ito's lemma:
% \begin{equation}
%     \mu_{p_k} = \frac{\tilde{p}_{k,s}}{\tilde{p}_{k}} \mu_s + \frac{1}{2}\frac{\tilde{p}_{k,ss}}{\tilde{p}_{k}} \tilde{\sigma}_s \tilde{\sigma}_s^\top, \qquad \qquad 
%     \sigma_{p_k} = \frac{\tilde{p}_{k,s}}{\tilde{p}_{k}} \tilde{\sigma}_{s},
% \end{equation}
% where $\tilde{\sigma}_s = s (1-s)\sigma_s$.

% In the case $\gamma = 1$, we have that
% \begin{equation}
%     \tilde{p}_D = \frac{1}{\delta} \left[ \frac{\delta}{\delta + \kappa_s} s + \frac{\kappa_s}{\delta + \kappa_s} \overline{s} \right],\qquad \qquad 
%     \tilde{p}_B = \frac{1}{\delta} \left[ \frac{\delta}{\delta + \kappa_s} (1-s) + \frac{\kappa_s}{\delta + \kappa_s} (1-\overline{s}) \right],
% \end{equation}

% This implies that
% \begin{equation}
%     \sigma_{\tilde{p}_D} = \frac{1}{\delta + \kappa_s} \frac{\tilde{\sigma}_s}{\tilde{p}_D},\qquad \qquad 
%     \sigma_{\tilde{p}_B} = -\frac{1}{\delta + \kappa_s} \frac{\tilde{\sigma}_s}{\tilde{p}_B}.
% \end{equation}

% The risk premium for the dividend and bond claims are given by
% \begin{equation}
%     \pi_{D,0} = \eta_0 \sigma_{R_D}^\top = \sigma \sigma^\top + \frac{1}{\delta + \kappa_s} \frac{\sigma \tilde{\sigma}_s^T}{\tilde{p}_D}, \qquad \qquad 
%     \pi_{B,0} = \eta_0 \sigma_{R_B}^\top = \sigma \sigma^\top -\frac{1}{\delta + \kappa_s} \frac{\sigma \tilde{\sigma}_s^T}{\tilde{p}_B}.
% \end{equation}
% using the fact that $\eta_0 = \sigma$ when $\gamma = 1$.

% In this case the drift of $\tilde{p}_D$ and $\tilde{p}_B$ are given by
% \begin{equation}
%     \mu_{\tilde{p}_D} = \frac{1}{\delta + \kappa_s} \frac{\mu_s}{\tilde{p}_D},\qquad \qquad 
%     \mu_{\tilde{p}_B} = -\frac{1}{\delta + \kappa_s} \frac{\mu_s}{\tilde{p}_B},
% \end{equation}
% which implies that
% \begin{equation}
%     \frac{s}{\tilde{p}_D} + \mu_{\tilde{p}_D} = \delta \frac{s}{ \frac{\delta}{\delta + \kappa_s} s + \frac{\kappa_s}{\delta + \kappa_s} \overline{s} }  +
%     \delta\frac{\kappa_s}{\delta+\kappa_s} \frac{\overline{s}-s}{ \frac{\delta}{\delta + \kappa_s} s + \frac{\kappa_s}{\delta + \kappa_s} \overline{s} } = \delta,
% \end{equation}
% and, similarly, $\frac{1-s}{\tilde{p}_B} + \mu_{\tilde{p}_B} = \delta$.

% This implies that the pricing condition for the dividend and bond claims reduces to
% \begin{equation}
%     r + \sigma \sigma^\top - \mu = \delta.
% \end{equation}
% % which is consistent with the definition of $\delta$. 



\subsection{Proof of Proposition~\ref{prop:first_order}}\label{app:prop_first_order}

\begin{proof}
Consistent with Proposition~\ref{prop:mft}, the consumption-wealth ratios and the price-dividend ratio of the aggregate claim are independent of the dividend-share shock $s$.
We can then focus on the reduced set of states $\hat{X} = (x, \overline{\alpha}_{p})$, where $x = (x_1, x_2)$ is the vector of wealth shares of the active investors.

We write the portfolio share of passive investors as $\overline{\alpha}_{p} = 1 + \hat{\alpha}_{p} \epsilon$, where $\hat{\alpha}_{p}$ follows a process independent of $\epsilon$:
\begin{equation}
    d \hat{\alpha}_{p} = \kappa_p (\hat{\alpha} - \hat{\alpha}_{p}) dt + \sigma_p dZ_t.
\end{equation}
This implies $\overline{\alpha} = 1+\hat{\alpha} \epsilon$ and that the diffusion of $\overline{\alpha}_{p}$ is of order $\mathcal{O}(\epsilon)$.

We consider the following expansion for the consumption-wealth ratio of investor $j$ and the price-dividend ratio of the aggregate claim:
\begin{align}
    c_{j}(x, \overline{\alpha}_p;\epsilon) &= c_{j,0}(x, \hat{\alpha}_p) + c_{j,1}(x, \hat{\alpha}_p) \epsilon+\mathcal{O}(\epsilon^2),\\
    p_{Y}(x, \overline{\alpha}_p;\epsilon) &= p_{Y,0}(x, \hat{\alpha}_p) + p_{Y,1}(x, \hat{\alpha}_p) \epsilon+\mathcal{O}(\epsilon^2).
\end{align}

The terms $c_{j,0}(x, \hat{\alpha}_p)$ and $p_{Y,0}(x, \hat{\alpha}_p)$ are independent of $(x, \hat{\alpha}_p)$ and are characterized in Lemma~\ref{lemma:benchmark}.
We focus on the first-order correction terms $c_{j,1}(x, \hat{\alpha}_p)$ and $p_{Y,1}(x, \hat{\alpha}_p)$.

\paragraph{Diffusion for consumption-wealth ratios and price-dividend ratio.}
Given the specification of $\overline{\alpha}_{p}$, we have $\mu_{\overline{\alpha}_{p}} = \mathcal{O}(\epsilon)$ and $\sigma_{\overline{\alpha}_{p}} = \mathcal{O}(\epsilon)$.
Similarly, $\mu_{x_{j}} = \mathcal{O}(\epsilon)$ and $\sigma_{x_{j}} = \mathcal{O}(\epsilon)$ for $j = 1,2$.
Hence, the drift and diffusion of the state variables are first-order in $\epsilon$, that is, $\mu_{\hat{X}} = \mathcal{O}(\epsilon)$ and $\sigma_{\hat{X}} = \mathcal{O}(\epsilon)$.

We characterize next the derivatives of the consumption-wealth ratios and the price-dividend ratio with respect to $\hat{X}$.
Given the consumption-wealth ratios and the price-dividend ratio for the aggregate claim are independent of $(x, \hat{\alpha}_{p})$ at $\epsilon = 0$, we have that
$c_{j,x_j}=\mathcal{O}(\epsilon)$ and $p_{Y,x_j}=\mathcal{O}(\epsilon)$, as these derivatives are zero at order zero.
Moreover, we guess-and-verify that the derivatives of the consumption-wealth ratios and the price-dividend ratio with respect to $\overline{\alpha}_{p}$ are first-order in $\epsilon$:
\begin{equation}
    c_{j,\overline{\alpha}_{p}} = \mathcal{O}(\epsilon); \qquad \qquad p_{Y,\overline{\alpha}_{p}} = \mathcal{O}(\epsilon).
\end{equation}

From the expansion for the consumption-wealth ratios, we can express the derivatives with respect to $\overline{\alpha}_p$ in terms of the derivatives of the first-order correction terms with respect to $\hat{\alpha}_p$:
\begin{equation}
    c_{j,\overline{\alpha}_{p}} = \frac{\partial c_{j,1}}{\partial \hat{\alpha}_{p}}  +\mathcal{O}(\epsilon),
\end{equation}
using the fact that $c_{j,0}$ is independent of $\hat{\alpha}_p$.

Hence, $c_{j,\overline{\alpha}_{p}} = \mathcal{O}(\epsilon)$ implies $\frac{\partial c_{j,1}}{\partial \hat{\alpha}_{p}} = 0$, so $c_{j,1}$ is independent of $\hat{\alpha}_{p}$.
A similar argument shows that $p_{Y,\overline{\alpha}_{p}} = \mathcal{O}(\epsilon)$ implies that $p_{Y,1}$ is independent of $\hat{\alpha}_{p}$.

We can then write the diffusion terms for the consumption-wealth ratios and the price-dividend ratio are second-order terms in $\epsilon$:
\begin{align}
    \sigma_{c_j} &= \underbrace{\frac{c_{j,\hat{X}}}{c_j}}_{\mathcal{O}(\epsilon)} \times \underbrace{\sigma_{\hat{X}}}_{\mathcal{O}(\epsilon)} = \mathcal{O}(\epsilon^2), \qquad \qquad 
    \sigma_{p_Y} = \underbrace{\frac{p_{Y,\hat{X}}}{p_Y}}_{\mathcal{O}(\epsilon)} \times \underbrace{\sigma_{\hat{X}}}_{\mathcal{O}(\epsilon)} = \mathcal{O}(\epsilon^2),
\end{align}
This implies that $\sigma_{R_Y,1} = 0_{1\times 3}$.

\paragraph{Risk exposure and market price of risk.} The risk exposure for the passive investor is given by
\begin{equation}
    \sigma_{0,1} = \hat{\alpha}_{p} \sigma. 
\end{equation}

The first-order correction to the risk exposure for an active investor is given by
\begin{equation}
    \sigma_{j,1} = \frac{\eta_1}{\gamma} - \frac{\eta_0}{\gamma} (\hat{\gamma}_j+ \hat{\nu}_{j}),
\end{equation}
where $\hat{\nu}_{j}$ is the Lagrange multiplier on the leverage constraint for investor $j$.

Expanding the market clearing condition for risky assets, we obtain
\begin{equation}
    \sum_{j=0}^2 x_{j} \sigma_{j,1} = 0
\end{equation}
using the fact that $\sigma_{R_Y,1} = 0$.

The first-order correction to the market price of risk is given by
\begin{equation}
   \eta_1  =  \gamma \sigma \left[\sum_{j=1}^2 \frac{x_j}{x_a}(\hat{\gamma}_j+ \hat{\nu}_{j}) -\frac{x_0}{x_a}\hat{\alpha}_{p} \right],
\end{equation}
where $x_a = 1 - x_0$.

The risk exposure for the first active investor is
\begin{equation}
    \sigma_{1,1} = - \left[ \frac{x_2}{x_a}\left(\hat{\gamma}_1 - \hat{\gamma}_2-\hat{\nu}_{2}\right) +\frac{x_0}{x_a}\hat{\alpha}_{p} \right]\sigma.
\end{equation}
and the risk exposure for the second active investor is
\begin{equation}
    \sigma_{2,1} = \left[ \frac{x_1}{x_a}\left(\hat{\gamma}_1 - \hat{\gamma}_2 - \hat{\nu}_{2}\right)-\frac{x_0}{x_a}\hat{\alpha}_{p}  \right]\sigma.
\end{equation}

\paragraph{Lagrange multiplier on the leverage constraint.}
The leverage constraint for investor $j$ is
\begin{equation}
    \|\sigma_j\| \leq \overline{\sigma}.
\end{equation}
Using $\sigma_j=\sigma+\sigma_{j,1}\epsilon+\mathcal{O}(\epsilon^2)$ and $\overline{\sigma}=\|\sigma\|(1+\hat{\sigma}\epsilon)$, the first-order expansion of the constraint is
\begin{equation}
    \sqrt{(\sigma+\sigma_{j,1}\epsilon)(\sigma+\sigma_{j,1}\epsilon)^\top}
    = \|\sigma\|\left(1+\frac{\sigma_{j,1}\sigma^\top}{\sigma\sigma^\top}\epsilon\right)
    +\mathcal{O}(\epsilon^2).
\end{equation}
Therefore, the first-order constraint is
\begin{equation}
    \frac{\sigma_{j,1}\sigma^\top}{\sigma\sigma^\top} \leq \hat{\sigma}.
\end{equation}

We assume that investor $1$ is unconstrained and investor $2$ may be constrained, so $\hat{\nu}_1=0$ and $\hat{\nu}_2\geq 0$.
For investor $2$, using the expression for $\sigma_{2,1}$ above, the first-order leverage constraint is
\begin{equation}
    \frac{x_1}{x_a}\left(\hat{\gamma}_1 - \hat{\gamma}_2 - \hat{\nu}_{2}\right) -\frac{x_0}{x_a}\hat{\alpha}_{p,1} \leq \hat{\sigma}.
\end{equation}
If the constraint binds, this condition holds with equality, which implies
\begin{equation}
    \hat{\nu}_{2} = \hat{\gamma}_1-
    \hat{\gamma}_2 -\frac{x_0}{x_1}\hat{\alpha}_{p,1} -\frac{x_a}{x_1} \hat{\sigma}.
\end{equation}
If the constraint is slack, then $\hat{\nu}_{2}=0$ and
\begin{equation}
    \frac{x_1}{x_a}\left(\hat{\gamma}_1 - \hat{\gamma}_2 \right) -\frac{x_0}{x_a}\hat{\alpha}_{p,1} < \hat{\sigma},
\end{equation}
or, equivalently,
\begin{equation}
    \hat{\gamma}_1-\hat{\gamma}_2 -\frac{x_0}{x_1}\hat{\alpha}_{p,1} -\frac{x_a}{x_1} \hat{\sigma} <0.
\end{equation}
Hence, the first-order Lagrange multiplier is
\begin{equation}
    \hat{\nu}_{2} = \max\left\{ \hat{\gamma}_1-\hat{\gamma}_2 -\frac{x_0}{x_1}\hat{\alpha}_{p,1} -\frac{x_a}{x_1} \hat{\sigma}, 0 \right\}.
\end{equation}

When the constraint binds, substituting the expression for $\hat{\nu}_2$ into the market price of risk yields
\begin{equation}
    \eta_1 = \gamma \sigma \left[\hat{\gamma}_1 - \frac{x_2}{x_1}\hat{\sigma}  -\frac{x_0}{x_1}\hat{\alpha}_{p,1}  \right].
\end{equation}
When the constraint is slack, $\hat{\nu}_2=0$, and the market price of risk is
\begin{equation}
    \eta_1  =  \gamma \sigma \left[\sum_{j=1}^2 \frac{x_j}{x_a} \hat{\gamma}_j -\frac{x_0}{x_a}\hat{\alpha}_{p,1} \right].
\end{equation}
Equivalently, both cases can be summarized as
\begin{equation}
    \eta_1 = \gamma \sigma \left[\sum_{j=1}^2 \frac{x_j}{x_a}(\hat{\gamma}_j+ \hat{\nu}_{j}) -\frac{x_0}{x_a}\hat{\alpha}_{p,1} \right],
\end{equation}
with $\hat{\nu}_1=0$ and $\hat{\nu}_2$ given above.

\paragraph{Risk premium.}
The first-order risk premium is
\begin{equation}
    \pi_{Y,1} = \eta_1 \sigma^\top = \gamma \|\sigma\|^2 \left[\sum_{j=1}^2 \frac{x_j}{x_a}(\hat{\gamma}_j+ \hat{\nu}_{j}) -\frac{x_0}{x_a}\hat{\alpha}_{p,1} \right].
\end{equation}
It increases with the average risk aversion of active investors, with tighter (binding) leverage constraints, and when active investors bear a larger share of aggregate risk.

\paragraph{Drift for consumption-wealth ratios and price-dividend ratio.}
It\^o's lemma implies that the drifts are
\begin{align}
    \mu_{c_j} &= \frac{c_{j,\hat{X}}}{c_j} \mu_{\hat{X}} + \frac{1}{2}\sum_{k=1}^3 \sigma_{\hat{X},k}^\top \frac{c_{j,\hat{X} \hat{X}}}{c_j} \sigma_{\hat{X},k}, \\
    \mu_{p_Y} &= \frac{p_{Y,\hat{X}}}{p_Y} \mu_{\hat{X}} + \frac{1}{2}\sum_{k=1}^3 \sigma_{\hat{X},k}^\top \frac{p_{Y,\hat{X} \hat{X}}}{p_Y} \sigma_{\hat{X},k},
\end{align}
where $\sigma_{\hat{X},k}$ is the $k$-th column of the matrix $\sigma_{\hat{X}}$.

Hence, $\mu_{c_j} = \mathcal{O}(\epsilon^2)$ and $\mu_{p_Y} = \mathcal{O}(\epsilon^2)$, since the gradients are $\mathcal{O}(\epsilon)$ and both $\mu_{\hat{X}}$ and $\sigma_{\hat{X}}$ are $\mathcal{O}(\epsilon)$.



\paragraph{Interest rate.}
The first-order correction to the interest rate is
\begin{align}
    r_1 &= (1-\psi^{-1}) \frac{\gamma  \|\sigma\|^2}{2}\sum_{j=0}^2 x_{j} \hat{\gamma}_j 
    + (1-\psi^{-1}) \gamma \|\sigma\|^2 \sum_{j=0}^2 x_j \frac{\sigma_{j,1}\sigma^\top}{\|\sigma\|^2}
    - \pi_{Y,1},
\end{align}
where we used the fact that $\xi_j = \mathcal{O}(\epsilon^2)$, as $\mu_{c_j} = \mathcal{O}(\epsilon^2)$ and $\sigma_{c_j} = \mathcal{O}(\epsilon^2)$.

Using $\sum_{j=0}^2 x_{j} \sigma_{j,1} = 0_{1\times 3}$, we obtain
\begin{equation}
    r_1 = 
    - \gamma \|\sigma\|^2 \left[\sum_{j=1}^2 \frac{x_j}{x_a}(\hat{\gamma}_j+ \hat{\nu}_{j})- \sum_{j=0}^2 x_{j} \hat{\gamma}_j  -\frac{x_0}{x_a}\hat{\alpha}_{p,1} +\frac{1+\psi^{-1}}{2} \sum_{j=0}^2 x_{j} \hat{\gamma}_j \right].
\end{equation}

Notice that we can write the expression above as
\begin{equation}
    r_1 = -\gamma \|\sigma\|^2 \frac{1+\psi^{-1}}{2} \sum_{j=0}^2 x_{j} \hat{\gamma}_j - \left[ \pi_{Y,1} - \pi_{Y,1}^{fb}\right],    
\end{equation}
where
\begin{equation}
    \pi_{Y,1}^{fb} = \gamma \|\sigma\|^2 \sum_{j=0}^2 x_{j} \hat{\gamma}_j.
\end{equation}

The first term captures the effect of heterogeneity in risk aversion in the frictionless benchmark (no passive investors or leverage constraints).
The second term captures the impact of frictions through the wedge between the equilibrium risk premium and its frictionless counterpart.

\paragraph{Consumption-wealth ratio.} The consumption-wealth ratio for investor $j$ is given by
\begin{align}
    c_{j,1} &= (1-\psi) \left[ r_1 + \eta_1\,\sigma^\top + \eta_0 \sigma_{j,1}^\top -\gamma \sigma \sigma_{j,1}^\top - \frac{\gamma}{2} \|\sigma\|^2 \hat{\gamma}_j \right].
\end{align}

% Using $\eta_0 = \gamma \sigma$ and $\eta_1 \sigma^\top = \pi_{Y,1}$, we obtain
% \begin{equation}
%     c_{j}(X,\epsilon) = c_{j,0} + c_{j,1}\epsilon + \mathcal{O}(\epsilon^2),
% \end{equation}

Substituting for $r_1$ and $\pi_{Y,1}$, the first-order correction to the consumption-wealth ratio is
\begin{equation}
    c_{j,1} = \frac{\gamma \|\sigma\|^2}{2} (1-\psi)\left[ \left(1-\psi^{-1}\right) \sum_{k=0}^2 x_k\hat{\gamma}_k-\hat{\gamma}_j\right],
\end{equation}
which is independent of $\hat{\alpha}_p$ as conjectured.

\paragraph{Price-dividend ratio of aggregate claim.}
The pricing condition implies
\begin{equation}
    -\frac{1}{(p_{Y}^*)^2} p_{Y,1} = r_1 + \pi_{Y,1}.
\end{equation}

Hence, the first-order correction to the price-dividend ratio is
\begin{equation}
    \frac{p_{Y,1}}{p_{Y}^*} = - p_{Y}^* (1-\psi^{-1}) \frac{\gamma  \|\sigma\|^2}{2}\sum_{j=0}^2 x_{j} \hat{\gamma}_j,
\end{equation}
which is independent of $\hat{\alpha}_p$ as conjectured.
If $\psi>1$, an increase in average risk aversion reduces the price-dividend ratio of the aggregate claim.

\paragraph{State dynamics.} The diffusion of the wealth share of investor $j$, $j=1,2$, is given by
\begin{align}
    \sigma_{x_j,1} &= x_j (\sigma_{j,1} - \sigma_{R,1})\nonumber\\
    &= x_j \left[ \sum_{k=1}^2 \frac{x_k}{x_a}(\hat{\gamma}_k+ \hat{\nu}_k) -\frac{x_0}{x_a}\hat{\alpha}_{p,1} - (\hat{\gamma}_j+ \hat{\nu}_j)\right] \sigma.
    % &= x_j \min\left\{ \sum_{k\in\mathcal{J}^u} \frac{x_k}{x_u} \hat{\gamma}_j -\left(\frac{\hat{\alpha}_p  x_0+\frac{\hat{\sigma}}{\|\sigma\|} x_c }{1-x_0-x_c}\right)- \hat{\gamma}_j  ,\frac{\hat{\sigma}}{\|\sigma\|}  \right\} \sigma \epsilon + \mathcal{O}(\epsilon^2).
\end{align}

The drift of the wealth share of investor $j$, $j=1,2$, is given by
\begin{align}
    \mu_{x_j,1} &= x_{j} \left[ 
        (\gamma-1)\sigma  \sigma_{j,1}^\top +\frac{\gamma \|\sigma\|^2}{2} (\psi-1) \left( \sum_{k=0}^2 x_k\hat{\gamma}_k  - \hat{\gamma}_j\right)
     \right] + \hat{\kappa} (\omega_j - x_{j}),
\end{align}

Finally, the law of motion of $\hat{\alpha}_p$ is
\begin{equation}
    d \hat{\alpha}_p = \kappa_p ( \hat{\alpha} - \hat{\alpha}_p )dt + \sigma_p d Z_t,
\end{equation}
so the drift and diffusion of $\overline{\alpha}_p$ are $\mu_{\overline{\alpha}_p} = \kappa_p ( \hat{\alpha} - \hat{\alpha}_p )\epsilon$ and $\sigma_{\overline{\alpha}_p} = \sigma_p \epsilon$.

\end{proof}

% REVISED SUBSECTION: First-order correction of individual claims
\subsection{First-order correction of individual claims}
We next price the dividend and bond claims, whose payoffs are $D=sY$ and $B=(1-s)Y$.
It is convenient to work with prices normalized by aggregate output,
\begin{equation}
    p_{k} \equiv \frac{P_k}{Y}, \qquad k\in\{D,B\}.
\end{equation}
We consider the expansion
\begin{equation}
    p_k(x, \overline{\alpha}_p, s;\epsilon) = p_{k,0} + p_{k,1}(x, \hat{\alpha}_p, \hat{s}) \epsilon + \mathcal{O}(\epsilon^2),
\end{equation}
where $p_{k,0}$ is constant.
Let $s_D=s$ and $s_B=1-s$, so that
\begin{equation}
    s_k = \overline{s}_k + \hat{s}_k \epsilon,
\end{equation}
with
\begin{equation}
    \overline{s}_D=\overline{s}, \quad \overline{s}_B=1-\overline{s}, \qquad
    \hat{s}_D=\hat{s}, \quad \hat{s}_B=-\hat{s}.
\end{equation}

The pricing condition for claim $k$ is
\begin{equation}
    r + \sigma_{R_k} \eta^\top = \frac{s_k}{p_{k}} + \mu + \mu_{p_k} + \sigma \sigma_{p_k}^\top.
\end{equation}

\paragraph{Zeroth-order term.}
At $\epsilon=0$, we have
\begin{equation}
    p_{k,0} = \frac{\overline{s}_k}{\delta}, \qquad
    \delta \equiv r^* + \gamma\|\sigma\|^2 - \mu.
\end{equation}

\paragraph{First-order correction.}
The first-order pricing equation implies
\begin{equation}
    \delta p_{k,1} + (r_1 + \pi_{Y,1}) p_{k,0} = \hat{s}_k - \kappa_s \hat{s} \frac{\partial p_{k,1}}{\partial \hat{s}} + (1-\gamma) \sigma \frac{\partial p_{k,1}}{\partial \hat{s}} \overline{s}(1-\overline{s})\sigma_s^\top.
\end{equation}

\paragraph{Dividend claim.}
We conjecture
\begin{equation}
    p_{D,1} = a_D(x) + b_D(x) \hat{s}.
\end{equation}
Matching coefficients gives
\begin{equation}
    b_D = \frac{1}{\delta + \kappa_s},
\end{equation}
and
\begin{equation}
    a_D = - (r_1 + \pi_{Y,1}) \frac{\overline{s}}{\delta^2} - \frac{\gamma-1}{\delta(\delta+\kappa_s)} \overline{s}(1-\overline{s}) \sigma \sigma_s^\top.
\end{equation}
Hence
\begin{equation}
    p_{D,1} = \frac{\hat{s}}{\delta + \kappa_s} - (r_1 + \pi_{Y,1}) \frac{\overline{s}}{\delta^2} - \frac{\gamma-1}{\delta(\delta+\kappa_s)} \overline{s}(1-\overline{s}) \sigma \sigma_s^\top.
\end{equation}

\paragraph{Bond claim.}
Similarly, the normalized price of the bond claim is given by
\begin{equation}
    p_{B,1} = -\frac{\hat{s}}{\delta + \kappa_s} - (r_1 + \pi_{Y,1}) \frac{1-\overline{s}}{\delta^2} + \frac{\gamma-1}{\delta(\delta+\kappa_s)} \overline{s}(1-\overline{s}) \sigma \sigma_s^\top.
\end{equation}

\paragraph{Consistency.}
Adding both claims yields
\begin{equation}
    p_{D,1} + p_{B,1} = -\frac{r_1 + \pi_{Y,1}}{\delta^2},
\end{equation}
which is consistent with the first-order correction of the aggregate claim.


% \subsection{Determining the set of constrained agents}

% In this section, we focus on the case where investors are in the inner region, that is, the following condition holds:
% \begin{align}\label{eq:inner-region}
%     \left|\sum_{k\in\mathcal{J}^u} \frac{x_k}{x_u} \hat{\gamma}_k -\left(\frac{\hat{\alpha}_p  x_0+\frac{\hat{\sigma}}{\|\sigma\|} x_c }{1-x_0-x_c}\right)- \hat{\gamma}_j -  \frac{\hat{\sigma}}{\|\sigma\|}\right| \gg 0,
% \end{align}
% for $j = 1,2,\ldots,J$.

% \paragraph{The case with one active investor} Suppose first that $J = 1$. In this case, from the market clearing condition for the risky asset, the risk exposure of investor 1 must satisfy the condition
% \begin{equation}
%     \hat{\alpha}_p \sigma x_0 + \sigma_{j,1} (1-x_0) = 0 \Rightarrow \sigma_{j,1} = - \frac{\hat{\alpha}_p \sigma x_0}{1-x_0}.
% \end{equation}

% From condition~\eqref{eq:inner-region}, this is feasible if the following condition is satisfied
% \begin{equation}
%     - \frac{\hat{\alpha}_p  x_0}{1-x_0} \ll \frac{\hat{\sigma}}{\sigma}.
% \end{equation}

% The condition above is always satisfied if $\hat{\alpha}_p \geq 0$. If $\hat{\alpha}_p < 0$, then the condition above imposes an upper bound on the wealth share of passive investors $x_0$:
% \begin{equation}
%      x_0 \ll \frac{\hat{\sigma}}{\hat{\sigma} - \hat{\alpha}_p \sigma}. 
% \end{equation}

% \paragraph{The case with two active investors} Suppose now that $J = 2$. Let's assume that $\hat{\gamma}_1\leq \hat{\gamma}_2$. Similarly, this case is feasible if 
% \begin{equation}
%     \frac{x_1}{x_1+x_2} \hat{\gamma}_1 + \frac{x_2}{x_1+x_2} \hat{\gamma}_2  - \frac{\hat{\alpha}_p x_0}{1-x_0} - \hat{\gamma}_1 - \frac{\hat{\sigma}}{\sigma} \ll 0.
% \end{equation}

% Rearranging the expression above, we obtain
% \begin{equation}\label{eq:two-agents}
%     \frac{x_2}{1-x_0} (\hat{\gamma}_2-\hat{\gamma}_1) - \frac{\hat{\alpha}_p x_0}{1-x_0}  \ll \frac{\hat{\sigma}}{\sigma}.
% \end{equation}

% In the case where $\hat{\alpha}_p < 0$, condition~\eqref{eq:two-agents} is more likely to be satisfied when both $x_0$ and $x_2 $ are small. This implies that the leverage constraint is more likely not be binding when the low risk averse investor ($ j=1 $ in this case) is better capitalized.

% Investor 1 will be constrained if the following condition is satisfied
% \begin{equation}
%     \hat{\gamma}_2 - \frac{\hat{\alpha}_p x_0}{x_2} - \frac{\hat{\sigma}}{\sigma} \frac{x_1}{x_2} - \hat{\gamma}_1  \gg \frac{\hat{\sigma}}{\sigma},
% \end{equation}
% which can be written as
% \begin{equation}
%     \frac{x_2}{1-x_0} (\hat{\gamma}_2 - \hat{\gamma}_1)- \frac{\hat{\alpha}_p x_0}{1-x_0} \gg \frac{\hat{\sigma}}{\sigma}.
% \end{equation}

% Whether investor 1 is constrained depends on second-order terms if the following condition holds:
% \begin{equation}
%     \left|\frac{x_2}{1-x_0} (\hat{\gamma}_2-\hat{\gamma}_1) - \frac{\hat{\alpha}_p x_0}{1-x_0}\right| < \frac{\hat{\sigma}}{\sigma} + b \epsilon,
% \end{equation}
% for some constant $b$. 

% \paragraph{The case with three active investors} Suppose now that $J = 3$. Let's assume that $\hat{\gamma}_1\leq \hat{\gamma}_2\leq \hat{\gamma}_3$. Investors will be unconstrained if 
% \begin{equation}
%     \frac{x_2}{1-x_0} (\hat{\gamma}_2-\hat{\gamma}_1) + \frac{x_3}{1-x_0} (\hat{\gamma}_3-\hat{\gamma}_1)  - \frac{\hat{\alpha}_p x_0}{1-x_0} \ll \frac{\hat{\sigma}}{\sigma},
% \end{equation}
% which is more likely to hold as $x_1$ is relatively high, that is, type-1 investors are better capitalized.

% Investor 1 will be constrained if the following condition holds
% \begin{equation}
%     \frac{x_2}{1-x_0-x_1} \hat{\gamma}_2 + \frac{x_3}{1-x_0-x_1} \hat{\gamma}_3 - \frac{\hat{\alpha}_p x_0}{1-x_0 - x_1} - \frac{x_1}{1-x_0-x_1} \frac{\hat{\sigma}}{\sigma} - \hat{\gamma}_1 \gg \frac{\hat{\sigma}}{\sigma},
% \end{equation}
% the condition above can be written as
% \begin{equation}
%     \frac{x_2}{1-x_0} (\hat{\gamma}_2-\hat{\gamma}_1) + \frac{x_3}{1-x_0} (\hat{\gamma}_3-\hat{\gamma}_2) - \frac{\hat{\alpha}_p x_0}{1-x_0 } \gg \frac{\hat{\sigma}}{\sigma},
% \end{equation}

% Investors 1 and 2 will be constrained if the following condition holds
% \begin{equation}
%      \hat{\gamma}_3 - \frac{\hat{\alpha}_p x_0}{x_3} - \frac{x_1+x_2}{x_3} \frac{\hat{\sigma}}{\sigma} - \hat{\gamma}_2 \gg \frac{\hat{\sigma}}{\sigma},
% \end{equation}
% which can be written as
% \begin{equation}
%     \frac{x_3}{1-x_0} (\hat{\gamma}_3 -\hat{\gamma}_2) - \frac{\hat{\alpha}_p x_0}{1-x_0} \gg \frac{\hat{\sigma}}{\sigma}.
% \end{equation}


\subsection{Second-order correction: no dynamics in passive portfolio share}\label{app:prop_second_order}

We next derive the second-order correction.
To keep the algebra transparent, we proceed in two steps.
First, we consider the case in which the portfolio share of passive investors is fixed, so
\begin{equation}
    \kappa_p = 0, \qquad \sigma_p = 0_{1\times 3}.
\end{equation}
Equivalently, $\hat{\alpha}_p$ is a fixed parameter rather than a stochastic state variable.
This removes all drift and diffusion terms associated with $\overline{\alpha}_p$ and leaves the endogenous wealth shares $x=(x_1,x_2)$ as the only state variables that contribute to the second-order drift and diffusion of equilibrium objects.
After deriving this benchmark second-order system, we then allow for non-zero drift and diffusion in the passive investors' portfolio share.

% In Proposition~\ref{prop:second_order}, we compute the second-order correction for our economy.

% \begin{prop}[Second-order correction]\label{prop:second_order}
% 	Suppose $\rho>\left(1-\psi^{-1}\right)\left( \mu - \frac{\gamma
% 		\sigma^2}{2}  \right)$.  Then,
% 	\begin{enumerate}[label=(\roman*)]
% 		\item The second-order correction for the consumption-wealth ratio and risk exposure are given by:
% 		\begin{align}
% 			c_{j,2}(X) &= (1-\psi) \frac{\gamma \sigma^2}{2} \left[ \left(1-\psi^{-1}\right) \sum_{k=0}^J x_{k,t} \hat{\gamma}_k -\hat{\gamma}_j \right]\\
% 			%		\sigma_{j,1}(X) &= \sigma\min\left\{ \sum_{k\in\mathcal{J}^u} \frac{x_k}{x_u} \hat{\gamma}_k- \hat{\gamma}_j  -\frac{\hat{\alpha}_p  x_0+\frac{\hat{\sigma}}{\sigma} x_c }{1-x_0-x_c},\frac{\hat{\sigma}}{\sigma}  \right\}.\nonumber
% 			\sigma_{j,2}(X) &= \frac{\theta_2(X)}{\gamma} - \frac{\theta_1(X)}{\gamma}\hat{\gamma}_j + \frac{\theta_0(X)}{\gamma}\hat{\gamma}_j^2 - (1-\gamma^{-1})   \sigma_{y,2}(X),
% 		\end{align}
% 		where
% 		\begin{equation}
% 			\sigma_{y,2}(X)  = \left(1-\psi^{-1}\right)\frac{\gamma\sigma^2}{2 y_{0}(X)} \sum_{k=1}^J \left( \hat{\gamma}_k - \hat{\gamma}_0  \right) x_k \sigma_{k,1}(X) \frac{\sigma}{\|\sigma\|}.
% 		\end{equation}
% 		\item The second-order correction for the Sharpe ratio, interest rate, and dividend yield are given by:
% 		\begin{align}		
% 			\theta_2(X) &=  - \frac{(\gamma-1) x_c + 1-x_0}{1-x_0-x_c} \sigma_{y,2}(X) - \gamma \sigma \mathbb{E}^u[\hat{\gamma}_j] \frac{ \hat{\alpha}_p x_0+ \frac{\hat{\sigma}}{\sigma}x_c }{1-x_0-x_c} - \gamma \sigma \var^u[\hat{\gamma}_j]\\	
% 			%		r_1(X) &=\gamma\sigma^2 \left[ -\sum_{j\in\mathcal{J}^u} \frac{x_j}{x_u} \hat{\gamma}_j +\left(1-\psi^{-1}\right) \sum_{j=0}^J x_{j,t}\frac{\hat{\gamma}_j}{2}+\frac{\hat{\alpha}_p  x_0+\frac{\hat{\sigma}}{\sigma} x_c }{1-x_0-x_c} \right]\nonumber\\
% 			r_2(X) &= - \theta_2(X) \sigma + \theta_0(X) \sigma_{y,2}(X)-\psi^{-1} (\mu_{y,2}(X) + \sigma \sigma_{y,2}(X))\nonumber\\
% 			&\quad+ \left(1-\psi^{-1}\right) \sum_{j=0}^J x_{j} \gamma\sigma \left[ \sigma_{j,2}(X) + \frac{\sigma_{j,1}^2(X)}{2 \sigma} + \hat{\gamma}_j \sigma_{j,1} \right]\nonumber\\
% 			&\quad+\psi^{-1} \sum_{j=0}^J x_{j}\left[ \mu_{\xi_j,2}(X) + (1-\gamma) \sigma_{\xi_j,2}(X) \sigma \right]\\
%             y_2(X) &= \left(1-\psi^{-1}\right)\sum_{j=0}^{J} x_{j} \gamma \sigma\left(\frac{\sigma_{j, 1}^{2}(X)}{2 \sigma}+\hat{\gamma}_{j} \sigma_{j, 1}(X)\right),\label{eq:y-second-order}
% 			% y_2(X) &= \left(1-\psi^{-1}\right)\left[\gamma \sigma_{y, 2}(X) \sigma+\sum_{j=0}^{J} x_{j} \gamma \sigma\left(\frac{\sigma_{j, 1}^{2}(X)}{2 \sigma}+\sigma_{j, 2}(X)+\hat{\gamma}_{j} \sigma_{j, 1}(X)\right)\right],\nonumber
% 		\end{align}
% 		where
% 		\[\mathbb{E}^{u}[\hat{\gamma}_{j}] \equiv \sum_{j \in \mathcal{J}^{u}} \frac{x_{j}}{x_{u}} \hat{\gamma}_{j}, \quad \text { and } \quad \var^{u}[\hat{\gamma}_{j}] \equiv \sum_{j \in \mathcal{J}^{u}} \frac{x_{j}}{x_{u}} \hat{\gamma}_{j}^{2}-\left(\sum_{j \in \mathcal{J}^{u}} \frac{x_{j}}{x_{u}} \hat{\gamma}_{j}\right)^{2}.\]
% 	\end{enumerate}
% \end{prop}

% \begin{proof}


\noindent\textbf{Step 1: Diffusion for $c_j$ and $y_Y$.}

The expansion of $\sigma_{c_j,t}$ in $\epsilon$ is given by\small
\begin{align}
    \sigma_{c_j}(\hat{X},\epsilon) &= \frac{c_{j,\hat{X}}(\hat{X},\epsilon)}{c_j(\hat{X},\epsilon)} \, \sigma_{\hat{X}}(\hat{X},\epsilon) \nonumber\\
    &= \frac{c_{j,1,x}(x, \hat{\alpha}_p)}{c_{j,0}(x, \hat{\alpha}_p)} \sigma_{x,1}(x, \hat{\alpha}_p) \epsilon^2 +\mathcal{O}(\epsilon^3)\nonumber\\
    &= -(\psi-1)\frac{\left(1-\psi^{-1}\right)\gamma \|\sigma\|^2}{2 c_{j,0}(x, \hat{\alpha}_p)} \sum_{k=1}^2 \left( \hat{\gamma}_k - \hat{\gamma}_0  \right) x_k\sigma_{k,1}(x, \hat{\alpha}_p) \epsilon^2  + \mathcal{O}(\epsilon^3),
\end{align}\normalsize
This follows from $c_{j,x}=\mathcal{O}(\epsilon)$ and $\sigma_x=\mathcal{O}(\epsilon)$, so their product is second order.

At second order, it is convenient to work with the dividend yield of the aggregate claim,
\begin{equation}
    y_Y \equiv \frac{1}{p_Y}.
\end{equation}
This is only a change of variables.
After solving for the second-order correction to $y_Y$, we recover the correction to $p_Y$ from
\begin{equation}
    p_Y = \frac{1}{y_Y}.
\end{equation}
If
\begin{equation}
    y_Y = y_{Y,0} + y_{Y,1}\epsilon + y_{Y,2}\epsilon^2 + \mathcal{O}(\epsilon^3),
\end{equation}
then
\begin{equation}
    p_Y = p_{Y,0} - p_{Y,0}^2 y_{Y,1}\epsilon + \left(p_{Y,0}^3 y_{Y,1}^2 - p_{Y,0}^2 y_{Y,2}\right)\epsilon^2 + \mathcal{O}(\epsilon^3).
\end{equation}

\noindent
We now derive the dynamics of $y_Y$ directly:
% It will be convenient to work with the dividend yield instead of the price-dividend ratio.
% Define the dividend yield as $y_Y = 1/p_{Y}$, and following a similar reasoning, we obtain
\begin{align}\label{eq:sig-y2}
    \sigma_{y_Y}(\hat{X},\epsilon) &= \frac{y_{Y,\hat{X}}(\hat{X},\epsilon)}{y_{Y}(\hat{X},\epsilon)} \sigma_{\hat{X}}(\hat{X},\epsilon) \nonumber\\
    &= \left(1-\psi^{-1}\right)\frac{\gamma\|\sigma\|^2}{2 y_{Y,0}(x, \hat{\alpha}_p)} \sum_{k=1}^2 \left( \hat{\gamma}_k - \hat{\gamma}_0  \right) x_k \sigma_{k,1}(x, \hat{\alpha}_p)  \epsilon^2 + \mathcal{O}(\epsilon^3).
\end{align}
The same scaling applies to $y_Y$, implying that its diffusion starts at order $\epsilon^2$.

The drift of $c_{j,t}$ is given by\small
\begin{align}
    \mu_{c_j}(\hat{X},\epsilon) &= \frac{c_{j,\hat{X}}(\hat{X},\epsilon)}{c_j(\hat{X},\epsilon)} \mu_{\hat{X}}(\hat{X},\epsilon)+ \sum_{k=1}^3\frac{1}{2}  \sigma_{\hat{X},k}^\top(\hat{X},\epsilon) \frac{c_{j,\hat{X}\hat{X}}(\hat{X},\epsilon)}{c_{j}(\hat{X},\epsilon)} \sigma_{\hat{X},k}(\hat{X},\epsilon)\nonumber\\
    &=  \frac{c_{j,1,x}(x, \hat{\alpha}_p)}{c_{j,0}(x, \hat{\alpha}_p)} \mu_{x,1}(x, \hat{\alpha}_p) \epsilon^2+\mathcal{O}(\epsilon^3)\nonumber\\
    &= -(\psi-1)\frac{\left(1-\psi^{-1}\right)\gamma \|\sigma\|^2}{2 c_{j,0}(x, \hat{\alpha}_p)} \sum_{k=1}^2 \left( \hat{\gamma}_k - \hat{\gamma}_0  \right) x_k\mu_{x_k,1}(x, \hat{\alpha}_p) \epsilon^2+ \mathcal{O}(\epsilon^3).
\end{align}\normalsize

The drift of $y_Y,t$ is given by\small
\begin{align}
    \mu_{y_Y}(\hat{X},\epsilon) &= \frac{y_{Y,\hat{X}}(\hat{X},\epsilon)}{y_{Y}(\hat{X},\epsilon)} \mu_{\hat{X}}(\hat{X},\epsilon)+ \sum_{k=1}^3\frac{1}{2}  \sigma_{\hat{X},k}^\top(\hat{X},\epsilon) \frac{y_{Y,\hat{X}\hat{X}}(\hat{X},\epsilon)}{y_{Y}(\hat{X},\epsilon)} \sigma_{\hat{X},k}(\hat{X},\epsilon)\nonumber\\
    &=  \frac{y_{Y,1,x}(x, \hat{\alpha}_p)}{y_{Y,0}(x, \hat{\alpha}_p)} \mu_{x,1}(x, \hat{\alpha}_p) \epsilon^2+\mathcal{O}(\epsilon^3)\nonumber\\
    &= \left(1-\psi^{-1}\right)\frac{\gamma \|\sigma\|^2}{2 y_{Y,0}(x, \hat{\alpha}_p)} \sum_{k=1}^2 \left( \hat{\gamma}_k - \hat{\gamma}_0  \right) x_k\mu_{x_k,1}(x, \hat{\alpha}_p) \epsilon^2+ \mathcal{O}(\epsilon^3).
\end{align}\normalsize


\noindent\textbf{Step 2: Risk exposures of investors.}\\
The risk exposure of an active investor is
\begin{equation}
    \sigma_{j,t}
    = \frac{1}{1+\nu_{j,t}} \left[ \frac{\eta_t}{\gamma_j} + \frac{1-\gamma_j^{-1}}{\psi-1} \, \sigma_{c_j,t} \right].
\end{equation}
It is useful to first define the desired exposure before imposing the leverage multiplier:
\begin{equation}
    q_{j,t}
    \equiv \frac{\eta_t}{\gamma_j} + \frac{1-\gamma_j^{-1}}{\psi-1} \, \sigma_{c_j,t}.
\end{equation}
We write
\begin{align}
    \eta_t &= \eta_0 + \eta_1 \epsilon + \eta_2 \epsilon^2 + \mathcal{O}(\epsilon^3),\\
    \nu_{j,t} &= \hat{\nu}_{j,1}\epsilon + \hat{\nu}_{j,2}\epsilon^2 + \mathcal{O}(\epsilon^3),\\
    \sigma_{c_j,t} &= \sigma_{c_j,2}\epsilon^2 + \mathcal{O}(\epsilon^3),
\end{align}
where $\hat{\nu}_{j,1}=\hat{\nu}_j$ is the first-order multiplier characterized above.
Since $\gamma_j=\gamma(1+\hat{\gamma}_j\epsilon)$, we have
\begin{equation}
    \gamma_j^{-1}=\gamma^{-1}\left(1-\hat{\gamma}_j\epsilon+\hat{\gamma}_j^2\epsilon^2\right)+\mathcal{O}(\epsilon^3).
\end{equation}
Therefore,
\begin{equation}
    q_{j,t}=q_{j,0}+q_{j,1}\epsilon+q_{j,2}\epsilon^2+\mathcal{O}(\epsilon^3),
\end{equation}
with
\begin{align}
    q_{j,0} &= \frac{\eta_0}{\gamma}=\sigma,\\
    q_{j,1} &= \frac{\eta_1}{\gamma}-\sigma\hat{\gamma}_j,\\
    q_{j,2} &= \frac{\eta_2}{\gamma}-\frac{\eta_1}{\gamma}\hat{\gamma}_j+
    \sigma\hat{\gamma}_j^2+\frac{1-\gamma^{-1}}{\psi-1}\sigma_{c_j,2}.
\end{align}
Expanding $\sigma_{j,t}=q_{j,t}/(1+\nu_{j,t})$ gives
\begin{align}
    \sigma_{j,1} &= q_{j,1}-\hat{\nu}_{j,1}\sigma,\\
    \sigma_{j,2} &= q_{j,2}-\hat{\nu}_{j,1}q_{j,1}+
    \left(\hat{\nu}_{j,1}^2-\hat{\nu}_{j,2}\right)\sigma.
\end{align}
Equivalently,
\begin{align}\label{eq:sigma-j2-general}
    \sigma_{j,2}
    &= \frac{\eta_2}{\gamma}-\frac{\eta_1}{\gamma}\hat{\gamma}_j+
    \sigma\hat{\gamma}_j^2+\frac{1-\gamma^{-1}}{\psi-1}\sigma_{c_j,2}
    -\hat{\nu}_{j,1}\left(\frac{\eta_1}{\gamma}-\sigma\hat{\gamma}_j\right)
    +\left(\hat{\nu}_{j,1}^2-\hat{\nu}_{j,2}\right)\sigma.
\end{align}
For an unconstrained investor, $\hat{\nu}_{j,1}=\hat{\nu}_{j,2}=0$, and this reduces to
\begin{align}\label{eq:sigma-j2-unconstrained}
    \sigma_{j,2}
    &= \frac{\eta_2}{\gamma}-\frac{\eta_1}{\gamma}\hat{\gamma}_j+
    \sigma\hat{\gamma}_j^2+\frac{1-\gamma^{-1}}{\psi-1}\sigma_{c_j,2}.
\end{align}

For a constrained investor, the multiplier is determined by the second-order expansion of the leverage constraint.
The leverage constraint is
\begin{equation}
    \|\sigma_j\| \leq \|\sigma\|(1+\hat{\sigma}\epsilon).
\end{equation}
Using
\begin{equation}
    \sigma_j=\sigma+\sigma_{j,1}\epsilon+\sigma_{j,2}\epsilon^2+\mathcal{O}(\epsilon^3),
\end{equation}
we obtain
\begin{align}
    \|\sigma_j\|
    &= \|\sigma\|\left[1+\frac{\sigma_{j,1}\sigma^\top}{\sigma\sigma^\top}\epsilon
    +\left(\frac{\sigma_{j,2}\sigma^\top}{\sigma\sigma^\top}
    +\frac{1}{2}\left(\frac{\sigma_{j,1}\sigma_{j,1}^\top}{\sigma\sigma^\top}
    -\left(\frac{\sigma_{j,1}\sigma^\top}{\sigma\sigma^\top}\right)^2\right)\right)\epsilon^2\right]
    +\mathcal{O}(\epsilon^3).
\end{align}
At first order, $\sigma_{j,1}$ is proportional to $\sigma$.
Therefore, the curvature term in the second line is zero.
If the constraint binds through second order and $\overline{\sigma}=\|\sigma\|(1+\hat{\sigma}\epsilon)$ has no second-order correction, then
\begin{equation}\label{eq:constraint-second-order-projection}
    \frac{\sigma_{j,2}\sigma^\top}{\sigma\sigma^\top}=0.
\end{equation}
Combining \eqref{eq:sigma-j2-general} and \eqref{eq:constraint-second-order-projection} yields
\begin{align}\label{eq:nu-j2-second-order}
    \hat{\nu}_{j,2}
    &= \frac{q_{j,2}\sigma^\top}{\sigma\sigma^\top}
    -\hat{\nu}_{j,1}\frac{q_{j,1}\sigma^\top}{\sigma\sigma^\top}
    +\hat{\nu}_{j,1}^2.
\end{align}
Substituting for $q_{j,1}$ and $q_{j,2}$ gives
\begin{align}\label{eq:nu-j2-expanded}
    \hat{\nu}_{j,2}
    &= \frac{\eta_2\sigma^\top}{\gamma\sigma\sigma^\top}
    -\hat{\gamma}_j\frac{\eta_1\sigma^\top}{\gamma\sigma\sigma^\top}
    +\hat{\gamma}_j^2
    +\frac{1-\gamma^{-1}}{\psi-1}\frac{\sigma_{c_j,2}\sigma^\top}{\sigma\sigma^\top}
    -\hat{\nu}_{j,1}\left(\frac{\eta_1\sigma^\top}{\gamma\sigma\sigma^\top}-\hat{\gamma}_j\right)
    +\hat{\nu}_{j,1}^2.
\end{align}

In the maintained two-active-investor case, investor $1$ is unconstrained and investor $2$ may be constrained.
Thus $\hat{\nu}_{1,1}=\hat{\nu}_{1,2}=0$.
If investor $2$ is constrained, then $\hat{\nu}_{2,1}=\hat{\nu}_2$ and $\hat{\nu}_{2,2}$ is given by \eqref{eq:nu-j2-expanded} with $j=2$.
If investor $2$ is slack, then $\hat{\nu}_{2,1}=\hat{\nu}_{2,2}=0$ and \eqref{eq:sigma-j2-unconstrained} applies to both active investors.

% \noindent\textbf{Step 3: Aggregate risk aversion.}\\
% The aggregate risk aversion can be written as
% \begin{equation}
%     \gamma_u(X,\epsilon) = \gamma \left[ 1 + \mathbb{E}^u\left[\hat{\gamma}_j\right] \epsilon - \var^u\left[\hat{\gamma}_j\right] \epsilon^2 \right] + \mathcal{O}(\epsilon^3),
% \end{equation}
% where
% \[\mathbb{E}^u\left[ \hat{\gamma}_j\right] \equiv\sum_{j\in\mathcal{J}^u} \frac{x_j}{x_u} \hat{\gamma}_j,\quad\text{and}\quad \var^u[\hat{\gamma}_j] \equiv \sum_{j\in\mathcal{J}^u} \frac{x_j}{x_u} \hat{\gamma}_j^2 - \left(\sum_{j\in\mathcal{J}^u} \frac{x_j}{x_u} \hat{\gamma}_j\right)^2.\]

\noindent\textbf{Step 3: Market price of risk.}\\
The market price of risk is given by
\begin{equation}
    \eta_t = \gamma_{a,t}\left[\chi_{a,t}\sigma_{R_Y,t} - h_{a,t}\right],
\end{equation}
where
\begin{equation}
    \chi_{a,t} \equiv \frac{1-x_0\overline{\alpha}_{p}}{x_a}
\end{equation}
is the aggregate risky exposure that must be absorbed by active investors per unit of active wealth,
\begin{equation}
    h_{a,t} \equiv \sum_{j=1}^2 \frac{x_j}{x_a}
    \frac{1-\gamma_j^{-1}}{\psi-1}\frac{\sigma_{c_j,t}}{1+\nu_{j,t}}
\end{equation}
is the average hedging demand of active investors, adjusted for leverage constraints, and
\begin{equation}
    \gamma_{a,t} \equiv \left[\sum_{j=1}^2 \frac{x_j}{x_a}
    \frac{1}{\gamma_j(1+\nu_{j,t})}\right]^{-1}
\end{equation}
is the average effective risk aversion of active investors.

We now expand each term to second order.
Let
\begin{equation}
    a_j \equiv \hat{\gamma}_j+\hat{\nu}_{j,1},
\end{equation}
and define active-wealth-weighted averages by
\begin{equation}
    \mathbb{E}_a[z_j]\equiv \sum_{j=1}^2 \frac{x_j}{x_a}z_j.
\end{equation}
Since
\begin{equation}
    \frac{1}{\gamma_j(1+\nu_{j,t})}
    = \frac{1}{\gamma}\left[1-a_j\epsilon+
    \left(a_j^2-\hat{\nu}_{j,2}-\hat{\gamma}_j\hat{\nu}_{j,1}\right)\epsilon^2\right]
    +\mathcal{O}(\epsilon^3),
\end{equation}
we have
\begin{equation}
    \gamma_{a,t}=\gamma\left[1+\Gamma_{a,1}\epsilon+\Gamma_{a,2}\epsilon^2\right]+\mathcal{O}(\epsilon^3),
\end{equation}
where
\begin{align}
    \Gamma_{a,1} &= \mathbb{E}_a[a_j],\\
    \Gamma_{a,2} &= \mathbb{E}_a[\hat{\nu}_{j,2}+\hat{\gamma}_j\hat{\nu}_{j,1}]
    -\operatorname{Var}_a(a_j).
\end{align}
Here
\begin{equation}
    \operatorname{Var}_a(a_j) \equiv \mathbb{E}_a[a_j^2]-\mathbb{E}_a[a_j]^2.
\end{equation}

Because $\overline{\alpha}_p=1+\hat{\alpha}_p\epsilon$ and $\kappa_p=\sigma_p=0$ in this step,
\begin{equation}
    \chi_{a,t}=1-\frac{x_0}{x_a}\hat{\alpha}_p\epsilon.
\end{equation}
Moreover, since $y_Y=1/p_Y$, the risk exposure of the aggregate claim is
\begin{equation}
    \sigma_{R_Y,t}=\sigma-\sigma_{y_Y,2}\epsilon^2+\mathcal{O}(\epsilon^3).
\end{equation}
Finally, since $\sigma_{c_j,t}=\sigma_{c_j,2}\epsilon^2+\mathcal{O}(\epsilon^3)$, hedging demand starts at second order:
\begin{equation}
    h_{a,t}=h_{a,2}\epsilon^2+\mathcal{O}(\epsilon^3),
\end{equation}
with
\begin{equation}
    h_{a,2}=\frac{1-\gamma^{-1}}{\psi-1}\mathbb{E}_a[\sigma_{c_j,2}].
\end{equation}
Therefore,
\begin{equation}
    \chi_{a,t}\sigma_{R_Y,t}-h_{a,t}
    = \sigma - \frac{x_0}{x_a}\hat{\alpha}_p\sigma\epsilon
    -\left(\sigma_{y_Y,2}+h_{a,2}\right)\epsilon^2
    +\mathcal{O}(\epsilon^3).
\end{equation}

Combining these expansions gives
\begin{align}
    \eta_0 &= \gamma\sigma,\\
    \eta_1 &= \gamma\sigma\left[\mathbb{E}_a[a_j]-\frac{x_0}{x_a}\hat{\alpha}_p\right],\\
    \eta_2 &= \gamma\left[
        \Gamma_{a,2}\sigma
        -\frac{x_0}{x_a}\hat{\alpha}_p\Gamma_{a,1}\sigma
        -\sigma_{y_Y,2}-h_{a,2}
    \right].
\end{align}
The first-order term coincides with the expression derived above because $a_j=\hat{\gamma}_j+\hat{\nu}_{j,1}$.
The second-order term depends on $\hat{\nu}_{j,2}$ through $\Gamma_{a,2}$.
Thus, when a leverage constraint binds at second order, \eqref{eq:nu-j2-expanded} and the expression for $\eta_2$ form a linear system for $\eta_2$ and the second-order multiplier.

% The market price of risk is given by
% \begin{small}
% 	\begin{equation}
% 	\theta(X,\epsilon) =  \frac{\gamma_u(X,\epsilon)}{x_{u}} \left[ (1-(1+\hat{\alpha}_p \epsilon) x_{0,t})\|\sigma - \sigma_{y}(X,\epsilon)\| - \|\sigma\|(1 + \hat{\sigma}\epsilon) x_{c,t}+ \sum_{j\in \mathcal{J}_t^u} x_{j,t} \frac{1-\gamma_j^{-1}}{1-\psi} \sigma_{c_j}(X,\epsilon) \frac{\sigma^\top}{\|\sigma\|}\right],
% 	\end{equation}	
% \end{small}

% The second-order term is then given by
% \begin{align}
%     \theta_2(X) &=  \frac{\gamma_{u,0}(X)}{x_{u}} \left[ -(1-x_{0,t})\sigma_{y,2}(X)\frac{\sigma^\top}{\|\sigma\|} + \sum_{j\in \mathcal{J}_t^u} x_{j,t} \frac{1-\gamma^{-1}}{1-\psi} \sigma_{c_j,2}(X) \frac{\sigma^\top}{\|\sigma\|}\right] + \nonumber\\
%     &+ \frac{\gamma_{u,1}(X)}{x_{u}} \|\sigma\|\left[ - \hat{\alpha}_p x_{0}  -\hat{\sigma} x_{c}\right]+ \frac{\gamma_{u,2}(X)}{x_{u}} \|\sigma\| \left[ 1-x_{0} -  x_{c} \right].
% \end{align}

% The expression above can be written as
% \begin{align}
%     \theta_2(X) =  - \frac{(\gamma-1) x_c + 1-x_0}{1-x_0-x_c} \sigma_{y,2}(X) \frac{\sigma^\top}{\|\sigma\|} - \gamma \|\sigma\| \mathbb{E}^u[\hat{\gamma}_j] \frac{ \hat{\alpha}_p x_0+ \hat{\sigma} x_c }{1-x_0-x_c} - \gamma \|\sigma\| \var^u[\hat{\gamma}_j].
% \end{align}
% using the fact that $\sigma_{c_j,2}(X) = (1-\psi)\sigma_{y,2}$.

\noindent\textbf{Step 4: Risk premium.}\\

The risk premium on the aggregate claim is given by
\begin{equation}
    \pi_{Y,t} = \sigma_{R_Y,t} \eta_t^\top.
\end{equation}

Using that $\sigma_{R_Y,1}=0$, expanding the risk premium to second order gives
\begin{equation}
    \pi_{Y,2} = \sigma_{R_Y,2} \eta_0^\top + \sigma \eta_2^\top = -\sigma_{y_Y,2} \eta_0^\top + \sigma \eta_2^\top.
\end{equation}
where we used $\sigma_{R_Y,t} = \sigma - \sigma_{y_Y,2}\epsilon^2 + \mathcal{O}(\epsilon^3)$.

% Step 5: Interest rate
\noindent\textbf{Step 5: Interest rate.} 
\begin{align}
    r_t &= \rho 
    + \psi^{-1} \left(\mu_{P_Y,t} + \xi_t \right)
    + (1-\psi^{-1}) \sum_{j=0}^2 x_{j,t} 
        \frac{\gamma_j}{2} \|\sigma_{j,t}\|^2
    - \pi_{Y,t},
\end{align}
where $\mu_{P_Y,t}$ is the drift of the aggregate claim value.

Since $P_Y=p_YY$, it satisfies
\begin{equation}
    \mu_{P_Y,t}=\mu+\mu_{p_Y,t}+\sigma\sigma_{p_Y,t}^\top.
\end{equation}
At second order, working with the dividend yield implies
\begin{equation}
    \mu_{P_Y,2}= -\mu_{y_Y,2}-\sigma\sigma_{y_Y,2}^\top,
\end{equation}
up to terms of order higher than $\epsilon^2$, since $\sigma_{y_Y,t}=\mathcal{O}(\epsilon^2)$ and $\sigma_{p_Y,2}=-\sigma_{y_Y,2}$.
Moreover, the intertemporal-hedging term satisfies
\begin{equation}
    \xi_2 = \sum_{j=0}^2 x_j\left[\mu_{c_j,2} + (1-\gamma)\sigma_{c_j,2}\sigma^\top\right],
\end{equation}
where terms involving $\|\sigma_{c_j,t}\|^2$ are fourth order and therefore do not enter at order $\epsilon^2$.

Expanding the risk-compensation term gives
\begin{align}
    \left[\sum_{j=0}^2 x_j\frac{\gamma_j}{2}\|\sigma_j\|^2\right]_2
    &= \gamma \sum_{j=0}^2 x_j \left[
        \sigma_{j,2}\sigma^\top
        + \frac{1}{2}\sigma_{j,1}\sigma_{j,1}^\top
        + \hat{\gamma}_j \sigma_{j,1}\sigma^\top
    \right].
\end{align}
Therefore, the second-order correction to the interest rate is
\begin{align}\label{eq:r2-second-order}
    r_2
    &= -\psi^{-1}\left(\mu_{y_Y,2}+\sigma\sigma_{y_Y,2}^\top\right)
    + \psi^{-1}\sum_{j=0}^2 x_j\left[\mu_{c_j,2} + (1-\gamma)\sigma_{c_j,2}\sigma^\top\right] \nonumber\\
    &\quad + (1-\psi^{-1})\gamma \sum_{j=0}^2 x_j \left[
        \sigma_{j,2}\sigma^\top
        + \frac{1}{2}\sigma_{j,1}\sigma_{j,1}^\top
        + \hat{\gamma}_j \sigma_{j,1}\sigma^\top
    \right]
    - \pi_{Y,2}.
\end{align}
Using market clearing for the risky claim,
\begin{equation}
    \sum_{j=0}^2 x_j \sigma_{j,2}=\sigma_{R_Y,2}=-\sigma_{y_Y,2},
\end{equation}
we can also write
\begin{align}\label{eq:r2-second-order-simplified}
    r_2
    &= -\psi^{-1}\left(\mu_{y_Y,2}+\sigma\sigma_{y_Y,2}^\top\right)
    + \psi^{-1}\sum_{j=0}^2 x_j\left[\mu_{c_j,2} + (1-\gamma)\sigma_{c_j,2}\sigma^\top\right] \nonumber\\
    &\quad + (1-\psi^{-1})\gamma \left[-\sigma_{y_Y,2}\sigma^\top
        + \sum_{j=0}^2 x_j \left(
        \frac{1}{2}\sigma_{j,1}\sigma_{j,1}^\top
        + \hat{\gamma}_j \sigma_{j,1}\sigma^\top
    \right)\right]
    - \pi_{Y,2}.
\end{align}
This expression separates the second-order interest-rate correction into the aggregate-claim drift term, the intertemporal-hedging term, the second-order change in average risk compensation, and the second-order risk premium.

% The interest rate is given by
% \begin{align}
%      r(X,\epsilon)   &  =  \rho - \theta(X,\epsilon) \|\sigma_{R}(X,\epsilon)\| +\psi^{-1} (\mu - \mu_{y}(X,\epsilon)- \sigma_y \sigma_{R}(X,\epsilon)^\top)\nonumber\\
%      &+ \left(1-\psi^{-1}\right)  \sum_{j=0}^J x_{j}\frac{\gamma_j}{2} \sigma_{j}^2(X,\epsilon)  \nonumber\\
%     &+\psi^{-1}\sum_{j=0}^J x_{j,t} \left[\mu_{c_j}(X,\epsilon) + (1-\gamma_j)  \sigma_{c_j}(X,\epsilon) \frac{\sigma_{R,t}^\top}{\|\sigma_R\|} \sigma_{j}(X,\epsilon) + \frac{\psi - \gamma_j}{1-\psi} \frac{\|\sigma_{c_j,t}(X,\epsilon)\|^2}{2}\right].
% \end{align}

% The second-order term is given by
% \begin{align}
%     r_2(X) &= - \theta_2(X) \|\sigma\| + \theta_0(X) \sigma_{y,2}(X) \frac{\sigma^\top}{\|\sigma\|}-\psi^{-1} (\mu_{y,2}(X) + \sigma_{y,2}(X)\sigma^\top)\nonumber\\
%     &+ \left(1-\psi^{-1}\right) \sum_{j=0}^J x_{j} \gamma\|\sigma\| \left[ \sigma_{j,2}(X) + \frac{\sigma_{j,1}^2(X)}{2 \|\sigma\|} + \hat{\gamma}_j \sigma_{j,1}(X) \right]\nonumber\\
%     &+\psi^{-1} \sum_{j=0}^J x_{j}\left[ \mu_{c_j,2}(X) + (1-\gamma) \sigma_{c_j,2}(X) \sigma^\top \right].
% \end{align}

\noindent\textbf{Step 6: Dividend yield.}\\
The dividend yield is given by
\begin{equation}
    y_Y = r + \eta (\sigma + \sigma_{p_Y})^\top - \left(\mu + \mu_{p_Y} + \sigma \sigma_{p_Y}^\top\right),        
\end{equation}
where $\sigma_{p_Y}$ and $\mu_{p_Y}$ denote the diffusion and drift of the price-dividend ratio.
Since $y_Y=1/p_Y$, we have
\begin{equation}
    \sigma_{p_Y,2}=-\sigma_{y_Y,2}, \qquad \mu_{p_Y,2}=-\mu_{y_Y,2},
\end{equation}
up to terms of order higher than $\epsilon^2$.
Expanding the dividend yield to second order gives
\begin{align}\label{eq:y2-pricing-condition}
    y_{Y,2}
    &= r_2 + \sigma \eta_2^\top - \eta_0\sigma_{y_Y,2}^\top
    + \mu_{y_Y,2} + \sigma\sigma_{y_Y,2}^\top \\
    &= r_2 + \pi_{Y,2} + \mu_{y_Y,2} + \sigma\sigma_{y_Y,2}^\top,
\end{align}
where the second equality uses the expression for $\pi_{Y,2}$ from Step 4.
Substituting the expression for $r_2$ from \eqref{eq:r2-second-order-simplified}, we obtain
% Step 6: Dividend yield
\begin{align}\label{eq:y2-second-order-general}
    y_{Y,2}
    &= (1-\psi^{-1})\mu_{y_Y,2}
    +(1-\psi^{-1})\sigma\sigma_{y_Y,2}^\top
    + \psi^{-1}\sum_{j=0}^2 x_j\left[\mu_{c_j,2} + (1-\gamma)\sigma_{c_j,2}\sigma^\top\right] \nonumber\\
    &\quad + (1-\psi^{-1})\gamma \left[-\sigma_{y_Y,2}\sigma^\top
        + \sum_{j=0}^2 x_j \left(
        \frac{1}{2}\sigma_{j,1}\sigma_{j,1}^\top
        + \hat{\gamma}_j \sigma_{j,1}\sigma^\top
    \right)\right].
\end{align}
From Step 1, the first-order corrections imply
\begin{equation}
    \mu_{c_j,2}=(1-\psi)\mu_{y_Y,2}, \qquad
    \sigma_{c_j,2}=(1-\psi)\sigma_{y_Y,2}.
\end{equation}
Using $\sum_{j=0}^2 x_j=1$, the terms involving $\mu_{y_Y,2}$ cancel.
The terms involving $\sigma_{y_Y,2}\sigma^\top$ also cancel.
Therefore,
\begin{equation}\label{eq:y2-second-order}
    y_{Y,2}
    = (1-\psi^{-1})\gamma
    \sum_{j=0}^2 x_j \left(
        \frac{1}{2}\sigma_{j,1}\sigma_{j,1}^\top
        + \hat{\gamma}_j \sigma_{j,1}\sigma^\top
    \right).
\end{equation}
Equivalently,
\begin{equation}\label{eq:y2-second-order-scalar-exposure}
    y_{Y,2}
    = (1-\psi^{-1})\gamma\|\sigma\|^2
    \sum_{j=0}^2 x_j \left(
        \frac{1}{2}m_{j,1}^2 + \hat{\gamma}_j m_{j,1}
    \right),
\end{equation}
where $\sigma_{j,1}=m_{j,1}\sigma$.
This expression shows that the second-order correction to the dividend yield depends only on first-order risk reallocations and preference heterogeneity.

The first-order exposure coefficients are affine in $\hat{\alpha}_p$.
Write
\begin{equation}
    m_{j,1}=\bar{m}_j-\ell_j\hat{\alpha}_p.
\end{equation}
For passive investors,
\begin{equation}
    \bar{m}_0=0,
    \qquad
    \ell_0=-1,
\end{equation}
so that $m_{0,1}=\hat{\alpha}_p$.
For active investors,
\begin{align}
    \bar{m}_1 &= \frac{x_2}{x_a}(\hat{\gamma}_2-\hat{\gamma}_1+\hat{\nu}_2),
    &
    \ell_1 &= \frac{x_0}{x_a},\\
    \bar{m}_2 &= \frac{x_1}{x_a}(\hat{\gamma}_1-\hat{\gamma}_2-\hat{\nu}_2),
    &
    \ell_2 &= \frac{x_0}{x_a}.
\end{align}
These definitions satisfy the first-order market-clearing condition
\begin{equation}
    \sum_{j=0}^2 x_j m_{j,1}=0.
\end{equation}

Substituting $m_{j,1}=\bar{m}_j-\ell_j\hat{\alpha}_p$ into \eqref{eq:y2-second-order-scalar-exposure}, the solution for $y_{Y,2}$ is quadratic in $\hat{\alpha}_p$:
\begin{equation}\label{eq:y2-quadratic-alpha-fixed}
    y_{Y,2}(x,\hat{\alpha}_p)
    = a_Y^{0}(x)+b_Y^{0}(x)\hat{\alpha}_p+d_Y^{0}(x)\hat{\alpha}_p^2,
\end{equation}
where
\begin{align}
    a_Y^{0}(x)
    &= (1-\psi^{-1})\gamma\|\sigma\|^2
    \sum_{j=0}^2 x_j\left(
        \frac{1}{2}\bar{m}_j^2+\hat{\gamma}_j\bar{m}_j
    \right),\\
    b_Y^{0}(x)
    &= -(1-\psi^{-1})\gamma\|\sigma\|^2
    \sum_{j=0}^2 x_j\ell_j\left(\bar{m}_j+\hat{\gamma}_j\right),\\
    d_Y^{0}(x)
    &= \frac{1}{2}(1-\psi^{-1})\gamma\|\sigma\|^2
    \sum_{j=0}^2 x_j\ell_j^2.
\end{align}
Using the values of $\ell_j$, the quadratic coefficient can be simplified to
\begin{equation}\label{eq:dY0-fixed-alpha-simplified}
    d_Y^{0}(x)
    = \frac{1}{2}(1-\psi^{-1})\gamma\|\sigma\|^2
    \left[x_0+\frac{x_0^2}{x_a}\right]
    = \frac{1}{2}(1-\psi^{-1})\gamma\|\sigma\|^2\frac{x_0}{x_a}.
\end{equation}

The second-order price-dividend-ratio correction follows from the change of variables above:
\begin{equation}
    p_{Y,2}=p_{Y,0}^3 y_{Y,1}^2-p_{Y,0}^2 y_{Y,2}.
\end{equation}


\noindent\textbf{Step 7: Consumption-wealth ratio.}\\
The consumption-wealth ratio is given by
\begin{equation}
    c_{j,t} = \psi \rho 
    + (1-\psi) \left[ r_t + \eta_{t}\,\sigma_{j,t}^\top - \frac{\gamma_j}{2} \|\sigma_{j,t}\|^2 \right] 
    + \mu_{c_j,t} 
    + (1-\gamma_j) \sigma_{c_j,t} \sigma_{j,t}^\top 
    + \frac{\psi - \gamma_j}{1-\psi} \frac{\|\sigma_{c_j,t}\|^2}{2}.
\end{equation}

We expand each term to second order.
Using
\begin{align}
    r_t &= r_0 + r_1\epsilon + r_2\epsilon^2 + \mathcal{O}(\epsilon^3),\\
    \eta_t &= \eta_0 + \eta_1\epsilon + \eta_2\epsilon^2 + \mathcal{O}(\epsilon^3),\\
    \sigma_{j,t} &= \sigma + \sigma_{j,1}\epsilon + \sigma_{j,2}\epsilon^2 + \mathcal{O}(\epsilon^3),
\end{align}
and
\begin{align}
    \gamma_j &= \gamma(1+\hat{\gamma}_j\epsilon), \\
    \sigma_{c_j,t} &= \sigma_{c_j,2}\epsilon^2 + \mathcal{O}(\epsilon^3), \\
    \mu_{c_j,t} &= \mu_{c_j,2}\epsilon^2 + \mathcal{O}(\epsilon^3),
\end{align}
we collect second-order terms.

The first-order term is
\begin{equation}
    c_{j,1}
    =
    (1-\psi)
    \left[
        r_1
        + \eta_1\sigma^\top
        - \frac{\gamma}{2}\|\sigma\|^2 \hat{\gamma}_j
    \right],
\end{equation}
which matches the previous result.

Before simplifying, the second-order term is
\begin{align}
    c_{j,2} &= (1-\psi)\Big[
        r_2
        + \eta_2\sigma^\top
        + \eta_1\sigma_{j,1}^\top
        + \eta_0\sigma_{j,2}^\top
        - \gamma\sigma\sigma_{j,2}^\top
        - \frac{\gamma}{2}\sigma_{j,1}\sigma_{j,1}^\top
        - \gamma\hat{\gamma}_j\sigma\sigma_{j,1}^\top
    \Big] \nonumber\\
    &\quad + \mu_{c_j,2}
    + (1-\gamma)\sigma_{c_j,2}\sigma^\top.
\end{align}
Using $\eta_0=\gamma\sigma$, the terms involving $\sigma_{j,2}$ cancel.
Therefore,
\begin{align}
    c_{j,2} &= (1-\psi)\Big[
        r_2
        + \eta_2\sigma^\top
        + \eta_1\sigma_{j,1}^\top
        - \frac{\gamma}{2}\sigma_{j,1}\sigma_{j,1}^\top
        - \gamma\hat{\gamma}_j\sigma\sigma_{j,1}^\top
    \Big] \nonumber\\
    &\quad + \mu_{c_j,2}
    + (1-\gamma)\sigma_{c_j,2}\sigma^\top.
\end{align}

Substituting the expressions for $r_2$, $\eta_2$, and $\sigma_{j,2}$ derived in Steps 2--5 yields a closed-form expression for $c_{j,2}$ in terms of first-order reallocations and heterogeneity.

\noindent\textbf{Step 8: The dividend claim.}\\
The normalized price of the dividend claim, $p_D=P_D/Y$, satisfies
\begin{equation}
    r_t + (\sigma+\sigma_{p_D,t})\eta_t^\top
    = \frac{s_t}{p_{D,t}} + \mu + \mu_{p_D,t} + \sigma\sigma_{p_D,t}^\top.
\end{equation}
Equivalently, after multiplying by $p_{D,t}$,
\begin{equation}\label{eq:pD-pricing-condition}
    \left[r_t+\sigma\eta_t^\top-\mu\right]p_{D,t}
    = s_t + p_{D,x,t}\mu_{x,t}+p_{D,s,t}\mu_{s,t}
    +\frac{1}{2}\sum_{k=1}^3 \Sigma_{D,k,t}^\top p_{D,XX,t}\Sigma_{D,k,t}
    +(\sigma-\eta_t)\left(p_{D,x,t}\sigma_{x,t}+p_{D,s,t}\sigma_{s,t}\right)^\top,
\end{equation}
where $X=(x,s)$, $p_{D,XX}$ is the Hessian with respect to $(x,s)$, and $\Sigma_{D,k,t}$ is the $k$-th column of the diffusion matrix of $(x,s)$.

We now introduce the local coordinate
\begin{equation}
    s=\bar{s}+\epsilon\hat{s}.
\end{equation}
The function $p_D$ is a function of the physical state $(x,s)$, but its perturbation coefficients are functions of $(x,\hat{s})$:
\begin{equation}\label{eq:pD-expansion-local-s}
    p_D(x,s;\epsilon)
    = p_{D,0}+p_{D,1}(x,\hat{s})\epsilon+p_{D,2}(x,\hat{s})\epsilon^2
    +\mathcal{O}(\epsilon^3).
\end{equation}
Because $\hat{s}=(s-\bar{s})/\epsilon$, derivatives with respect to the physical state $s$ satisfy
\begin{align}
    p_{D,s}
    &= p_{D,1,\hat{s}}+\epsilon p_{D,2,\hat{s}}+\mathcal{O}(\epsilon^2),\\
    p_{D,ss}
    &= \epsilon^{-1}p_{D,1,\hat{s}\hat{s}}+p_{D,2,\hat{s}\hat{s}}+\mathcal{O}(\epsilon).
\end{align}
In the first-order solution below, $p_{D,1}$ is linear in $\hat{s}$, so $p_{D,1,\hat{s}\hat{s}}=0$ and the second derivative is well behaved:
\begin{equation}
    p_{D,ss}=p_{D,2,\hat{s}\hat{s}}+\mathcal{O}(\epsilon).
\end{equation}
Similarly,
\begin{equation}
    p_{D,x}=\epsilon p_{D,1,x}+\epsilon^2p_{D,2,x}+\mathcal{O}(\epsilon^3),
    \qquad
    p_{D,xs}=p_{D,1,x\hat{s}}+\epsilon p_{D,2,x\hat{s}}+\mathcal{O}(\epsilon^2).
\end{equation}

The dividend-share process implies
\begin{align}
    \mu_{s,t} &= -\kappa_s(s_t-\bar{s}) = -\kappa_s\hat{s}\epsilon,\\
    \sigma_{s,t} &= s_t(1-s_t)\sigma_s\epsilon 
    = \bar{s}(1-\bar{s})\sigma_s\epsilon
    +(1-2\bar{s})\hat{s}\sigma_s\epsilon^2
    +\mathcal{O}(\epsilon^3).
\end{align}

The wealth-share dynamics satisfy
\begin{equation}
    \mu_{x,t}=\mu_{x,1}\epsilon+\mathcal{O}(\epsilon^2),
    \qquad
    \sigma_{x,t}=\sigma_{x,1}\epsilon+\mathcal{O}(\epsilon^2).
\end{equation}

Define
\begin{equation}
    A_t\equiv r_t+\sigma\eta_t^\top-\mu
    =\delta+A_1\epsilon+A_2\epsilon^2+\mathcal{O}(\epsilon^3),
\end{equation}
where
\begin{equation}
    \delta=r_0+\sigma\eta_0^\top-\mu,
    \qquad
    A_1=r_1+\sigma\eta_1^\top,
    \qquad
    A_2=r_2+\sigma\eta_2^\top.
\end{equation}
The payoff satisfies $s_t=\bar{s}+\hat{s}\epsilon$.

The zeroth-order condition is
\begin{equation}
    \delta p_{D,0}=\bar{s},
    \qquad\text{so}\qquad
    p_{D,0}=\frac{\bar{s}}{\delta}.
\end{equation}

At first order, the only state-derivative terms come from the drift of $s$ and from the order-$\epsilon$ diffusion of $s$ evaluated at $\bar{s}$.
The first-order correction satisfies
\begin{equation}\label{eq:pD-first-order-PDE}
    \delta p_{D,1}+A_1p_{D,0}
    =\hat{s}-\kappa_s\hat{s}p_{D,1,\hat{s}}
    +(\sigma-\eta_0)\left[p_{D,1,\hat{s}}\bar{s}(1-\bar{s})\sigma_s\right]^\top.
\end{equation}
This equation is the same first-order pricing equation derived above.
Its solution is
\begin{equation}
    p_{D,1}
    = \frac{\hat{s}}{\delta+\kappa_s}
    -(r_1+\pi_{Y,1})\frac{\bar{s}}{\delta^2}
    -\frac{\gamma-1}{\delta(\delta+\kappa_s)}\bar{s}(1-\bar{s})\sigma\sigma_s^\top.
\end{equation}

At second order, using $\mu_{s,t}=-\kappa_s\hat{s}\epsilon$ and $p_{D,s}=p_{D,1,\hat{s}}+\epsilon p_{D,2,\hat{s}}+\mathcal{O}(\epsilon^2)$, the correction satisfies
\begin{align}\label{eq:pD-second-order-PDE}
    \delta p_{D,2}+A_1p_{D,1}+A_2p_{D,0}
    &= -\kappa_s\hat{s}p_{D,2,\hat{s}}
    +(\sigma-\eta_0)\left[p_{D,2,\hat{s}}\bar{s}(1-\bar{s})\sigma_s\right]^\top \nonumber\\
    &\quad +p_{D,1,x}\mu_{x,1}
    +(\sigma-\eta_0)\left[p_{D,1,x}\sigma_{x,1}\right]^\top
    -\eta_1\left[p_{D,1,\hat{s}}\bar{s}(1-\bar{s})\sigma_s\right]^\top \nonumber\\
    &\quad +(\sigma-\eta_0)\left[p_{D,1,\hat{s}}(1-2\bar{s})\hat{s}\sigma_s\right]^\top
    +\frac{1}{2}p_{D,2,\hat{s}\hat{s}}\bar{s}^2(1-\bar{s})^2\sigma_s\sigma_s^\top.
\end{align}
The potential second-order Hessian terms involving $p_{D,1,x\hat{s}}$ and $p_{D,1,\hat{s}\hat{s}}$ vanish because $p_{D,1}$ is additively separable in $(x,\hat{s})$ and linear in $\hat{s}$.
Thus, after these simplifications, the right-hand side of \eqref{eq:pD-second-order-PDE} is at most affine in $\hat{s}$, apart from the operator applied to $p_{D,2}$.
We can solve \eqref{eq:pD-second-order-PDE} by matching coefficients in $\hat{s}$.
Let
\begin{equation}
    B_s \equiv \bar{s}(1-\bar{s}),
    \qquad
    \Lambda_s \equiv (\sigma-\eta_0)\sigma_s^\top=(1-\gamma)\sigma\sigma_s^\top.
\end{equation}
Write the first-order solution as
\begin{equation}
    p_{D,1}(x,\hat{s})=a_1(x,\hat{\alpha}_p)+b_1\hat{s},
\end{equation}
where
\begin{equation}
    b_1=\frac{1}{\delta+\kappa_s},
\end{equation}
and
\begin{equation}
    a_1(x,\hat{\alpha}_p)
    =-(r_1+\pi_{Y,1})\frac{\bar{s}}{\delta^2}
    -\frac{\gamma-1}{\delta(\delta+\kappa_s)}B_s\sigma\sigma_s^\top.
\end{equation}
Equivalently, using $A_1=r_1+\sigma\eta_1^\top=r_1+\pi_{Y,1}$,
\begin{equation}
    a_1(x,\hat{\alpha}_p)
    =-\frac{A_1p_{D,0}}{\delta}+\frac{\Lambda_s B_s}{\delta(\delta+\kappa_s)}.
\end{equation}

Conjecture that the second-order correction is quadratic in $\hat{s}$:
\begin{equation}
    p_{D,2}(x,\hat{s})=a_2(x,\hat{\alpha}_p)+b_2(x,\hat{\alpha}_p)\hat{s}+c_2(x,\hat{\alpha}_p)\hat{s}^2.
\end{equation}
Matching the coefficient on $\hat{s}^2$ gives
\begin{equation}
    (\delta+\kappa_s)c_2=0,
\end{equation}
so
\begin{equation}
    c_2=0.
\end{equation}
Thus, in this benchmark case, the second-order correction is affine in $\hat{s}$.

Matching the coefficient on $\hat{s}$ gives
\begin{equation}
    (\delta+\kappa_s)b_2
    =-A_1b_1+\Lambda_s(1-2\bar{s})b_1,
\end{equation}
so
\begin{equation}
    b_2(x,\hat{\alpha}_p)
    =\frac{\Lambda_s(1-2\bar{s})-A_1}{(\delta+\kappa_s)^2}.
\end{equation}
Finally, matching the constant term gives
\begin{equation}
    \delta a_2
    =-A_1a_1-A_2p_{D,0}+\Lambda_sB_sb_2+\mathcal{M}_x
    -\eta_1\left(b_1B_s\sigma_s\right)^\top,
\end{equation}
where
\begin{equation}
    \mathcal{M}_x
    \equiv a_{1,x}\mu_{x,1}
    +(\sigma-\eta_0)\left[a_{1,x}\sigma_{x,1}\right]^\top.
\end{equation}
Therefore,
\begin{equation}
    a_2(x,\hat{\alpha}_p)
    =\frac{1}{\delta}\left[
        -A_1a_1-A_2p_{D,0}+\Lambda_sB_sb_2+\mathcal{M}_x
        -\eta_1\left(b_1B_s\sigma_s\right)^\top
    \right].
\end{equation}
Combining the coefficient restrictions, the second-order correction is
\begin{equation}
    p_{D,2}(x,\hat{s})=a_2(x,\hat{\alpha}_p)+b_2(x,\hat{\alpha}_p)\hat{s}.
\end{equation}
The term proportional to $\eta_1(b_1B_s\sigma_s)^\top$ captures the interaction between first-order changes in the market price of risk and the first-order dividend-share exposure of the dividend claim.

\noindent\textbf{Step 9: Pricing impact for the dividend claim.}\\
We now compute the derivative of the second-order dividend-claim price with respect to the passive investors' portfolio share $\hat{\alpha}_p$.
The coefficient representation from Step 8 gives
\begin{equation}
    p_{D,2}(x,\hat{s})=a_2(x,\hat{\alpha}_p)+b_2(x,\hat{\alpha}_p)\hat{s},
\end{equation}
where
\begin{equation}
    b_2(x,\hat{\alpha}_p)
    =\frac{\Lambda_s(1-2\bar{s})-A_1}{(\delta+\kappa_s)^2}
\end{equation}
and
\begin{equation}
    a_2(x,\hat{\alpha}_p)
    =\frac{1}{\delta}\left[
        -A_1a_1-A_2p_{D,0}+\Lambda_sB_sb_2+\mathcal{M}_x
        -\eta_1\left(b_1B_s\sigma_s\right)^\top
    \right].
\end{equation}
Recall that
\begin{equation}
    A_1=r_1+\sigma\eta_1^\top=r_1+\pi_{Y,1}.
\end{equation}
The first-order aggregate-price term is independent of $\hat{\alpha}_p$, so
\begin{equation}
    \frac{\partial A_1}{\partial \hat{\alpha}_p}=0.
\end{equation}
It follows that $a_1$ and $b_2$ are also independent of $\hat{\alpha}_p$, except through the wealth-share dynamics term $\mathcal{M}_x$ in $a_2$.
Therefore,
\begin{equation}
    \frac{\partial p_{D,2}}{\partial \hat{\alpha}_p}
    =\frac{\partial a_2}{\partial \hat{\alpha}_p},
\end{equation}
and
\begin{equation}\label{eq:pD2-alpha-derivative-coefficients}
    \frac{\partial p_{D,2}}{\partial \hat{\alpha}_p}
    =\frac{1}{\delta}\left[
        -p_{D,0}\frac{\partial A_2}{\partial \hat{\alpha}_p}
        +\frac{\partial \mathcal{M}_x}{\partial \hat{\alpha}_p}
        -\frac{\partial \eta_1}{\partial \hat{\alpha}_p}
        \left(b_1B_s\sigma_s\right)^\top
    \right].
\end{equation}
Using $p_{D,0}=\bar{s}/\delta$ and $b_1=1/(\delta+\kappa_s)$, this can be written as
\begin{equation}\label{eq:pD2-alpha-derivative-final}
    \frac{\partial p_{D,2}}{\partial \hat{\alpha}_p}
    =-\frac{\bar{s}}{\delta^2}\frac{\partial A_2}{\partial \hat{\alpha}_p}
    -\frac{B_s}{\delta(\delta+\kappa_s)}
        \frac{\partial \eta_1}{\partial \hat{\alpha}_p}\sigma_s^\top
    +\frac{1}{\delta}\frac{\partial \mathcal{M}_x}{\partial \hat{\alpha}_p}.
\end{equation}
The first term is the aggregate-price-impact channel.
The second term is specific to the dividend claim: passive demand changes the first-order market price of risk, and the dividend claim loads on dividend-share risk through $b_1B_s\sigma_s$.
The last term captures the indirect effect operating through first-order wealth-share dynamics.
Since
\begin{equation}
    \mathcal{M}_x
    =a_{1,x}\mu_{x,1}+ (\sigma-\eta_0)\left[a_{1,x}\sigma_{x,1}\right]^\top,
\end{equation}
and $a_1$ is independent of $\hat{\alpha}_p$, this indirect term is
%
\begin{align}
    \frac{\partial \mathcal{M}_x}{\partial \hat{\alpha}_p}
    &=a_{1,x}\frac{\partial \mu_{x,1}}{\partial \hat{\alpha}_p}
    +(\sigma-\eta_0)\left[a_{1,x}\frac{\partial \sigma_{x,1}}{\partial \hat{\alpha}_p}\right]^\top.
\end{align}
This term is equal to zero.
To see this, recall that the first-order diffusion of wealth shares is
\begin{equation}
    \sigma_{x_j,1}=x_j\sigma_{j,1},
\end{equation}
and the first-order drift of wealth shares is
\begin{equation}
    \mu_{x_j,1}=x_j\left[(\eta_0-\sigma)\sigma_{j,1}^\top-c_{j,1}+y_{Y,1}\right].
\end{equation}
Since $c_{j,1}$ and $y_{Y,1}$ are independent of $\hat{\alpha}_p$, differentiating with respect to $\hat{\alpha}_p$ gives
\begin{equation}
    \frac{\partial \mu_{x_j,1}}{\partial \hat{\alpha}_p}
    =x_j(\eta_0-\sigma)\frac{\partial \sigma_{j,1}^\top}{\partial \hat{\alpha}_p},
    \qquad
    \frac{\partial \sigma_{x_j,1}}{\partial \hat{\alpha}_p}
    =x_j\frac{\partial \sigma_{j,1}}{\partial \hat{\alpha}_p}.
\end{equation}
Therefore,
\begin{align}
    \frac{\partial \mathcal{M}_x}{\partial \hat{\alpha}_p}
    &=\sum_{j=1}^2 a_{1,x_j}x_j(\eta_0-\sigma)\frac{\partial \sigma_{j,1}^\top}{\partial \hat{\alpha}_p}
    +(\sigma-\eta_0)\sum_{j=1}^2 a_{1,x_j}x_j\frac{\partial \sigma_{j,1}^\top}{\partial \hat{\alpha}_p} \\
    &=0.
\end{align}
Thus the indirect wealth-share-dynamics channel drops out of the dividend-claim price impact.
Finally, since
\begin{equation}
    p_D(x,s;\epsilon)=p_{D,0}+p_{D,1}(x,\hat{s})\epsilon+p_{D,2}(x,\hat{s})\epsilon^2+\mathcal{O}(\epsilon^3)
\end{equation}
and $p_{D,1}$ is independent of $\hat{\alpha}_p$, the price impact of passive flows is
\begin{equation}
    \frac{\partial p_D(x,s;\epsilon)}{\partial F(x,s;\epsilon)}
    =\frac{\epsilon}{x_0}\frac{\partial p_{D,2}}{\partial \hat{\alpha}_p}
    +\mathcal{O}(\epsilon^2).
\end{equation}
Thus, the dividend-claim price impact is the sum of an aggregate-price-impact channel and a direct dividend-risk-exposure channel.


\subsection{Second-order correction: time-varying passive portfolio share}\label{app:prop_second_order}

We next derive the second-order correction in the case the passive portfolio share is time-varying. 
Now, $\overline{\alpha}_p$ becomes a state variable. 
The dynamics of $\overline{\alpha}_p = 1 + \hat{\alpha}_p \epsilon$ is given by
\begin{align}
    \mu_{\overline{\alpha}_p,t} &= \kappa_p (\hat{\alpha} - \hat{\alpha}_p) \epsilon, \\
    \sigma_{\overline{\alpha}_p,t} &= \sigma_p \epsilon. 
\end{align}

We consider the following expansion of $c_j(x, \overline{\alpha}_p;\epsilon)$:
\begin{equation}
    c_{j}(x, \overline{\alpha}_p;\epsilon) = c_{j,0} + c_{j,1}(x, \hat{\alpha}_p) \epsilon + c_{j,2}(x, \hat{\alpha}_p) \epsilon^2 + \mathcal{O}(\epsilon^3). 
\end{equation}

Because $\hat{\alpha}_p=(\overline{\alpha}_p-1)/\epsilon$, derivatives with respect to the physical state $\overline{\alpha}_p$ satisfy the same chain-rule structure as for $s$.
For a generic object $f(x,\overline{\alpha}_p;\epsilon)=f_0+f_1(x,\hat{\alpha}_p)\epsilon+f_2(x,\hat{\alpha}_p)\epsilon^2+\mathcal{O}(\epsilon^3)$, we have
\begin{align}
    f_{\overline{\alpha}_p}
    &= f_{1,\hat{\alpha}_p} + \epsilon f_{2,\hat{\alpha}_p} + \mathcal{O}(\epsilon^2),\\
    f_{\overline{\alpha}_p\overline{\alpha}_p}
    &= \epsilon^{-1} f_{1,\hat{\alpha}_p\hat{\alpha}_p} + f_{2,\hat{\alpha}_p\hat{\alpha}_p} + \mathcal{O}(\epsilon).
\end{align}

In our setting, $c_{j,1}$ and $y_{Y,1}$ are independent of $\hat{\alpha}_p$, so
\begin{equation}
    f_{1,\hat{\alpha}_p} = 0, \qquad f_{1,\hat{\alpha}_p\hat{\alpha}_p} = 0,
\end{equation}
and therefore
\begin{equation}
    f_{\overline{\alpha}_p} = \epsilon f_{2,\hat{\alpha}_p} + \mathcal{O}(\epsilon^2),
    \qquad
    f_{\overline{\alpha}_p\overline{\alpha}_p} = f_{2,\hat{\alpha}_p\hat{\alpha}_p} + \mathcal{O}(\epsilon).
\end{equation}

Similarly, cross-derivatives satisfy
\begin{equation}
    f_{x\overline{\alpha}_p} = f_{1,x\hat{\alpha}_p} + \epsilon f_{2,x\hat{\alpha}_p} + \mathcal{O}(\epsilon^2) = \epsilon f_{2,x\hat{\alpha}_p} + \mathcal{O}(\epsilon^2),
\end{equation}
since $f_{1}$ does not depend on $\hat{\alpha}_p$.

\noindent\textbf{Step 1: Drift and diffusion for $c_j$ and $y_Y$.}

We avoid repeating the derivation for the case in which $\hat{\alpha}_p$ is fixed.
For any second-order object $z_2$, write
\begin{equation}
    z_2=z_2^{0}+\Delta z_2,
\end{equation}
where $z_2^{0}$ denotes the second-order correction derived in the fixed-$\hat{\alpha}_p$ economy, and $\Delta z_2$ collects only the additional terms generated by the drift and diffusion of $\overline{\alpha}_p$.
By construction, $\Delta z_2$ vanishes when $\kappa_p=0$ and $\sigma_p=0_{1\times 3}$.

Because $c_{j,1}$ and $y_{Y,1}$ are independent of $\hat{\alpha}_p$, the diffusion terms coming from the time variation in $\overline{\alpha}_p$ are
\begin{align}
    \Delta \sigma_{c_j,2}
    &= \frac{c_{j,2,\hat{\alpha}_p}}{c_{j,0}}\sigma_p,\\
    \Delta \sigma_{y_Y,2}
    &= \frac{y_{Y,2,\hat{\alpha}_p}}{y_{Y,0}}\sigma_p.
\end{align}
Therefore,
\begin{align}
    \sigma_{c_j,2} &= \sigma_{c_j,2}^{0}+\frac{c_{j,2,\hat{\alpha}_p}}{c_{j,0}}\sigma_p,\\
    \sigma_{y_Y,2} &= \sigma_{y_Y,2}^{0}+\frac{y_{Y,2,\hat{\alpha}_p}}{y_{Y,0}}\sigma_p.
\end{align}

Similarly, the new drift terms generated by the time variation in $\overline{\alpha}_p$ are
\begin{align}
    \Delta \mu_{c_j,2}
    &= \frac{c_{j,2,\hat{\alpha}_p}}{c_{j,0}}\kappa_p(\hat{\alpha}-\hat{\alpha}_p)
    +\frac{1}{2}\frac{c_{j,2,\hat{\alpha}_p\hat{\alpha}_p}}{c_{j,0}}\sigma_p\sigma_p^\top,\\
    \Delta \mu_{y_Y,2}
    &= \frac{y_{Y,2,\hat{\alpha}_p}}{y_{Y,0}}\kappa_p(\hat{\alpha}-\hat{\alpha}_p)
    +\frac{1}{2}\frac{y_{Y,2,\hat{\alpha}_p\hat{\alpha}_p}}{y_{Y,0}}\sigma_p\sigma_p^\top.
\end{align}

Thus the only new second-order terms are proportional to $\kappa_p$ or $\sigma_p$.

\noindent\textbf{Step 2 (time-varying $\alpha_p$): Risk exposures of investors.}\\
Write
\begin{equation}
    \sigma_{j,2} = \sigma_{j,2}^{0} + \Delta\sigma_{j,2},
\end{equation}
where $\sigma_{j,2}^{0}$ is the second-order exposure from the fixed-$\hat{\alpha}_p$ economy and $\Delta\sigma_{j,2}$ collects only the new terms generated by the time variation in $\overline{\alpha}_p$.

The desired exposure remains
\begin{equation}
    q_{j,t}
    \equiv \frac{\eta_t}{\gamma_j}+\frac{1-\gamma_j^{-1}}{\psi-1}\,\sigma_{c_j,t}.
\end{equation}
Relative to the fixed-$\hat{\alpha}_p$ economy, the second-order desired exposure changes for two reasons: the market price of risk changes by $\Delta\eta_2$, and the consumption-wealth ratio has an additional diffusion induced by $\sigma_p$.
Thus,
\begin{equation}
    \Delta q_{j,2}
    =\frac{\Delta\eta_2}{\gamma}
    +\frac{1-\gamma^{-1}}{\psi-1}\Delta\sigma_{c_j,2}.
\end{equation}
Using Steps 1 and 2,
\begin{equation}
    \Delta q_{j,2}
    =\frac{\Delta\eta_2}{\gamma}
    +\frac{1-\gamma^{-1}}{\psi-1}\frac{c_{j,2,\hat{\alpha}_p}}{c_{j,0}}\sigma_p.
\end{equation}

For an unconstrained investor, the multiplier is zero, so the incremental second-order exposure is simply
\begin{equation}
    \Delta\sigma_{j,2}=\Delta q_{j,2}.
\end{equation}
Therefore,
\begin{equation}\label{eq:delta-sigma-j2-unconstrained-alpha-p}
    \Delta\sigma_{j,2}
    =\frac{\Delta\eta_2}{\gamma}
    +\frac{1-\gamma^{-1}}{\psi-1}\frac{c_{j,2,\hat{\alpha}_p}}{c_{j,0}}\sigma_p.
\end{equation}

For a constrained investor, the second-order multiplier adjusts to keep the constraint binding.
Writing
\begin{equation}
    \hat{\nu}_{j,2}=\hat{\nu}_{j,2}^{0}+\Delta\hat{\nu}_{j,2},
\end{equation}
the incremental exposure is
\begin{equation}
    \Delta\sigma_{j,2}
    =\Delta q_{j,2}-\Delta\hat{\nu}_{j,2}\sigma.
\end{equation}
The second-order constraint condition requires the component of $\Delta\sigma_{j,2}$ in the direction of $\sigma$ to vanish:
\begin{equation}
    \Delta\sigma_{j,2}\sigma^\top=0.
\end{equation}
Hence
\begin{equation}
    \Delta\hat{\nu}_{j,2}
    =\frac{\Delta q_{j,2}\sigma^\top}{\sigma\sigma^\top},
\end{equation}
and therefore
\begin{equation}\label{eq:delta-sigma-j2-constrained-alpha-p}
    \Delta\sigma_{j,2}
    =\Delta q_{j,2}
    -\frac{\Delta q_{j,2}\sigma^\top}{\sigma\sigma^\top}\sigma.
\end{equation}
Thus, for constrained investors, only the component of the new desired exposure orthogonal to $\sigma$ affects actual risk exposure; the component parallel to $\sigma$ is absorbed by the second-order multiplier.

Combining the results, the time-varying passive portfolio share affects second-order risk exposures through the new market-price-of-risk term $\Delta\eta_2$ and the new consumption-wealth diffusion term proportional to $\sigma_p$.

\noindent\textbf{Step 3 (time-varying $\alpha_p$): Market price of risk.}\\
Write
\begin{equation}
    \eta_2=\eta_2^{0}+\Delta\eta_2,
\end{equation}
where $\eta_2^{0}$ is the second-order market price of risk from the fixed-$\hat{\alpha}_p$ economy and $\Delta\eta_2$ collects only the new terms generated by the time variation in $\overline{\alpha}_p$.

Recall that
\begin{equation}
    \eta_t=\gamma_{a,t}\left[\chi_{a,t}\sigma_{R_Y,t}-h_{a,t}\right].
\end{equation}
Relative to the fixed-$\hat{\alpha}_p$ economy, there are three new second-order effects.
First, the aggregate claim has an additional diffusion induced by $\sigma_p$.
Second, active investors have additional hedging demand because their consumption-wealth ratios load on $\sigma_p$.
Third, if the leverage constraint binds for investor $2$, the second-order multiplier changes, which changes effective aggregate risk aversion.

From Step 1,
\begin{equation}
    \Delta\sigma_{y_Y,2}
    =\frac{y_{Y,2,\hat{\alpha}_p}}{y_{Y,0}}\sigma_p,
\end{equation}
and the additional hedging demand is
\begin{equation}
    \Delta h_{a,2}
    =\frac{1-\gamma^{-1}}{\psi-1}\mathbb{E}_a[\Delta\sigma_{c_j,2}]
    =\frac{1-\gamma^{-1}}{\psi-1}\mathbb{E}_a\left[\frac{c_{j,2,\hat{\alpha}_p}}{c_{j,0}}\right]\sigma_p.
\end{equation}

The new contribution to effective aggregate risk aversion comes from the second-order multiplier.
Using the expansion for $\gamma_{a,t}$ from the fixed-$\hat{\alpha}_p$ case,
\begin{equation}
    \Delta\Gamma_{a,2}=\mathbb{E}_a[\Delta\hat{\nu}_{j,2}].
\end{equation}
Therefore, the incremental market price of risk is
\begin{equation}\label{eq:delta-eta2-system-alpha-p}
    \Delta\eta_2
    =\gamma\left[
        \mathbb{E}_a[\Delta\hat{\nu}_{j,2}]\sigma
        -\Delta\sigma_{y_Y,2}
        -\Delta h_{a,2}
    \right].
\end{equation}
If no leverage constraint binds, then $\Delta\hat{\nu}_{j,2}=0$ for all active investors and \eqref{eq:delta-eta2-system-alpha-p} reduces to
\begin{equation}
    \Delta\eta_2
    =-\gamma\left[
        \frac{y_{Y,2,\hat{\alpha}_p}}{y_{Y,0}}
        +\frac{1-\gamma^{-1}}{\psi-1}\mathbb{E}_a\left[\frac{c_{j,2,\hat{\alpha}_p}}{c_{j,0}}\right]
    \right]\sigma_p.
\end{equation}

If the leverage constraint binds for investor $2$, then $\Delta\hat{\nu}_{1,2}=0$ and $\Delta\hat{\nu}_{2,2}$ is pinned down by the incremental second-order constraint.
Let
\begin{equation}
    H_j\equiv\frac{1-\gamma^{-1}}{\psi-1}\frac{c_{j,2,\hat{\alpha}_p}}{c_{j,0}}.
\end{equation}
The incremental desired exposure of investor $2$ is
\begin{equation}
    \Delta q_{2,2}=\frac{\Delta\eta_2}{\gamma}+H_2\sigma_p.
\end{equation}
The binding constraint requires the component of $\Delta\sigma_{2,2}$ in the direction of $\sigma$ to be zero.
Thus,
\begin{equation}\label{eq:delta-nu22-alpha-p}
    \Delta\hat{\nu}_{2,2}
    =\frac{\Delta q_{2,2}\sigma^\top}{\sigma\sigma^\top}
    =\frac{\Delta\eta_2\sigma^\top}{\gamma\sigma\sigma^\top}
    +H_2\frac{\sigma_p\sigma^\top}{\sigma\sigma^\top}.
\end{equation}
\begin{align}\label{eq:delta-eta2-constrained-system-alpha-p}
    \Delta\eta_2
    &=\gamma\left[
        \frac{x_2}{x_a}\Delta\hat{\nu}_{2,2}\sigma
        -\frac{y_{Y,2,\hat{\alpha}_p}}{y_{Y,0}}\sigma_p
        -\frac{1-\gamma^{-1}}{\psi-1}\mathbb{E}_a\left[\frac{c_{j,2,\hat{\alpha}_p}}{c_{j,0}}\right]\sigma_p
    \right],\\
    \Delta\hat{\nu}_{2,2}
    &=\frac{\Delta\eta_2\sigma^\top}{\gamma\sigma\sigma^\top}
    +H_2\frac{\sigma_p\sigma^\top}{\sigma\sigma^\top}.
\end{align}
Equations~\eqref{eq:delta-eta2-constrained-system-alpha-p} jointly determine the new market-price-of-risk component and the new second-order multiplier when investor $2$ is constrained.
The key difference relative to the slack case is that part of the new risk is absorbed by the multiplier, which changes effective aggregate risk aversion.

\bigskip

This system can be solved explicitly.
Let
\begin{equation}
    C_p
    \equiv
    \frac{y_{Y,2,\hat{\alpha}_p}}{y_{Y,0}}
    +\frac{1-\gamma^{-1}}{\psi-1}\mathbb{E}_a\left[\frac{c_{j,2,\hat{\alpha}_p}}{c_{j,0}}\right],
    \qquad
    \zeta_p\equiv\frac{\sigma_p\sigma^\top}{\sigma\sigma^\top}.
\end{equation}
Then \eqref{eq:delta-eta2-constrained-system-alpha-p} becomes
\begin{align}
    \Delta\eta_2
    &=\gamma\left[\frac{x_2}{x_a}\Delta\hat{\nu}_{2,2}\sigma-C_p\sigma_p\right],\\
    \Delta\hat{\nu}_{2,2}
    &=\frac{\Delta\eta_2\sigma^\top}{\gamma\sigma\sigma^\top}+H_2\zeta_p.
\end{align}
Projecting the first equation on $\sigma$ and substituting into the second equation gives
\begin{equation}
    \Delta\hat{\nu}_{2,2}
    =\frac{x_2}{x_a}\Delta\hat{\nu}_{2,2}-C_p\zeta_p+H_2\zeta_p.
\end{equation}
Since $1-x_2/x_a=x_1/x_a$, we obtain
\begin{equation}\label{eq:delta-nu22-alpha-p-solution}
    \Delta\hat{\nu}_{2,2}
    =\frac{x_a}{x_1}(H_2-C_p)\zeta_p.
\end{equation}
Substituting back into the first equation yields
\begin{equation}\label{eq:delta-eta2-constrained-alpha-p-solution}
    \Delta\eta_2
    =-\gamma C_p\sigma_p
    +\gamma\frac{x_2}{x_1}(H_2-C_p)\zeta_p\sigma.
\end{equation}
Equivalently, using
\begin{equation}
    \mathbb{E}_a[H_j]=\frac{x_1}{x_a}H_1+\frac{x_2}{x_a}H_2,
    \qquad
    C_p=\frac{y_{Y,2,\hat{\alpha}_p}}{y_{Y,0}}+\mathbb{E}_a[H_j],
\end{equation}
we can write
\begin{equation}
    H_2-C_p
    =\frac{x_1}{x_a}(H_2-H_1)-\frac{y_{Y,2,\hat{\alpha}_p}}{y_{Y,0}}.
\end{equation}

Combining the expressions above, we obtain
\begin{equation}
    \Delta\eta_2
    =-\gamma \left[ \frac{y_{Y,2,\hat{\alpha}_p}}{y_{Y,0}}+\mathbb{E}_a[H_j] \right]\sigma_p
    -\gamma\frac{x_2}{x_1}\left[\frac{y_{Y,2,\hat{\alpha}_p}}{y_{Y,0}}+\frac{x_1}{x_a}\left(H_1-H_2\right)\right]\frac{\sigma_p\sigma^\top}{\sigma\sigma^\top}\sigma.
\end{equation}
Thus the constrained investor absorbs the component of the new desired exposure in the direction of $\sigma$, while the orthogonal component of $\sigma_p$ remains priced through $\Delta\eta_2$.

\noindent\textbf{Step 4 (time-varying $\alpha_p$): Risk premium.}\\

The risk premium on the aggregate claim is given by
\begin{equation}
    \pi_{Y,t}=\sigma_{R_Y,t}\eta_t^\top.
\end{equation}
We write
\begin{equation}
    \pi_{Y,2}=\pi_{Y,2}^{0}+\Delta\pi_{Y,2},
\end{equation}
where $\pi_{Y,2}^{0}$ is the second-order risk premium in the fixed-$\hat{\alpha}_p$ economy.

Since
\begin{equation}
    \sigma_{R_Y,t}=\sigma-\sigma_{y_Y,2}\epsilon^2+\mathcal{O}(\epsilon^3),
\end{equation}
the second-order risk premium is
\begin{equation}
    \pi_{Y,2}=-\sigma_{y_Y,2}\eta_0^\top+\sigma\eta_2^\top.
\end{equation}
Therefore, the incremental risk premium induced by time variation in $\overline{\alpha}_p$ is
\begin{equation}\label{eq:delta-piY2-alpha-p}
    \Delta\pi_{Y,2}
    =-\Delta\sigma_{y_Y,2}\eta_0^\top+\sigma\Delta\eta_2^\top.
\end{equation}
Using $\eta_0=\gamma\sigma$ and Step 1,
\begin{equation}
    \Delta\pi_{Y,2}
    =-\gamma\frac{y_{Y,2,\hat{\alpha}_p}}{y_{Y,0}}\sigma_p\sigma^\top
    +\sigma\Delta\eta_2^\top.
\end{equation}

If no leverage constraint binds, substituting the expression for $\Delta\eta_2$ from Step 3 gives
\begin{equation}
    \Delta\pi_{Y,2}
    =-\gamma\left[
        \frac{y_{Y,2,\hat{\alpha}_p}}{y_{Y,0}}\sigma_p\sigma^\top
        +\left(
        \frac{y_{Y,2,\hat{\alpha}_p}}{y_{Y,0}}
        +\frac{1-\gamma^{-1}}{\psi-1}\mathbb{E}_a\left[\frac{c_{j,2,\hat{\alpha}_p}}{c_{j,0}}\right]
        \right)\sigma\sigma_p^\top
    \right].
\end{equation}
When the leverage constraint binds for investor $2$, the same formula \eqref{eq:delta-piY2-alpha-p} applies, but $\Delta\eta_2$ is jointly determined with $\Delta\hat{\nu}_{2,2}$ by \eqref{eq:delta-eta2-constrained-system-alpha-p}.

\noindent\textbf{Step 5 (time-varying $\alpha_p$): Interest rate.}\\
We write the second-order interest rate as
\begin{equation}
    r_2=r_2^{0}+\Delta r_2,
\end{equation}
where $r_2^{0}$ is the fixed-$\hat{\alpha}_p$ expression in \eqref{eq:r2-second-order}, and $\Delta r_2$ collects only the new terms generated by the dynamics of $\overline{\alpha}_p$.

From \eqref{eq:r2-second-order}, the incremental correction is
\begin{align}\label{eq:delta-r2-alpha-p}
    \Delta r_2
    &= -\psi^{-1}\left(\Delta\mu_{y_Y,2}+\sigma\Delta\sigma_{y_Y,2}^\top\right) 
\nonumber\\
    &\quad +\psi^{-1}\sum_{j=0}^2 x_j\left[
        \Delta\mu_{c_j,2}+(1-\gamma)\Delta\sigma_{c_j,2}\sigma^\top
    \right] \nonumber\\
    &\quad +(1-\psi^{-1})\gamma\sum_{j=0}^2 x_j
        \Delta\sigma_{j,2}\sigma^\top
    -\Delta\pi_{Y,2}.
\end{align}
All first-order terms in $\sigma_{j,1}$ and $\hat{\gamma}_j$ are already included in $r_2^{0}$ and therefore do not appear in $\Delta r_2$.

Using Step 1, the new drift and diffusion terms are
\begin{align}
    \Delta \sigma_{c_j,2}
    &=\frac{c_{j,2,\hat{\alpha}_p}}{c_{j,0}}\sigma_p,\\
    \Delta \sigma_{y_Y,2}
    &=\frac{y_{Y,2,\hat{\alpha}_p}}{y_{Y,0}}\sigma_p,\\
    \Delta \mu_{c_j,2}
    &=\frac{c_{j,2,\hat{\alpha}_p}}{c_{j,0}}\kappa_p(\hat{\alpha}-\hat{\alpha}_p)
    +\frac{1}{2}\frac{c_{j,2,\hat{\alpha}_p\hat{\alpha}_p}}{c_{j,0}}\sigma_p\sigma_p^\top,\\
    \Delta \mu_{y_Y,2}
    &=\frac{y_{Y,2,\hat{\alpha}_p}}{y_{Y,0}}\kappa_p(\hat{\alpha}-\hat{\alpha}_p)
    +\frac{1}{2}\frac{y_{Y,2,\hat{\alpha}_p\hat{\alpha}_p}}{y_{Y,0}}\sigma_p\sigma_p^\top.
\end{align}
Moreover, from Step 4,
\begin{equation}
    \Delta\pi_{Y,2}
    =-\Delta\sigma_{y_Y,2}\eta_0^\top+\sigma\Delta\eta_2^\top.
\end{equation}
Substituting this expression into \eqref{eq:delta-r2-alpha-p} gives
\begin{align}\label{eq:delta-r2-alpha-p-expanded}
    \Delta r_2
    &= -\psi^{-1}\Delta\mu_{y_Y,2}
    -\psi^{-1}\sigma\Delta\sigma_{y_Y,2}^\top 
\nonumber\\
    &\quad +\psi^{-1}\sum_{j=0}^2 x_j\left[
        \Delta\mu_{c_j,2}+(1-\gamma)\Delta\sigma_{c_j,2}\sigma^\top
    \right] 
\nonumber\\
    &\quad +(1-\psi^{-1})\gamma\sum_{j=0}^2 x_j
        \Delta\sigma_{j,2}\sigma^\top
    +\Delta\sigma_{y_Y,2}\eta_0^\top
    -\sigma\Delta\eta_2^\top.
\end{align}
Since $\eta_0=\gamma\sigma$, the terms involving $\Delta\sigma_{y_Y,2}$ can also be combined as
\begin{equation}
    \left(\gamma-\psi^{-1}\right)\sigma\Delta\sigma_{y_Y,2}^\top.
\end{equation}
Thus, equivalently,
\begin{align}\label{eq:delta-r2-alpha-p-compact-before-simplification}
    \Delta r_2
    &= -\psi^{-1}\Delta\mu_{y_Y,2}
    +\psi^{-1}\sum_{j=0}^2 x_j\Delta\mu_{c_j,2} 
    \nonumber\\
    &\quad +\psi^{-1}(1-\gamma)\sum_{j=0}^2 x_j\Delta\sigma_{c_j,2}\sigma^\top
    +(1-\psi^{-1})\gamma\sum_{j=0}^2 x_j\Delta\sigma_{j,2}\sigma^\top 
    \nonumber\\
    &\quad +\left(\gamma-\psi^{-1}\right)\sigma\Delta\sigma_{y_Y,2}^\top
    -\sigma\Delta\eta_2^\top.
\end{align}
This expression simplifies substantially because the baseline economy is homogeneous.
In particular, $c_{j,0}=y_{Y,0}$ for all $j$, and the aggregation identity $\Delta y_{Y,2}=\sum_{j=0}^2x_j\Delta c_{j,2}$ implies
\begin{equation}
    \sum_{j=0}^2x_j\Delta\mu_{c_j,2}=\Delta\mu_{y_Y,2},
    \qquad
    \sum_{j=0}^2x_j\Delta\sigma_{c_j,2}=\Delta\sigma_{y_Y,2}.
\end{equation}
The drift terms in \eqref{eq:delta-r2-alpha-p-compact-before-simplification} therefore cancel.
Moreover, the consumption-diffusion terms combine with the aggregate-claim diffusion term as
\begin{equation}
    \psi^{-1}(1-\gamma)\Delta\sigma_{y_Y,2}\sigma^\top
    +\left(\gamma-\psi^{-1}\right)\sigma\Delta\sigma_{y_Y,2}^\top
    =(1-\psi^{-1})\gamma\Delta\sigma_{y_Y,2}\sigma^\top.
\end{equation}
Using market clearing for the risky claim,
\begin{equation}
    \sum_{j=0}^2x_j\Delta\sigma_{j,2}= -\Delta\sigma_{y_Y,2},
\end{equation}
the remaining exposure terms cancel as well.
Hence the incremental interest-rate correction reduces to
\begin{equation}\label{eq:delta-r2-alpha-p-compact}
    \Delta r_2=-\sigma\Delta\eta_2^\top.
\end{equation}
Thus, after aggregation and market clearing, the only effect of time-varying passive demand on the second-order interest rate operates through the induced change in the market price of risk.
When no leverage constraint binds, $\Delta\eta_2$ is given directly by Step 3.
When the leverage constraint binds for investor $2$, $\Delta\eta_2$ and $\Delta\hat{\nu}_{2,2}$ are jointly determined by \eqref{eq:delta-eta2-constrained-system-alpha-p}.


\noindent\textbf{Step 6 (time-varying $\alpha_p$): Dividend yield.}\\
We write the second-order dividend yield as
\begin{equation}
    y_{Y,2}=y_{Y,2}^{0}+\Delta y_{Y,2},
\end{equation}
where $y_{Y,2}^{0}$ is the fixed-$\hat{\alpha}_p$ correction in \eqref{eq:y2-second-order}, and $\Delta y_{Y,2}$ collects only the new terms generated by the dynamics of $\overline{\alpha}_p$.

From the pricing identity,
\begin{equation}
    y_{Y,2}=r_2+\pi_{Y,2}+\mu_{y_Y,2}+\sigma\sigma_{y_Y,2}^\top,
\end{equation}
so the incremental correction satisfies
\begin{equation}\label{eq:delta-yY2-alpha-p-basic}
    \Delta y_{Y,2}
    =\Delta r_2+\Delta\pi_{Y,2}+\Delta\mu_{y_Y,2}
    +\sigma\Delta\sigma_{y_Y,2}^\top.
\end{equation}
Using \eqref{eq:delta-r2-alpha-p-compact} and \eqref{eq:delta-piY2-alpha-p}, we obtain
\begin{align}
    \Delta y_{Y,2}
    &= -\sigma\Delta\eta_2^\top
    +\left(-\Delta\sigma_{y_Y,2}\eta_0^\top+\sigma\Delta\eta_2^\top\right)
    +\Delta\mu_{y_Y,2}
    +\sigma\Delta\sigma_{y_Y,2}^\top \\
    &= \Delta\mu_{y_Y,2}
    +(1-\gamma)\Delta\sigma_{y_Y,2}\sigma^\top,
\end{align}
where the terms involving $\Delta\eta_2$ cancel and we use $\eta_0=\gamma\sigma$.

Using Step 1, this becomes
\begin{equation}\label{eq:delta-yY2-alpha-p-simplified}
    \Delta y_{Y,2}
    =\frac{y_{Y,2,\hat{\alpha}_p}}{y_{Y,0}}\kappa_p(\hat{\alpha}-\hat{\alpha}_p)
    +\frac{1}{2}\frac{y_{Y,2,\hat{\alpha}_p\hat{\alpha}_p}}{y_{Y,0}}\sigma_p\sigma_p^\top
    +(1-\gamma)\frac{y_{Y,2,\hat{\alpha}_p}}{y_{Y,0}}\sigma_p\sigma^\top.
\end{equation}
Equivalently, since $y_{Y,2}=y_{Y,2}^{0}+\Delta y_{Y,2}$, the incremental dividend-yield correction solves
\begin{equation}\label{eq:delta-yY2-alpha-p-operator}
    \Delta y_{Y,2}
    =\frac{(y_{Y,2}^{0}+\Delta y_{Y,2})_{\hat{\alpha}_p}}{y_{Y,0}}
    \left[\kappa_p(\hat{\alpha}-\hat{\alpha}_p)+(1-\gamma)\sigma_p\sigma^\top\right]
    +\frac{1}{2}\frac{(y_{Y,2}^{0}+\Delta y_{Y,2})_{\hat{\alpha}_p\hat{\alpha}_p}}{y_{Y,0}}\sigma_p\sigma_p^\top.
\end{equation}
This is the time-varying-$\alpha_p$ counterpart of \eqref{eq:y2-second-order}: the new correction is the drift and risk-adjusted diffusion of the total second-order dividend yield induced by the state variable $\overline{\alpha}_p$.

The second-order price-dividend-ratio correction follows from
\begin{equation}
    p_{Y,2}=p_{Y,0}^3y_{Y,1}^2-p_{Y,0}^2y_{Y,2},
\end{equation}
so the incremental term is
\begin{equation}
    \Delta p_{Y,2}=-p_{Y,0}^2\Delta y_{Y,2}.
\end{equation}


\noindent\textbf{Step 7 (time-varying $\alpha_p$): Solving for coefficients of dividend yield.}\\
We can solve the coefficients of $\Delta y_{Y,2}$ directly from \eqref{eq:delta-yY2-alpha-p-operator}.
Write the fixed-$\hat{\alpha}_p$ correction as
\begin{equation}
    y_{Y,2}^{0}(x,\hat{\alpha}_p)
    =a_Y^{0}(x)+b_Y^{0}(x)\hat{\alpha}_p+d_Y^{0}(x)\hat{\alpha}_p^2,
\end{equation}
and conjecture
\begin{equation}
    \Delta y_{Y,2}(x,\hat{\alpha}_p)
    =a_Y(x)+b_Y(x)\hat{\alpha}_p+d_Y(x)\hat{\alpha}_p^2.
\end{equation}
Let
\begin{equation}
    b_Y^T(x)\equiv b_Y^{0}(x)+b_Y(x),
    \qquad
    d_Y^T(x)\equiv d_Y^{0}(x)+d_Y(x)
\end{equation}
denote the linear and quadratic coefficients of the total second-order dividend-yield correction $y_{Y,2}^{0}+\Delta y_{Y,2}$.
Then
\begin{equation}
    (y_{Y,2}^{0}+\Delta y_{Y,2})_{\hat{\alpha}_p}
    =b_Y^T+2d_Y^T\hat{\alpha}_p,
    \qquad
    (y_{Y,2}^{0}+\Delta y_{Y,2})_{\hat{\alpha}_p\hat{\alpha}_p}
    =2d_Y^T.
\end{equation}
For compactness, define
\begin{equation}
    M_p\equiv \kappa_p\hat{\alpha}+(1-\gamma)\sigma_p\sigma^\top,
    \qquad
    V_p\equiv \sigma_p\sigma_p^\top.
\end{equation}
Then \eqref{eq:delta-yY2-alpha-p-operator} can be written as
\begin{equation}\label{eq:delta-yY2-alpha-p-coeff-operator}
    \Delta y_{Y,2}
    =\frac{1}{y_{Y,0}}(b_Y^T+2d_Y^T\hat{\alpha}_p)
    \left(M_p-\kappa_p\hat{\alpha}_p\right)
    +\frac{d_Y^T}{y_{Y,0}}V_p.
\end{equation}

We first match the coefficient on $\hat{\alpha}_p^2$.
The only term in \eqref{eq:delta-yY2-alpha-p-coeff-operator} that generates a quadratic coefficient is the product of $2d_Y^T\hat{\alpha}_p$ with $-\kappa_p\hat{\alpha}_p$.
Therefore,
\begin{equation}
    d_Y
    =-\frac{2\kappa_p}{y_{Y,0}}d_Y^T
    =-\frac{2\kappa_p}{y_{Y,0}}(d_Y^{0}+d_Y).
\end{equation}
Solving for $d_Y$ yields
\begin{equation}\label{eq:dY-solution-alpha-p-simple}
    d_Y
    =-\frac{\frac{2\kappa_p}{y_{Y,0}}}{1+\frac{2\kappa_p}{y_{Y,0}}}\,d_Y^{0}.
\end{equation}
Equivalently, the quadratic coefficient of the total second-order correction is
\begin{equation}\label{eq:dY-total-solution-alpha-p-simple}
    d_Y^T
    =\frac{1}{1+\frac{2\kappa_p}{y_{Y,0}}}\,d_Y^{0}.
\end{equation}
Thus, mean reversion in passive demand attenuates the quadratic dependence of the dividend yield on $\hat{\alpha}_p$.
When $\kappa_p=0$, we recover $d_Y=0$ and $d_Y^T=d_Y^{0}$, as in the fixed-$\hat{\alpha}_p$ economy.

We can similarly solve for the coefficient on $\hat{\alpha}_p$.
Matching the linear coefficient in \eqref{eq:delta-yY2-alpha-p-coeff-operator} gives
\begin{equation}
    b_Y
    =\frac{1}{y_{Y,0}}\left(2d_Y^T M_p-\kappa_p b_Y^T\right)
    =\frac{1}{y_{Y,0}}\left(2d_Y^T M_p-\kappa_p(b_Y^{0}+b_Y)\right).
\end{equation}
Solving for $b_Y$ yields
\begin{equation}\label{eq:bY-solution-alpha-p-simple}
    b_Y
    =\frac{2d_Y^T M_p-\kappa_p b_Y^{0}}{y_{Y,0}+\kappa_p}.
\end{equation}
Substituting \eqref{eq:dY-total-solution-alpha-p-simple}, this can also be written as
\begin{equation}\label{eq:bY-solution-alpha-p-substituted}
    b_Y
    =\frac{\frac{2M_p}{1+\frac{2\kappa_p}{y_{Y,0}}}d_Y^{0}-\kappa_p b_Y^{0}}{y_{Y,0}+\kappa_p}.
\end{equation}
Equivalently, the linear coefficient of the total second-order correction is
\begin{equation}\label{eq:bY-total-solution-alpha-p-simple}
    b_Y^T
    =b_Y^{0}+b_Y
    =\frac{y_{Y,0}b_Y^{0}+2d_Y^T M_p}{y_{Y,0}+\kappa_p}.
\end{equation}

Finally, we solve for the constant coefficient.
Matching the constant term in \eqref{eq:delta-yY2-alpha-p-coeff-operator} gives
\begin{equation}
    a_Y
    =\frac{1}{y_{Y,0}}\left(b_Y^T M_p+d_Y^T V_p\right).
\end{equation}
Using \eqref{eq:bY-total-solution-alpha-p-simple} and \eqref{eq:dY-total-solution-alpha-p-simple}, this becomes
\begin{equation}\label{eq:aY-solution-alpha-p-simple}
    a_Y
    =\frac{1}{y_{Y,0}}\left[
        \frac{y_{Y,0}b_Y^{0}+2d_Y^T M_p}{y_{Y,0}+\kappa_p}M_p
        +d_Y^T V_p
    \right],
\end{equation}
where
\begin{equation}
    d_Y^T=\frac{1}{1+\frac{2\kappa_p}{y_{Y,0}}}d_Y^{0}.
\end{equation}
Equivalently, substituting out $d_Y^T$,
\begin{equation}\label{eq:aY-solution-alpha-p-substituted}
    a_Y
    =\frac{1}{y_{Y,0}}\left[
        \frac{y_{Y,0}b_Y^{0}+\frac{2M_p}{1+\frac{2\kappa_p}{y_{Y,0}}}d_Y^{0}}{y_{Y,0}+\kappa_p}M_p
        +\frac{1}{1+\frac{2\kappa_p}{y_{Y,0}}}d_Y^{0}V_p
    \right].
\end{equation}
Therefore, the incremental dividend-yield correction is
\begin{equation}\label{eq:delta-yY2-coeff-solution-alpha-p}
    \Delta y_{Y,2}(x,\hat{\alpha}_p)
    =a_Y(x)+b_Y(x)\hat{\alpha}_p+d_Y(x)\hat{\alpha}_p^2,
\end{equation}
with $d_Y$, $b_Y$, and $a_Y$ given by \eqref{eq:dY-solution-alpha-p-simple}, \eqref{eq:bY-solution-alpha-p-simple}, and \eqref{eq:aY-solution-alpha-p-simple}.

\noindent\textbf{Step 8 (time-varying $\alpha_p$): Consumption-wealth ratio.}\\
We write the second-order consumption-wealth ratio as
\begin{equation}
    c_{j,2}=c_{j,2}^{0}+\Delta c_{j,2},
\end{equation}
where $c_{j,2}^{0}$ is the fixed-$\hat{\alpha}_p$ correction from the previous subsection and $\Delta c_{j,2}$ collects only the new terms generated by the dynamics of $\overline{\alpha}_p$.

The consumption-wealth ratio satisfies
\begin{equation}
    c_{j,t} = \psi \rho 
    + (1-\psi) \left[ r_t + \eta_{t}\,\sigma_{j,t}^\top - \frac{\gamma_j}{2} \|\sigma_{j,t}\|^2 \right] 
    + \mu_{c_j,t} 
    + (1-\gamma_j) \sigma_{c_j,t} \sigma_{j,t}^\top 
    + \frac{\psi - \gamma_j}{1-\psi} \frac{\|\sigma_{c_j,t}\|^2}{2}.
\end{equation}
The last term is fourth order and does not enter the second-order expansion.

From the fixed-$\hat{\alpha}_p$ derivation, the second-order term can be written as
\begin{align}
    c_{j,2} &= (1-\psi)\Big[
        r_2
        + \eta_2\sigma^\top
        + \eta_1\sigma_{j,1}^\top
        - \frac{\gamma}{2}\sigma_{j,1}\sigma_{j,1}^\top
        - \gamma\hat{\gamma}_j\sigma\sigma_{j,1}^\top
    \Big] \nonumber\\
    &\quad + \mu_{c_j,2}
    + (1-\gamma)\sigma_{c_j,2}\sigma^\top.
\end{align}
Here the terms involving $\sigma_{j,2}$ cancel because $\eta_0=\gamma\sigma$.
Therefore, relative to the fixed-$\hat{\alpha}_p$ economy,
\begin{align}\label{eq:delta-cj2-alpha-p}
    \Delta c_{j,2}
    &= (1-\psi)\left[\Delta r_2+\Delta\eta_2\sigma^\top\right]
    +\Delta\mu_{c_j,2}
    +(1-\gamma)\Delta\sigma_{c_j,2}\sigma^\top.
\end{align}
All first-order terms involving $\eta_1$, $\sigma_{j,1}$, and $\hat{\gamma}_j$ are already included in $c_{j,2}^{0}$ and therefore do not appear in $\Delta c_{j,2}$.

Using Step 1, this becomes
\begin{align}\label{eq:delta-cj2-alpha-p-expanded}
    \Delta c_{j,2}
    &= (1-\psi)\left[\Delta r_2+\Delta\eta_2\sigma^\top\right] \nonumber\\
    &\quad +\frac{c_{j,2,\hat{\alpha}_p}}{c_{j,0}}\kappa_p(\hat{\alpha}-\hat{\alpha}_p)
    +\frac{1}{2}\frac{c_{j,2,\hat{\alpha}_p\hat{\alpha}_p}}{c_{j,0}}\sigma_p\sigma_p^\top 
\nonumber\\
    &\quad +(1-\gamma)\frac{c_{j,2,\hat{\alpha}_p}}{c_{j,0}}\sigma_p\sigma^\top.
\end{align}
This equation is the time-varying-$\alpha_p$ counterpart of the fixed-$\hat{\alpha}_p$ expression for $c_{j,2}$.
It shows that the new correction to the consumption-wealth ratio is driven by the changes in the interest rate and market price of risk, together with the direct drift and diffusion terms induced by the dynamics of passive demand.

Equations~\eqref{eq:delta-r2-alpha-p-compact}, \eqref{eq:delta-eta2-system-alpha-p}, and \eqref{eq:delta-cj2-alpha-p-expanded} form the closed system for the incremental second-order corrections when the passive portfolio share is time-varying.

% Step 9: Solve for coefficients of consumption-wealth ratios
\noindent\textbf{Step 9 (time-varying $\alpha_p$): Solving for coefficients of consumption-wealth ratios.}\\
We now solve for the coefficients of the incremental correction to the consumption-wealth ratios.
The simplification in Step 5 implies
\begin{equation}
    \Delta r_2+\Delta\eta_2\sigma^\top=0.
\end{equation}
Using \eqref{eq:delta-cj2-alpha-p-expanded}, the incremental correction to $c_{j,2}$ satisfies
\begin{equation}\label{eq:delta-cj2-alpha-p-operator-simple}
    \Delta c_{j,2}
    =\frac{(c_{j,2}^{0}+\Delta c_{j,2})_{\hat{\alpha}_p}}{c_{j,0}}
    \left[\kappa_p(\hat{\alpha}-\hat{\alpha}_p)+(1-\gamma)\sigma_p\sigma^\top\right]
    +\frac{1}{2}\frac{(c_{j,2}^{0}+\Delta c_{j,2})_{\hat{\alpha}_p\hat{\alpha}_p}}{c_{j,0}}\sigma_p\sigma_p^\top.
\end{equation}
This is the same scalar coefficient problem as for the aggregate dividend yield, but applied investor by investor.

Write the fixed-$\hat{\alpha}_p$ correction as
\begin{equation}
    c_{j,2}^{0}(x,\hat{\alpha}_p)
    =a_j^{0}(x)+b_j^{0}(x)\hat{\alpha}_p+d_j^{0}(x)\hat{\alpha}_p^2,
\end{equation}
and conjecture
\begin{equation}
    \Delta c_{j,2}(x,\hat{\alpha}_p)
    =a_j(x)+b_j(x)\hat{\alpha}_p+d_j(x)\hat{\alpha}_p^2.
\end{equation}
Let
\begin{equation}
    b_j^T(x)\equiv b_j^{0}(x)+b_j(x),
    \qquad
    d_j^T(x)\equiv d_j^{0}(x)+d_j(x)
\end{equation}
denote the linear and quadratic coefficients of the total second-order correction $c_{j,2}^{0}+\Delta c_{j,2}$.
As before, define
\begin{equation}
    M_p\equiv\kappa_p\hat{\alpha}+(1-\gamma)\sigma_p\sigma^\top,
    \qquad
    V_p\equiv\sigma_p\sigma_p^\top.
\end{equation}
Then \eqref{eq:delta-cj2-alpha-p-operator-simple} can be written as
\begin{equation}\label{eq:delta-cj2-alpha-p-coeff-operator}
    \Delta c_{j,2}
    =\frac{1}{c_{j,0}}(b_j^T+2d_j^T\hat{\alpha}_p)
    (M_p-\kappa_p\hat{\alpha}_p)
    +\frac{d_j^T}{c_{j,0}}V_p.
\end{equation}

Matching the coefficient on $\hat{\alpha}_p^2$ gives
\begin{equation}
    d_j=-\frac{2\kappa_p}{c_{j,0}}d_j^T
    =-\frac{2\kappa_p}{c_{j,0}}(d_j^0+d_j).
\end{equation}
Hence
\begin{equation}\label{eq:dj-solution-alpha-p-simple}
    d_j
    =-\frac{\frac{2\kappa_p}{c_{j,0}}}{1+\frac{2\kappa_p}{c_{j,0}}}\,d_j^0,
\end{equation}
and the quadratic coefficient of the total correction is
\begin{equation}\label{eq:dj-total-solution-alpha-p-simple}
    d_j^T
    =\frac{1}{1+\frac{2\kappa_p}{c_{j,0}}}\,d_j^0.
\end{equation}

Matching the coefficient on $\hat{\alpha}_p$ gives
\begin{equation}
    b_j
    =\frac{1}{c_{j,0}}\left(2d_j^T M_p-\kappa_p b_j^T\right)
    =\frac{1}{c_{j,0}}\left(2d_j^T M_p-\kappa_p(b_j^0+b_j)\right).
\end{equation}
Solving for $b_j$ yields
\begin{equation}\label{eq:bj-solution-alpha-p-simple}
    b_j
    =\frac{2d_j^T M_p-\kappa_p b_j^0}{c_{j,0}+\kappa_p}.
\end{equation}
Equivalently, the linear coefficient of the total correction is
\begin{equation}\label{eq:bj-total-solution-alpha-p-simple}
    b_j^T
    =b_j^0+b_j
    =\frac{c_{j,0}b_j^0+2d_j^T M_p}{c_{j,0}+\kappa_p}.
\end{equation}

Finally, matching the constant term gives
\begin{equation}
    a_j
    =\frac{1}{c_{j,0}}\left(b_j^T M_p+d_j^T V_p\right).
\end{equation}
Using \eqref{eq:bj-total-solution-alpha-p-simple} and \eqref{eq:dj-total-solution-alpha-p-simple}, this becomes
\begin{equation}\label{eq:aj-solution-alpha-p-simple}
    a_j
    =\frac{1}{c_{j,0}}\left[
        \frac{c_{j,0}b_j^0+2d_j^T M_p}{c_{j,0}+\kappa_p}M_p
        +d_j^T V_p
    \right].
\end{equation}
Therefore,
\begin{equation}\label{eq:delta-cj2-coeff-solution-alpha-p}
    \Delta c_{j,2}(x,\hat{\alpha}_p)
    =a_j(x)+b_j(x)\hat{\alpha}_p+d_j(x)\hat{\alpha}_p^2,
\end{equation}
with $d_j$, $b_j$, and $a_j$ given by \eqref{eq:dj-solution-alpha-p-simple}, \eqref{eq:bj-solution-alpha-p-simple}, and \eqref{eq:aj-solution-alpha-p-simple}.
Because $c_{j,0}=y_{Y,0}$ in the homogeneous benchmark economy and $y_{Y,2}=\sum_{j=0}^2x_jc_{j,2}$, aggregating \eqref{eq:delta-cj2-coeff-solution-alpha-p} across investors recovers the coefficient formulas for $\Delta y_{Y,2}$ derived above.

\noindent\textbf{Step 10 (time-varying $\alpha_p$): Pricing the dividend claim.}\\
The dividend claim satisfies the pricing condition
\begin{align}\label{eq:pD-pricing-timevarying-alpha}
    \left[r_t+\sigma\eta_t^\top-\mu\right]p_{D,t}
    &= s_t + p_{D,x,t}\mu_{x,t}+p_{D,s,t}\mu_{s,t}
    +p_{D,\overline{\alpha}_p,t}\mu_{\overline{\alpha}_p,t}
    +\frac{1}{2}\sum_{k=1}^3 \Sigma_{D,k,t}^\top p_{D,XX,t}\Sigma_{D,k,t}
    \\
    &+(\sigma-\eta_t)\left(p_{D,x,t}\sigma_{x,t}+p_{D,s,t}\sigma_{s,t}
    +p_{D,\overline{\alpha}_p,t}\sigma_p\epsilon\right)^\top,
\end{align}
where now $X=(x,s,\overline{\alpha}_p)$.
Relative to the fixed-$\hat{\alpha}_p$ case, the new terms are those involving derivatives with respect to $\overline{\alpha}_p$ and the covariance between $s$ and $\overline{\alpha}_p$.

We write
\begin{equation}
    p_{D,2}^{T}(x,\hat{s},\hat{\alpha}_p)
    \equiv p_{D,2}^{0}(x,\hat{s},\hat{\alpha}_p)+\Delta p_{D,2}(x,\hat{s},\hat{\alpha}_p),
\end{equation}
where $p_{D,2}^{0}$ is the fixed-$\hat{\alpha}_p$ correction derived above.
From the fixed-$\hat{\alpha}_p$ solution,
\begin{equation}
    p_{D,2}^{0}(x,\hat{s},\hat{\alpha}_p)
    =a_2^{0}(x,\hat{\alpha}_p)+b_2^{0}(x)\hat{s},
\end{equation}
where $b_2^{0}$ is independent of $\hat{\alpha}_p$.
We conjecture that the time-varying-$\alpha_p$ correction has the same form,
\begin{equation}
    p_{D,2}^{T}(x,\hat{s},\hat{\alpha}_p)
    =a_2^{T}(x,\hat{\alpha}_p)+b_2^{T}(x,\hat{\alpha}_p)\hat{s}.
\end{equation}

The simplification in Step 5 implies
\begin{equation}
    \Delta A_2
    \equiv \Delta r_2+\sigma\Delta\eta_2^\top=0.
\end{equation}
Thus the aggregate-price term in the dividend-claim pricing equation is unchanged relative to the fixed-$\hat{\alpha}_p$ economy.
The only new terms come from the generator associated with the additional state variable $\overline{\alpha}_p$.
Using the chain rule for $\overline{\alpha}_p=1+\epsilon\hat{\alpha}_p$, the additional second-order operator applied to a function $f(\hat{s},\hat{\alpha}_p)$ is
\begin{equation}\label{eq:L-alpha-pD}
    \mathcal{L}_{\alpha}f
    =f_{\hat{\alpha}_p}\left[\kappa_p(\hat{\alpha}-\hat{\alpha}_p)+(\sigma-\eta_0)\sigma_p^\top\right]
    +\frac{1}{2}f_{\hat{\alpha}_p\hat{\alpha}_p}\sigma_p\sigma_p^\top
    +B_s f_{\hat{s}\hat{\alpha}_p}\sigma_s\sigma_p^\top,
\end{equation}
where $B_s\equiv\bar{s}(1-\bar{s})$.
Using $\eta_0=\gamma\sigma$, define
\begin{equation}
    M_p\equiv\kappa_p\hat{\alpha}+(1-\gamma)\sigma_p\sigma^\top,
    \qquad
    V_p\equiv\sigma_p\sigma_p^\top.
\end{equation}
Then
\begin{equation}
    \mathcal{L}_{\alpha}f
    =f_{\hat{\alpha}_p}(M_p-\kappa_p\hat{\alpha}_p)
    +\frac{1}{2}f_{\hat{\alpha}_p\hat{\alpha}_p}V_p
    +B_s f_{\hat{s}\hat{\alpha}_p}\sigma_s\sigma_p^\top.
\end{equation}

Subtracting the fixed-$\hat{\alpha}_p$ equation from the time-varying-$\alpha_p$ equation gives
\begin{equation}\label{eq:delta-pD2-alpha-operator}
    \delta\Delta p_{D,2}
    =-\kappa_s\hat{s}\Delta p_{D,2,\hat{s}}
    +\Lambda_s B_s\Delta p_{D,2,\hat{s}}
    +\frac{1}{2}B_s^2\sigma_s\sigma_s^\top\Delta p_{D,2,\hat{s}\hat{s}}
    +\mathcal{L}_{\alpha}p_{D,2}^{T},
\end{equation}
where $\Lambda_s\equiv(\sigma-\eta_0)\sigma_s^\top$.
Because $p_{D,2}^{0}$ is affine in $\hat{s}$ and $b_2^{0}$ is independent of $\hat{\alpha}_p$, the solution remains affine in $\hat{s}$.
Moreover, matching the coefficient on $\hat{s}$ gives
\begin{equation}
    (\delta+\kappa_s)\Delta b_2=\mathcal{L}_{\alpha}b_2^{T} = \mathcal{L}_{\alpha}\Delta b_2,
\end{equation}
as $b_2^0 $ is independent of $\hat{\alpha}_p$.
Then, $\Delta b_2 = 0$ satisfies the equation above, so
\begin{equation}\label{eq:pD-b2-timevarying-alpha}
    b_2^{T}=b_2^{0}.
\end{equation}
Thus, time variation in passive demand affects only the constant term in the dividend-claim correction.
The cross term $B_s f_{\hat{s}\hat{\alpha}_p}\sigma_s\sigma_p^\top$ in \eqref{eq:L-alpha-pD} drops out because $b_2^{T}$ is independent of $\hat{\alpha}_p$.

The constant term satisfies the scalar equation
\begin{equation}\label{eq:a2-timevarying-alpha-operator}
    \Delta a_2
    =\frac{1}{\delta}\left[
        a_{2,\hat{\alpha}_p}^{T}(M_p-\kappa_p\hat{\alpha}_p)
        +\frac{1}{2}a_{2,\hat{\alpha}_p\hat{\alpha}_p}^{T}V_p
    \right].
\end{equation}
This equation is identical in form to the coefficient equation for the aggregate dividend yield.
Indeed, in the benchmark economy $y_{Y,0}=\delta$, where $\delta\equiv r_0+\sigma\eta_0^\top-\mu$.
Write
\begin{equation}
    a_2^{0}(x,\hat{\alpha}_p)=a_D^{0}(x)+b_D^{0}(x)\hat{\alpha}_p+d_D^{0}(x)\hat{\alpha}_p^2,
\end{equation}
and
\begin{equation}
    \Delta a_2(x,\hat{\alpha}_p)=a_D(x)+b_D(x)\hat{\alpha}_p+d_D(x)\hat{\alpha}_p^2.
\end{equation}
Let $b_D^T\equiv b_D^0+b_D$ and $d_D^T\equiv d_D^0+d_D$ denote the linear and quadratic coefficients of $a_2^T=a_2^0+\Delta a_2$.
Matching coefficients in \eqref{eq:a2-timevarying-alpha-operator} gives
\begin{align}
    d_D^T
    &=\frac{1}{1+\frac{2\kappa_p}{\delta}}d_D^{0},\\
    b_D^T
    &=\frac{\delta b_D^{0}+2d_D^T M_p}{\delta+\kappa_p},\\
    a_D
    &=\frac{1}{\delta}\left(b_D^T M_p+d_D^T V_p\right).
\end{align}
Therefore,
\begin{equation}\label{eq:pD2-timevarying-alpha-solution}
    p_{D,2}^{T}(x,\hat{s},\hat{\alpha}_p)
    =a_D^{0}(x)+a_D(x)+b_D^T(x)\hat{\alpha}_p+d_D^T(x)\hat{\alpha}_p^2
    +b_2^{0}(x)\hat{s}.
\end{equation}
Finally, since $p_{D,1}$ is independent of $\hat{\alpha}_p$, the price impact of passive flows for the dividend claim is
\begin{equation}\label{eq:pD-price-impact-timevarying-alpha}
    \frac{\partial p_D(x,s,\overline{\alpha}_p;\epsilon)}{\partial F(x,s,\overline{\alpha}_p;\epsilon)}
    =\frac{\epsilon}{x_0}\left(b_D^T+2d_D^T\hat{\alpha}_p\right)
    +\mathcal{O}(\epsilon^2).
\end{equation}
We now relate this expression explicitly to the price impact of the aggregate claim.
Let
\begin{equation}
    G_Y^T
    \equiv \frac{\partial y_{Y,2}^T}{\partial\hat{\alpha}_p}
    = b_Y^T+2d_Y^T\hat{\alpha}_p,
\end{equation}
where $y_{Y,2}^T\equiv y_{Y,2}^0+\Delta y_{Y,2}$.
Because $p_Y=1/y_Y$ and $y_{Y,1}$ is independent of $\hat{\alpha}_p$,
\begin{equation}\label{eq:pY-price-impact-timevarying-alpha}
    \frac{\partial p_Y(x,\overline{\alpha}_p;\epsilon)}{\partial F(x,\overline{\alpha}_p;\epsilon)}
    =-\frac{\epsilon}{x_0}\frac{1}{\delta^2}G_Y^T
    +\mathcal{O}(\epsilon^2).
\end{equation}

The fixed-$\hat{\alpha}_p$ calculation gives the derivative of the dividend-claim correction as
\begin{equation}\label{eq:pD-fixed-alpha-decomposition-reuse}
    \frac{\partial p_{D,2}^{0}}{\partial\hat{\alpha}_p}
    =-\frac{\bar{s}}{\delta^2}\frac{\partial y_{Y,2}^{0}}{\partial\hat{\alpha}_p}
    -\frac{B_s}{\delta(\delta+\kappa_s)}
        \frac{\partial\eta_1}{\partial\hat{\alpha}_p}\sigma_s^\top.
\end{equation}
The first term is the aggregate-claim channel, while the second term is the dividend-risk-premium channel.
Equivalently, define the fixed-$\hat{\alpha}_p$ excess slope
\begin{equation}\label{eq:GD-fixed-excess-slope-alpha}
    G_D^{\mathrm{ex},0}
    \equiv
    \frac{\partial p_{D,2}^{0}}{\partial\hat{\alpha}_p}
    +\frac{\bar{s}}{\delta^2}\frac{\partial y_{Y,2}^{0}}{\partial\hat{\alpha}_p}
    =-\frac{B_s}{\delta(\delta+\kappa_s)}
        \frac{\partial\eta_1}{\partial\hat{\alpha}_p}\sigma_s^\top.
\end{equation}
This excess slope is independent of $\hat{\alpha}_p$.
Therefore, when passive portfolios are time-varying, the same coefficient equation attenuates this slope by the factor $\delta/(\delta+\kappa_p)$, but it does not create an additional $M_p$ term in the dividend-specific component:
\begin{equation}\label{eq:GD-timevarying-excess-slope-alpha}
    G_D^{\mathrm{ex},T}
    =\frac{\delta}{\delta+\kappa_p}G_D^{\mathrm{ex},0}
    =-\frac{\delta}{\delta+\kappa_p}
    \frac{B_s}{\delta(\delta+\kappa_s)}
        \frac{\partial\eta_1}{\partial\hat{\alpha}_p}\sigma_s^\top.
\end{equation}
Combining \eqref{eq:pY-price-impact-timevarying-alpha} and \eqref{eq:GD-timevarying-excess-slope-alpha}, the dividend-claim price impact is
\begin{equation}\label{eq:pD-price-impact-explicit-final-alpha}
    \frac{\partial p_D}{\partial F}
    =\bar{s}\frac{\partial p_Y}{\partial F}
    -\frac{\epsilon}{x_0}
    \frac{\delta}{\delta+\kappa_p}
    \frac{B_s}{\delta(\delta+\kappa_s)}
        \frac{\partial\eta_1}{\partial\hat{\alpha}_p}\sigma_s^\top
    +\mathcal{O}(\epsilon^2).
\end{equation}
Equation~\eqref{eq:pD-price-impact-explicit-final-alpha} is the time-varying-passive-share analogue of the fixed-$\hat{\alpha}_p$ decomposition.
There are two effects of portfolio dynamics.
First, $\partial p_Y/\partial F$ is itself modified by mean reversion and passive-demand risk through the coefficients $b_Y^T$ and $d_Y^T$.
Second, the dividend-specific risk-premium channel is attenuated by $\delta/(\delta+\kappa_p)$.
There is no separate third dividend-specific term in this benchmark specification because the fixed-$\hat{\alpha}_p$ excess slope in \eqref{eq:GD-fixed-excess-slope-alpha} is constant in $\hat{\alpha}_p$.
A separate $M_p$ term would appear only if the fixed-$\hat{\alpha}_p$ excess component had a nonzero quadratic coefficient in $\hat{\alpha}_p$.


\section{Special Case: No Preference Heterogeneity}

In this section, we consider the special case of no preference heterogeneity or leverage constraints, that is, $\hat{\gamma}_j=0$ for $j=0,1,2$ and $\hat{\nu}_2=0$.
This will allow us to derive the cleaner expressions for all equilibrium objects.
For simplicity, we focus on the case where $\hat{\alpha}=0$.

\begin{prop}[Time-varying passive demand without heterogeneity or constraints]
    Consider the case with no preference heterogeneity and no leverage constraints, that is, $\hat{\gamma}_j=0$ for $j=0,1,2$ and $\hat{\nu}_2=0$.
    Let
    \begin{equation}
        M_p\equiv (1-\gamma)\sigma_p\sigma^\top,
        \qquad
        d_Y^T
        \equiv
        \frac{1}{1+2\kappa_p/y_{Y,0}}
        \frac{1}{2}(1-\psi^{-1})\gamma\|\sigma\|^2\frac{x_0}{x_a}.
    \end{equation}
    Then the second-order correction to the dividend yield is
    \begin{equation}
        y_{Y,2}^T
        =
        \frac{2d_Y^T}{y_{Y,0}}
        \left[
            \frac{M_p^2}{y_{Y,0}+\kappa_p}
            +\frac{\sigma_p\sigma_p^\top}{2}
            +\frac{y_{Y,0}}{y_{Y,0}+\kappa_p}M_p\hat{\alpha}_p
            +\frac{y_{Y,0}}{2}\hat{\alpha}_p^2
        \right].
    \end{equation}
    The inverse aggregate market elasticity is
    \begin{equation}
        \epsilon_M^{-1}
        =
        \left(1-\psi^{-1}\right)
        \frac{\gamma\|\sigma\|^2}{y_{Y,0}+2\kappa_p}
        \frac{1}{x_a}
        \left[
            1-\overline{\alpha}_p
            +\frac{(\gamma-1)\sigma_p\epsilon\sigma^\top}{y_{Y,0}+\kappa_p}
        \right]
        +\mathcal{O}(\epsilon^2).
    \end{equation}
    The return volatility satisfies
    \begin{equation}
        \sigma_{R_Y}
        =
        \sigma
        +x_0\epsilon_M^{-1}\sigma_p\epsilon
        +\mathcal{O}(\epsilon^3).
    \end{equation}
    The market price of risk is
    \begin{equation}
        \eta
        = \gamma \sigma \left( \frac{1- x_0 \overline{\alpha}_p}{x_a}\right)
        +\left[
            1 -\frac{1-\gamma^{-1}}{1-\psi^{-1}} \left(\frac{1}{x_a} - \psi^{-1}\right)
        \right]\gamma x_0\epsilon_M^{-1}\sigma_p\epsilon
        +\mathcal{O}(\epsilon^3).
    \end{equation}
    The risk premium is
    \begin{equation}
        \pi_Y
        =
        \gamma\|\sigma\|^2
        \left(\frac{1-x_0\overline{\alpha}_p}{x_a}\right)
        +(2-H_a)\gamma x_0\epsilon_M^{-1}\sigma\sigma_p^\top\epsilon
        +\mathcal{O}(\epsilon^3),
    \end{equation}
    where
    \begin{equation}
        H_a
        \equiv
        \frac{1-\gamma^{-1}}{1-\psi^{-1}} \left(\frac{1}{x_a} - \psi^{-1}\right).
    \end{equation}
    The incremental interest-rate correction is
    \begin{equation}
        \Delta r_2
        =(1-H_a)\gamma
        \frac{G_p}{y_{Y,0}}\sigma\sigma_p^\top,
    \end{equation}
    where
    \begin{equation}
        G_p
        \equiv
        \frac{\partial y_{Y,2}^T}{\partial \hat{\alpha}_p}
        =2d_Y^T\left(\hat{\alpha}_p+\frac{M_p}{y_{Y,0}+\kappa_p}\right).
    \end{equation}
    Finally, the normalized price impact of the dividend claim satisfies
    \begin{equation}
        \frac{1}{p_D}\frac{\partial p_D}{\partial F}
        =
        \frac{1}{p_Y}\frac{\partial p_Y}{\partial F}
        + \frac{\delta\gamma(1-\bar{s})}{(\delta+\kappa_p)(\delta+\kappa_s)}
        \frac{\epsilon}{x_a}\sigma_s\sigma^\top
        +\mathcal{O}(\epsilon^2).
    \end{equation}
\end{prop}

\begin{proof}
    To derive the results, we rely on the general coefficient formulas derived in Appendix~\ref{app:derivations-perturbation-solution}.
    \paragraph{Dividend yield.}
In the absence of preference heterogeneity and leverage constraints,
\[
a_Y^0=b_Y^0=0,
\qquad
d_Y^0
=
\frac12(1-\psi^{-1})\gamma\|\sigma\|^2\frac{x_0}{x_a}.
\]

Applying the general coefficient formulas from
\eqref{eq:dY-total-solution-alpha-p-simple}--
\eqref{eq:aY-solution-alpha-p-simple},
the total second-order dividend-yield correction is
\[
y_{Y,2}^T
=
a_Y^T+b_Y^T\hat\alpha_p+d_Y^T\hat\alpha_p^2,
\]
where
\[
d_Y^T
=
\frac{1}{1+2\kappa_p/y_{Y,0}}d_Y^0,
\]
\[
b_Y^T
=
\frac{2d_Y^TM_p}{y_{Y,0}+\kappa_p},
\]
\[
a_Y^T
=
\frac{1}{y_{Y,0}}
\left(
b_Y^TM_p+d_Y^TV_p
\right).
\]
Thus, substituting the coefficients above,
\begin{equation}\label{eq:yY2-total-noheterog-closed-form}
    y_{Y,2}^T
    =
    \frac{2d_Y^T}{y_{Y,0}}
    \left[
        \frac{M_p^2}{y_{Y,0}+\kappa_p}
        +\frac{V_p}{2}
        +\frac{y_{Y,0}}{y_{Y,0}+\kappa_p}M_p\hat{\alpha}_p
        +\frac{y_{Y,0}}{2}\hat{\alpha}_p^2
    \right].
\end{equation}
Using the expression for $d_Y^T$, this becomes
\begin{equation}\label{eq:yY2-total-noheterog-primitives}
    y_{Y,2}^T
    =
    \left(1-\psi^{-1}\right)
    \frac{\gamma\|\sigma\|^2}{y_{Y,0}}
    \frac{x_0}{x_a}
    \frac{1}{1+2\kappa_p/y_{Y,0}}
    \left[
        \frac{M_p^2}{y_{Y,0}+\kappa_p}
        +\frac{V_p}{2}
        +\frac{y_{Y,0}}{y_{Y,0}+\kappa_p}M_p\hat{\alpha}_p
        +\frac{y_{Y,0}}{2}\hat{\alpha}_p^2
    \right].
\end{equation}

\paragraph{Aggregate market elasticity.}
The aggregate market elasticity is given by
\begin{equation}
    \epsilon_{M}^{-1}
    = \frac{1}{p(X,\epsilon)}\frac{\partial p(X,\epsilon)}{\partial F(X,\epsilon)}
    = -\frac{1}{y_{Y,0}}
    \frac{\partial y_{Y,2}^{T}}{\partial \hat{\alpha}_p}
    \frac{\epsilon}{x_0}
    + \mathcal{O}(\epsilon^2).
\end{equation}
Using the closed-form slope
\begin{equation}
    G_p(x,\hat{\alpha}_p)
    \equiv
    \frac{\partial y_{Y,2}^{T}}{\partial \hat{\alpha}_p}
    =2d_Y^T\left(\hat{\alpha}_p+\frac{M_p}{y_{Y,0}+\kappa_p}\right),
\end{equation}
we obtain
\begin{equation}
    \epsilon_{M}^{-1}
    =-\frac{G_p}{y_{Y,0}}\frac{\epsilon}{x_0}
    +\mathcal{O}(\epsilon^2).
\end{equation}
Substituting the expression for $G_p$ and using
\begin{equation}
    d_Y^T
    =\frac{1}{1+2\kappa_p/y_{Y,0}}d_Y^0,
    \qquad
    d_Y^0
    =\frac{1}{2}(1-\psi^{-1})\gamma\|\sigma\|^2\frac{x_0}{x_a},
\end{equation}
gives
\begin{equation}
    \epsilon_{M}^{-1}
    =-
    \left(1-\psi^{-1}\right)
    \frac{\gamma\|\sigma\|^2}{y_{Y,0}}
    \frac{1}{x_a}
    \frac{1}{1+2\kappa_p/y_{Y,0}}
    \left(\hat{\alpha}_p+\frac{M_p}{y_{Y,0}+\kappa_p}\right)
    \epsilon
    +\mathcal{O}(\epsilon^2).
\end{equation}
Finally, using
\begin{equation}
    \overline{\alpha}_p=1+\hat{\alpha}_p\epsilon,
    \qquad
    M_p=(1-\gamma)\sigma_p\sigma^\top,
\end{equation}
where the second equality follows from the assumption $\hat{\alpha}=0$, we obtain
\begin{equation}\label{eq:elasticity-timevarying-noheterog-final}
    \epsilon_{M}^{-1}
    =
    \left(1-\psi^{-1}\right)
    \frac{\gamma\|\sigma\|^2}{y_{Y,0}+2\kappa_p}
    \frac{1}{x_a}
    \left[
        1-\overline{\alpha}_p
        +\frac{(\gamma-1)\sigma_p\epsilon\sigma^\top}{y_{Y,0}+\kappa_p}
    \right]
    +\mathcal{O}(\epsilon^2).
\end{equation}
Thus, mean reversion in passive demand attenuates the aggregate market elasticity through the factor $(y_{Y,0}+2\kappa_p)^{-1}$, while passive-demand risk introduces an additional correction proportional to $(\gamma-1)\sigma_p\sigma^\top$.

\paragraph{Return volatility.}
The return volatility is given by
\begin{equation}
    \sigma_{R_Y}
    =
    \sigma-\sigma_{y_Y},
\end{equation}
where
\begin{equation}
    \sigma_{y_Y}
    =
    \frac{y_{Y,x,t}}{y_{Y,t}}\sigma_{x,t}
    +
    \frac{y_{Y,\overline{\alpha}_p,t}}{y_{Y,t}}\sigma_p
    =
    \frac{y_{Y,2,\hat{\alpha}_p}}{y_{Y,0}}
    \sigma_p\epsilon^2
    +
    \mathcal O(\epsilon^3).
\end{equation}
Using the expression for the aggregate market elasticity,
\begin{equation}
    \epsilon_M^{-1}
    =
    -\frac{1}{y_{Y,0}}
    \frac{\partial y_{Y,2}^T}{\partial\hat{\alpha}_p}
    \frac{\epsilon}{x_0},
\end{equation}
we obtain
\begin{equation}
    \sigma_{R_Y}
    =
    \sigma
    +
    x_0\epsilon_M^{-1}\sigma_p\epsilon
    +
    \mathcal O(\epsilon^3).
\end{equation}

\paragraph{Risk exposure.} The risk exposure of passive investors is given by
\begin{equation}
    \sigma_0 = \overline{\alpha}_p \sigma_{R_Y} = \sigma + \sigma \hat{\alpha}_p \epsilon - \sigma_{y_Y,2} \epsilon^2 + \mathcal{O}(\epsilon^3), \qquad
\end{equation}
From the expression for $\sigma_{y_Y,2}$, we have
\begin{equation}
    \sigma_0 =  \sigma + \hat{\alpha}_p \sigma \epsilon + x_0\epsilon_M^{-1}\sigma_p\epsilon+ \mathcal{O}(\epsilon^3), \qquad
\end{equation}

Using the fact that $\sum_{j=0}^2 x_j\sigma_j = \sigma_{R_Y}$, we can write the risk exposure of active investor $j=1,2$ as
\begin{equation}
    \sigma_j = \sigma - \frac{x_0\hat{\alpha}_p}{x_a} \sigma  \epsilon +x_0\epsilon_M^{-1}\sigma_p\epsilon + \mathcal{O}(\epsilon^3), \qquad j=1,2.
\end{equation}

We then write the first-order correction to the risk exposure as $\sigma_{j,1}=m_j \sigma$, where
\begin{equation}
    m_{a} = - \frac{x_0\hat{\alpha}_p}{x_a}, \qquad \qquad 
    m_{0} = \hat{\alpha}_p,
\end{equation}
where $m_{1} = m_{2} = m_a$.


\paragraph{Consumption-wealth ratio.}
The consumption-wealth ratio of investor $j=0,1,2$ can be written as
\begin{equation}
    c_j = c_j^0 + \Delta c_j,
\end{equation}
where $c_j^0$ is the fixed-$\hat{\alpha}_p$ consumption-wealth ratio and $\Delta c_j$ is the time-varying-$\alpha_p$ correction.

The consumption-wealth ratio in the fixed-$\hat{\alpha}_p$ economy is given by
\begin{equation}
    c_{j}^0 = c_{j,0} + c_{j,2}^0 \epsilon^2 + \mathcal{O}(\epsilon^3),
\end{equation}
where 
\begin{equation}
    c_{j,2}^0 = (1-\psi) \left[ y_{Y,2}^0 + \left(m_{j}m_{a}  - \frac{1}{2}m_j^2  \right)\gamma \|\sigma\|^2 \right].
\end{equation}

We can then write $c_{j,2}^0 = d_{j}^0 \hat{\alpha}_p^2$, where
\begin{align}
    d_{a}^0 &= (1-\psi)\frac{\gamma \|\sigma\|^2}{2} \left[ (1-\psi^{-1}) \frac{x_0}{x_a} + \left( \frac{x_0}{x_a} \right)^2  \right], \qquad 
    d_{0}^0 = (1-\psi)\gamma \|\sigma\|^2 \left[ \frac{1-\psi^{-1}}{2} \frac{x_0}{x_a} - \left( \frac{x_0}{x_a} +\frac{1}{2}\right)  \right].
\end{align}

The individual consumption-wealth ratios solve the same coefficient equation as the aggregate dividend yield, but investor by investor.
In the no-heterogeneity, no-constraint case, active investors are identical, so write
\begin{equation}
    c_{0,2}^{T}(x,\hat{\alpha}_p)
    =a_0^T+b_0^T\hat{\alpha}_p+d_0^T\hat{\alpha}_p^2,
    \qquad
    c_{a,2}^{T}(x,\hat{\alpha}_p)
    =a_a^T+b_a^T\hat{\alpha}_p+d_a^T\hat{\alpha}_p^2,
\end{equation}
where $c_{a,2}^{T}$ denotes the common correction for active investors $j=1,2$.
The coefficients are obtained from the fixed-$\hat{\alpha}_p$ coefficients using
\begin{align}
    d_j^T
    &=\frac{1}{1+2\kappa_p/y_{Y,0}}d_j^0,\\
    b_j^T
    &=\frac{2d_j^TM_p}{y_{Y,0}+\kappa_p},\\
    a_j^T
    &=\frac{1}{y_{Y,0}}\left(b_j^TM_p+d_j^TV_p\right),
    \qquad j\in\{0,a\},
\end{align}
where we use $c_{j,0}=y_{Y,0}$ in the homogeneous benchmark economy.

\paragraph{Consumption-wealth ratio drift and diffusion.}
The slope of each individual consumption-wealth ratio with respect to the passive portfolio share is
\begin{equation}\label{eq:Gcj-noheterog}
    G_j(x,\hat{\alpha}_p)
    \equiv \frac{\partial c_{j,2}^{T}}{\partial\hat{\alpha}_p}
    =2d_j^T\left(\hat{\alpha}_p+\frac{M_p}{y_{Y,0}+\kappa_p}\right),
    \qquad j\in\{0,a\}.
\end{equation}
Hence the diffusion of the consumption-wealth ratio of investor $j$ is
\begin{equation}\label{eq:sigmacj-noheterog}
    \sigma_{c_j,t}
    =\frac{G_j}{y_{Y,0}}\sigma_p\epsilon^2+\mathcal{O}(\epsilon^3),
    \qquad j\in\{0,a\}.
\end{equation}

The drift of the consumption-wealth ratio is computed in the same way.
Since
\begin{equation}
    c_{j,2,\hat{\alpha}_p}^{T}=G_j,
    \qquad
    c_{j,2,\hat{\alpha}_p\hat{\alpha}_p}^{T}=2d_j^T,
\end{equation}
we have
\begin{equation}\label{eq:mucj-noheterog}
    \mu_{c_j,t}
    =\left[
        \frac{G_j}{y_{Y,0}}\kappa_p(\hat{\alpha}-\hat{\alpha}_p)
        +\frac{d_j^T}{y_{Y,0}}V_p
    \right]\epsilon^2
    +\mathcal{O}(\epsilon^3),
    \qquad j\in\{0,a\}.
\end{equation}
Under the maintained normalization $\hat{\alpha}=0$, this becomes
\begin{equation}\label{eq:mucj-noheterog-alpha-zero}
    \mu_{c_j,t}
    =\left[
        -\frac{G_j}{y_{Y,0}}\kappa_p\hat{\alpha}_p
        +\frac{d_j^T}{y_{Y,0}}V_p
    \right]\epsilon^2
    +\mathcal{O}(\epsilon^3),
    \qquad j\in\{0,a\}.
\end{equation}

By contrast, the diffusion of the aggregate dividend yield is
\begin{equation}\label{eq:sigmay-noheterog-from-Gp}
    \sigma_{y_Y,t}
    =\frac{G_p}{y_{Y,0}}\sigma_p\epsilon^2+\mathcal{O}(\epsilon^3),
    \qquad
    G_p=x_0G_0+x_aG_a.
\end{equation}
Thus, aggregation implies
\begin{equation}
    x_0\sigma_{c_0,t}+x_a\sigma_{c_a,t}=\sigma_{y_Y,t},
\end{equation}


\paragraph{Market price of risk.}
The market price of risk can be written as
\begin{equation}
    \eta_t = \eta_0 + \eta_1 \epsilon + \eta_2 \epsilon^2 + \mathcal{O}(\epsilon^3).
\end{equation}
At the benchmark,
\begin{equation}
    \eta_0=\gamma\sigma.
\end{equation}
At first order, active investors absorb the residual portfolio exposure created by passive investors.
Since
\begin{equation}
    m_a=-\frac{x_0}{x_a}\hat{\alpha}_p,
\end{equation}
we have
\begin{equation}
    \eta_1=\gamma m_a\sigma
    =-\gamma\sigma\frac{x_0}{x_a}\hat{\alpha}_p.
\end{equation}

We next derive the second-order term.
Because active investors are unconstrained and identical, their optimal exposure satisfies
\begin{equation}
    \sigma_{a,t}
    =\frac{\eta_t}{\gamma}
    +\frac{1-\gamma^{-1}}{\psi-1}\sigma_{c_a,t}.
\end{equation}
Therefore,
\begin{equation}\label{eq:eta2-noheterog-from-active-exposure}
    \eta_2
    =\gamma\left[
        \sigma_{a,2}
        -\frac{1-\gamma^{-1}}{\psi-1}\sigma_{c_a,2}
    \right].
\end{equation}
where
\begin{equation}
    \sigma_{a,2}=-\sigma_{y_Y,2}.
\end{equation}
Substituting into \eqref{eq:eta2-noheterog-from-active-exposure} gives
\begin{equation}\label{eq:eta2-noheterog-total}
    \eta_2
    =-\gamma\left[
        \sigma_{y_Y,2}
        +\frac{1-\gamma^{-1}}{\psi-1}\sigma_{c_a,2}
    \right].
\end{equation}
Using the diffusion formulas above,
\begin{equation}
    \sigma_{y_Y,2}=\frac{G_p}{y_{Y,0}}\sigma_p,
    \qquad
    \sigma_{c_a,2}=\frac{G_a}{y_{Y,0}}\sigma_p,
\end{equation}
we obtain the closed-form expression
\begin{equation}\label{eq:eta2-noheterog-closed-form}
    \eta_2
    =-\frac{\gamma}{y_{Y,0}}
    \left[
        G_p+\frac{1-\gamma^{-1}}{\psi-1}G_a
    \right]\sigma_p.
\end{equation}
Equivalently, using the explicit slopes
\begin{equation}
    G_p=2d_Y^T\left(\hat{\alpha}_p+\frac{M_p}{y_{Y,0}+\kappa_p}\right),
    \qquad
    G_a=2d_a^T\left(\hat{\alpha}_p+\frac{M_p}{y_{Y,0}+\kappa_p}\right),
\end{equation}
we can write
\begin{equation}\label{eq:eta2-noheterog-slopes}
    \eta_2
    =-\frac{2\gamma}{y_{Y,0}}
    \left[
        d_Y^T+\frac{1-\gamma^{-1}}{\psi-1}d_a^T
    \right]
    \left(\hat{\alpha}_p+\frac{M_p}{y_{Y,0}+\kappa_p}\right)\sigma_p.
\end{equation}

It is useful to express this result in terms of the aggregate market elasticity.
Because $d_a^T$ and $d_Y^T$ are attenuated by the same factor, we have
\begin{equation}
    \frac{G_a}{G_p}=\frac{d_a^T}{d_Y^T}=\frac{d_a^0}{d_Y^0}.
\end{equation}
Using the expressions for $d_a^0$ and $d_Y^0$, this ratio is
\begin{equation}\label{eq:Ga-Gp-ratio-noheterog}
    \chi_a
    \equiv \frac{G_a}{G_p}
    =\frac{d_a^0}{d_Y^0}
    =(1-\psi)
    \left[
        1+\frac{x_0/x_a}{1-\psi^{-1}}
    \right].
\end{equation}
Therefore, \eqref{eq:eta2-noheterog-closed-form} can be written as
\begin{equation}\label{eq:eta2-noheterog-Gp}
    \eta_2
    =-\frac{\gamma G_p}{y_{Y,0}}
    \left[
        1+\frac{1-\gamma^{-1}}{\psi-1}\chi_a
    \right]\sigma_p.
\end{equation}
Using the expression for $\chi_a$,
\begin{equation}
        \frac{1-\gamma^{-1}}{\psi-1}\chi_a
    =
    -\frac{1-\gamma^{-1}}{1-\psi^{-1}} \left(\frac{1}{x_a} - \psi^{-1}\right)
\end{equation}


Finally, using the aggregate market elasticity relation
\begin{equation}
    \epsilon_M^{-1}
    =-\frac{G_p}{y_{Y,0}}\frac{\epsilon}{x_0}
    +\mathcal{O}(\epsilon^2),
\end{equation}
we obtain
\begin{equation}\label{eq:eta2-noheterog-elasticity}
    \eta
    = \gamma \sigma \left( \frac{1- x_0 \overline{\alpha}_p}{x_a}\right)
    +\left[
        1 -\frac{1-\gamma^{-1}}{1-\psi^{-1}} \left(\frac{1}{x_a} - \psi^{-1}\right)
    \right]\gamma x_0\epsilon_M^{-1}\sigma_p\epsilon
    +\mathcal{O}(\epsilon^3).
\end{equation}

This expression shows that the second-order market price of risk is proportional to the aggregate market elasticity, with an adjustment for the active investors' hedging demand.

\paragraph{Risk premium.}
The risk premium on the aggregate claim is
\begin{equation}
    \pi_Y=\sigma_{R_Y}\eta^\top.
\end{equation}
Using the expressions derived above,
\begin{equation}
    \sigma_{R_Y}
    =
    \sigma+x_0\epsilon_M^{-1}\sigma_p\epsilon
    +\mathcal O(\epsilon^3),
\end{equation}
and
\begin{equation}
    \eta
    =
    \gamma\sigma\left(\frac{1-x_0\overline{\alpha}_p}{x_a}\right)
    +
    (1-H_a) \gamma x_0\epsilon_M^{-1}\sigma_p\epsilon
    +\mathcal O(\epsilon^3),
\end{equation}
where
\begin{equation}
    H_a
    \equiv
    \frac{1-\gamma^{-1}}{1-\psi^{-1}} \left(\frac{1}{x_a} - \psi^{-1}\right).
\end{equation}
Therefore,
\begin{align}
    \pi_Y
    &=
    \gamma\|\sigma\|^2
    \left(\frac{1-x_0\overline{\alpha}_p}{x_a}\right) 
    +x_0\epsilon_M^{-1}\epsilon
    \left[
        2-H_a
    \right] \gamma \sigma\sigma_p^\top
    +\mathcal O(\epsilon^3).
\end{align}

The term $H_a$ captures the hedging-demand response of active investors. 
Under the usual assumptions $\gamma>1$ and $\psi>1$, $H_a>0$, so hedging demand dampens the effect of passive-demand risk on the aggregate risk premium.

\paragraph{Interest rate.}
The interest rate follows from
\begin{equation}
    r=y_Y-\pi_Y-\mu_{y_Y}-\sigma\sigma_{y_Y}^\top .
\end{equation}
Using the expressions above,
\begin{align}
    r
    &=y_{Y,0}
    -\gamma\|\sigma\|^2
    \left(\frac{1-x_0\overline{\alpha}_p}{x_a}\right)
    +y_{Y,2}^T\epsilon^2
    -\mu_{y_Y} 
    -\left[(2-H_a)\gamma-1\right]
    x_0\epsilon_M^{-1}\sigma\sigma_p^\top\epsilon
    +\mathcal O(\epsilon^3),
\end{align}
where
\begin{equation}
    \mu_{y_Y}
    =
    \left[
        -\frac{G_p}{y_{Y,0}}\kappa_p\hat{\alpha}_p
        +\frac{d_Y^T}{y_{Y,0}}V_p
    \right]\epsilon^2
    +\mathcal O(\epsilon^3).
\end{equation}

Equivalently, isolating only the incremental effect of time-varying passive demand, Step 5 gives
\begin{equation}
    \Delta r_2=-\sigma\Delta\eta_2^\top.
\end{equation}
Since
\begin{equation}
    \Delta\eta_2\epsilon^2
    =(1-H_a)\gamma x_0\epsilon_M^{-1}\sigma_p\epsilon
    +\mathcal O(\epsilon^3),
\end{equation}
we obtain
\begin{equation}
    \Delta r_2\epsilon^2
    =-(1-H_a)\gamma x_0\epsilon_M^{-1}
    \sigma\sigma_p^\top\epsilon
    +\mathcal O(\epsilon^3).
\end{equation}
Equivalently,
\begin{equation}
    \Delta r_2
    =(1-H_a)\gamma
    \frac{G_p}{y_{Y,0}}\sigma\sigma_p^\top.
\end{equation}


\paragraph{Price of dividend claim.}
We now price the dividend claim, whose payoff is $D_t=s_tY_t$.
Write
\begin{equation}
    s_t=\bar{s}+\epsilon\hat{s}_t,
\end{equation}
and conjecture the expansion
\begin{equation}
    p_D(x,s,\overline{\alpha}_p;\epsilon)
    =p_{D,0}+p_{D,1}(\hat{s})\epsilon+p_{D,2}^T(\hat{s},\hat{\alpha}_p)\epsilon^2+\mathcal{O}(\epsilon^3).
\end{equation}
At the benchmark,
\begin{equation}
    p_{D,0}=\frac{\bar{s}}{\delta},
    \qquad
    \delta\equiv y_{Y,0}.
\end{equation}
The first-order term is the present value of the transitory dividend-share component:
\begin{equation}
    p_{D,1}(\hat{s})=\frac{\hat{s}}{\delta+\kappa_s}.
\end{equation}

To obtain the second-order term, it is useful to decompose the dividend claim into an aggregate-claim component and a dividend-specific component.
The aggregate component is simply $\bar{s}$ times the aggregate claim.
Because
\begin{equation}
    p_Y=\frac{1}{y_Y}
    =\frac{1}{\delta}-\frac{y_{Y,2}^T}{\delta^2}\epsilon^2+\mathcal{O}(\epsilon^3),
\end{equation}
this component contributes
\begin{equation}
    -\frac{\bar{s}}{\delta^2}y_{Y,2}^T
\end{equation}
to $p_{D,2}^T$.

The dividend-specific component comes from the risk premium on the dividend-share shock.
The fixed-$\hat{\alpha}_p$ calculation gives
\begin{equation}
    p_{D,2}^{0}(\hat{s},\hat{\alpha}_p)
    =-\frac{\bar{s}}{\delta^2}y_{Y,2}^{0}(\hat{\alpha}_p)
    -\frac{B_s}{\delta(\delta+\kappa_s)}\eta_1\sigma_s^\top
    +b_2^0\hat{s},
\end{equation}
where $B_s\equiv\bar{s}(1-\bar{s})$ and $b_2^0$ is independent of $\hat{\alpha}_p$.
In the present special case,
\begin{equation}
    \eta_1=-\gamma\sigma\frac{x_0}{x_a}\hat{\alpha}_p,
\end{equation}
so the fixed-$\hat{\alpha}_p$ excess component is linear in $\hat{\alpha}_p$:
\begin{equation}
    e_D^0(\hat{\alpha}_p)
    \equiv
    -\frac{B_s}{\delta(\delta+\kappa_s)}\eta_1\sigma_s^\top
    =\frac{\gamma B_s}{\delta(\delta+\kappa_s)}\frac{x_0}{x_a}\sigma\sigma_s^\top\hat{\alpha}_p.
\end{equation}

With time-varying passive portfolios, this excess component solves the same scalar coefficient equation as the aggregate dividend-yield correction.
Because $e_D^0$ is linear in $\hat{\alpha}_p$, there is no quadratic term and hence no new $M_p$ term in its slope.
Mean reversion attenuates the linear coefficient by $\delta/(\delta+\kappa_p)$.
Thus the total excess component is
\begin{equation}
    e_D^T(\hat{\alpha}_p)
    =\frac{\delta}{\delta+\kappa_p}e_D^0(\hat{\alpha}_p)
    +\frac{1}{\delta+\kappa_p}
    \frac{\gamma B_s}{\delta(\delta+\kappa_s)}\frac{x_0}{x_a}\sigma\sigma_s^\top M_p,
\end{equation}
where the second term is the constant term generated by the drift and risk-adjustment operator.
Equivalently,
\begin{equation}\label{eq:eD-total-noheterog}
    e_D^T(\hat{\alpha}_p)
    =\frac{\gamma B_s}{(\delta+\kappa_p)(\delta+\kappa_s)}\frac{x_0}{x_a}\sigma\sigma_s^\top
    \left(\hat{\alpha}_p+\frac{M_p}{\delta}\right).
\end{equation}

Combining the aggregate and dividend-specific components, the second-order price of the dividend claim is
\begin{equation}\label{eq:pD2-total-noheterog}
    p_{D,2}^T(\hat{s},\hat{\alpha}_p)
    =-\frac{\bar{s}}{\delta^2}y_{Y,2}^T(\hat{\alpha}_p)
    +e_D^T(\hat{\alpha}_p)
    +b_2^0\hat{s}.
\end{equation}
Therefore,
\begin{equation}\label{eq:pD-total-noheterog}
    p_D(x,s,\overline{\alpha}_p;\epsilon)
    =\frac{\bar{s}}{\delta}
    +\frac{\hat{s}}{\delta+\kappa_s}\epsilon
    +\left[-\frac{\bar{s}}{\delta^2}y_{Y,2}^T(\hat{\alpha}_p)
    +e_D^T(\hat{\alpha}_p)+b_2^0\hat{s}\right]\epsilon^2
    +\mathcal{O}(\epsilon^3).
\end{equation}

Differentiating with respect to passive flows gives the price impact of passive demand on the dividend claim:
\begin{align}\label{eq:pD-price-impact-noheterog-final}
    \frac{\partial p_D}{\partial F}
    &=\bar{s}\frac{\partial p_Y}{\partial F}
    +\frac{\epsilon}{x_0}\frac{\partial e_D^T}{\partial\hat{\alpha}_p}
    +\mathcal{O}(\epsilon^2)\\
    &=\bar{s}\frac{\partial p_Y}{\partial F}
    +\frac{\epsilon}{x_0}
    \frac{\gamma B_s}{(\delta+\kappa_p)(\delta+\kappa_s)}\frac{x_0}{x_a}\sigma\sigma_s^\top
    +\mathcal{O}(\epsilon^2).
\end{align}

Equivalently, using $p_{D,0}=\bar{s}/\delta$ and $p_{Y,0}=1/\delta$, the normalized price impact of the dividend claim can be related to the normalized price impact of the aggregate claim as
\begin{equation}\label{eq:pD-pY-normalized-price-impact-noheterog}
    \frac{1}{p_D}\frac{\partial p_D}{\partial F}
    =
    \frac{1}{p_Y}\frac{\partial p_Y}{\partial F}
    +\frac{\delta}{\bar{s}}
    \frac{\epsilon}{x_0}
    \frac{\gamma B_s}{(\delta+\kappa_p)(\delta+\kappa_s)}
    \frac{x_0}{x_a}\sigma\sigma_s^\top
    +\mathcal{O}(\epsilon^2).
\end{equation}
Since $B_s=\bar{s}(1-\bar{s})$, this simplifies to
\begin{equation}\label{eq:pD-pY-normalized-price-impact-noheterog-simplified}
    \frac{1}{p_D}\frac{\partial p_D}{\partial F}
    =
    \frac{1}{p_Y}\frac{\partial p_Y}{\partial F}
    + \frac{1}{\bar{s}} \frac{\delta\gamma\bar{s}(1-\bar{s})\sigma_s\sigma^\top}{(\delta+\kappa_p)(\delta+\kappa_s)}
    \frac{\epsilon}{x_a}
    +\mathcal{O}(\epsilon^2).
\end{equation}

Thus, the dividend-claim price impact is the sum of two terms.
The first term is the mechanical exposure to the aggregate claim.
The second term is the dividend-specific risk-premium channel, attenuated by passive-demand mean reversion through the factor $(\delta+\kappa_p)^{-1}$.
\end{proof}


\subsection{The role of the wealth distribution}

The aggregate market elasticity computed above is a partial derivative: it changes the passive portfolio share while holding the wealth distribution fixed.
A portfolio-flow shock, however, also changes wealth shares because passive and active investors have different exposures to the risky claim.
We now show that, in the special case without preference heterogeneity or leverage constraints, this wealth-distribution channel is one order smaller than the direct portfolio-flow channel.

Let $x_0$ denote the wealth share of passive investors and $x_a=1-x_0$ the wealth share of active investors.
Because passive investors hold the portfolio share $\overline{\alpha}_p$, their wealth-share diffusion is
\begin{equation}\label{eq:sigmax0-portfolio-flow}
    \sigma_{x_0,t}
    =x_{0,t}\left(\sigma_{0,t}-\sigma_{R_Y,t}\right)
    =x_{0,t}\left(\overline{\alpha}_{p,t}-1\right)\sigma_{R_Y,t}.
\end{equation}
Using $\overline{\alpha}_{p,t}=1+\hat{\alpha}_{p,t}\epsilon$ and $\sigma_{R_Y,t}=\sigma+\mathcal{O}(\epsilon^2)$, this gives
\begin{equation}\label{eq:sigmax0-leading}
    \sigma_{x_0,t}
    =x_{0,t}\hat{\alpha}_{p,t}\sigma\epsilon+\mathcal{O}(\epsilon^3).
\end{equation}
Thus, the wealth-distribution state moves in response to the same aggregate return shock that redistributes wealth between passive and active investors.

The aggregate dividend yield satisfies
\begin{equation}
    y_Y(x_0,\hat{\alpha}_p;\epsilon)
    =y_{Y,0}+y_{Y,2}^{T}(x_0,\hat{\alpha}_p)\epsilon^2+\mathcal{O}(\epsilon^3).
\end{equation}
Therefore, allowing both $x_0$ and $\hat{\alpha}_p$ to move, its diffusion is
\begin{equation}\label{eq:sigmay-total-wealth-dist}
    \sigma_{y_Y,t}
    =\frac{1}{y_{Y,0}}
    \left[
        y_{Y,2,\hat{\alpha}_p}^{T}\sigma_p
        +y_{Y,2,x_0}^{T}\sigma_{x_0,t}
    \right]\epsilon^2
    +\mathcal{O}(\epsilon^4).
\end{equation}
The first term is the direct portfolio-flow channel studied above.
The second term is the wealth-distribution channel.
Using \eqref{eq:sigmax0-leading},
\begin{equation}\label{eq:sigmay-total-wealth-dist-expanded}
    \sigma_{y_Y,t}
    =\frac{G_p}{y_{Y,0}}\sigma_p\epsilon^2
    +\frac{x_0\hat{\alpha}_p}{y_{Y,0}}y_{Y,2,x_0}^{T}\sigma\epsilon^3
    +\mathcal{O}(\epsilon^4),
\end{equation}
where
\begin{equation}
    G_p\equiv y_{Y,2,\hat{\alpha}_p}^{T}.
\end{equation}
Consequently, the price-dividend-ratio diffusion induced by a portfolio-flow shock is
\begin{equation}\label{eq:sigmapY-total-wealth-dist}
    \sigma_{p_Y,t}
    =-\sigma_{y_Y,t}
    =-\frac{G_p}{y_{Y,0}}\sigma_p\epsilon^2
    -\frac{x_0\hat{\alpha}_p}{y_{Y,0}}y_{Y,2,x_0}^{T}\sigma\epsilon^3
    +\mathcal{O}(\epsilon^4).
\end{equation}
Using the aggregate market elasticity relation
\begin{equation}
    \epsilon_M^{-1}
    =-\frac{G_p}{y_{Y,0}}\frac{\epsilon}{x_0}+\mathcal{O}(\epsilon^2),
\end{equation}
we can write
\begin{equation}\label{eq:sigmapY-total-wealth-dist-elasticity}
    \sigma_{p_Y,t}
    =x_0\epsilon_M^{-1}\sigma_p\epsilon
    -\frac{x_0\hat{\alpha}_p}{y_{Y,0}}y_{Y,2,x_0}^{T}\sigma\epsilon^3
    +\mathcal{O}(\epsilon^4).
\end{equation}
Thus, to the order at which the aggregate market elasticity is computed, the equilibrium price response to a portfolio-flow shock is the direct partial-equilibrium effect holding the wealth distribution fixed.
The endogenous movement in the wealth distribution contributes only at the next order.

In the closed-form no-heterogeneity solution, $y_{Y,2}^{T}$ depends on the wealth distribution only through $x_0/x_a$.
Hence
\begin{equation}\label{eq:y2x0-noheterog}
    y_{Y,2,x_0}^{T}
    =\frac{y_{Y,2}^{T}}{x_0x_a}.
\end{equation}
Substituting \eqref{eq:y2x0-noheterog} into \eqref{eq:sigmapY-total-wealth-dist-elasticity} gives the explicit wealth-distribution correction
\begin{equation}\label{eq:sigmapY-total-wealth-dist-final}
    \sigma_{p_Y,t}
    =x_0\epsilon_M^{-1}\sigma_p\epsilon
    -\frac{\hat{\alpha}_p}{x_a}\frac{y_{Y,2}^{T}}{y_{Y,0}}\sigma\epsilon^3
    +\mathcal{O}(\epsilon^4).
\end{equation}
Equation~\eqref{eq:sigmapY-total-wealth-dist-final} separates the two effects of a portfolio-flow shock.
The first term is the direct demand effect captured by the aggregate market elasticity.
The second term is the equilibrium wealth-distribution effect: the price response to the change in $x_0$ caused by the shock itself.
Because this second term is of order $\epsilon^3$ in price diffusion, it affects the flow-price impact only at order $\epsilon^2$ and does not change the first-order aggregate elasticity formula above.



% The dividend yield, $y$, is given by
% \begin{align}
%     y(X,\epsilon) &= \psi \rho + (1-\psi) \left[ r(X,\epsilon) + \theta(X,\epsilon) \|\sigma_{R}(X,\epsilon)\| - \sum_{j=0}^J x_{j}\frac{\gamma_j}{2} \sigma_{j}^2(X,\epsilon) \right] \nonumber\\
%     &\qquad +\sum_{j=0}^J x_{j} \left[\mu_{c_j}(X,\epsilon) + (1-\gamma_j)  \sigma_{c_j}(X,\epsilon) \frac{\sigma_{R}(X;\epsilon)^\top}{\|\sigma_{R}(X;\epsilon)\|} \sigma_{j}(X,\epsilon) + \frac{\psi - \gamma_j}{1-\psi} \frac{\|\sigma_{c_j}(X,\epsilon)\|^2}{2}\right].
% \end{align}

% The second-order term in the expansion of $y(X,\epsilon)$ is given by
% \begin{small}
% 	\begin{align}
% 	y_2(X) &=  (1-\psi) \left[ r_2(X) + \theta_2(X) \|\sigma\| + \theta_0(X) \sigma_{R,2}(X) \frac{\sigma^\top}{\|\sigma\|} - \sum_{j=0}^J x_{j}\frac{\gamma}{2} \left( \sigma_{j,1}^2(X) + 2 \sigma_{j,0} \sigma_{j,2}(X) + 2 \hat{\gamma}_j \sigma_{j,0} \sigma_{j,1}(X)  \right)\right] \nonumber\\
% 	&\qquad +\sum_{j=0}^J x_{j} \left[\mu_{c_j,2}(X) + (1-\gamma)  \sigma_{c_j,2}(X)\sigma^\top\right].
% 	\end{align}
% \end{small}


% Using the expression for the interest rate, we obtain
% \begin{align}
%     y_2(X) &=  \left(1-\psi^{-1}\right) \left[ \mu_{y,2}(X) + \sigma_{y,2}(X) \sigma^\top + \sum_{j=0}^J x_{j}\gamma \|\sigma\| \left( \frac{\sigma_{j,1}^2(X)}{2 \|\sigma\|} + \sigma_{j,2}(X) + \hat{\gamma}_j \sigma_{j,1}(X)  \right)\right] \nonumber\\
%     &\qquad +\psi^{-1}\sum_{j=0}^J x_{j} \left[\mu_{c_j,2}(X) + (1-\gamma)  \sigma_{c_j,2}(X) \sigma^\top\right].
% \end{align}

% Given that $\mu_{c_j,2} = (1-\psi) \mu_{y,2} $ and $\sigma_{c_j,2} = (1-\psi) \sigma_{y,2} $, we obtain
% \begin{align}
%     y_2(X) &=  \left(1-\psi^{-1}\right) \left[\gamma \sigma_{y,2}(X) \sigma^\top+\sum_{j=0}^J x_{j}\gamma \|\sigma\| \left( \frac{\sigma_{j,1}^2(X)}{2 \|\sigma\|} + \sigma_{j,2}(X) + \hat{\gamma}_j \sigma_{j,1}(X)  \right)\right] \nonumber\\
%      &=\left(1-\psi^{-1}\right) \sum_{j=0}^J x_{j}\gamma \|\sigma\| \left( \frac{\sigma_{j,1}^2(X)}{2 \|\sigma\|}  + \hat{\gamma}_j \sigma_{j,1}(X)  \right),    
% \end{align}
%  where in the second equality, we use the fact that, from the market clearing condition for the risky asset, we have $\sum_{j=0}^J x_j\sigma_{j,2}(X)= -\sigma_{y,2}(X) \frac{\sigma^\top}{\|\sigma\|}.$

%  \noindent\textbf{Step 7: Risk premium.}\\
% Since the risk premium is given by $ \pi_t=\theta_t\|\sigma_{R,t}\| $, we can write
% \begin{equation*}
% \pi(X,\epsilon) = \theta_0(X)\|\sigma\| + \theta_1(X)\|\sigma\|\epsilon+
% \left(-\theta_{0}(X) \sigma_{y,2}(X) \frac{\sigma^\top}{\|\sigma\|} +\theta_2(X)\|\sigma\|\right)\epsilon^2+\mathcal{O}\left(\epsilon^3\right).
% \end{equation*}
% We can then write
% \begin{equation}\label{eq:rp-second-order}
% 	\pi_2(X) = -\theta_{0}(X)\sigma_{y, 2}(X) \frac{\sigma^\top}{\|\sigma\|}+\theta_2(X)\|\sigma\|.
% \end{equation}


% \end{proof}






% \begin{proof}
% \noindent\textbf{Step1: Laws of motion for $\xi_j$ and $y$.}\\
% We start by considering the diffusion terms for $\xi_j$ and $y$. Given that we are abstracting from portfolio shocks, we assume that $d=1$ without further loss of generality, so that we can treat diffusion terms as scalars. The expansion of $\sigma_{\xi_j,t}$ in $\epsilon$ is given by
% \begin{align}
%     \sigma_{\xi_j}(X,\epsilon) &= \frac{\xi_{j,X}(X,\epsilon)}{\xi_j(X,\epsilon)} \sigma_{X}(X,\epsilon) \nonumber\\
%     &= \frac{\xi_{j,X,1}(X)}{\xi_{j,0}(X)} \sigma_{X,1}(X) \epsilon^2+\mathcal{O}(\epsilon^3)\nonumber\\
%     &= -(\psi-1)\left(1-\psi^{-1}\right)\frac{\gamma \sigma^2}{2 \xi_{j,0}(X)} \sum_{k=1}^J \left( \hat{\gamma}_k - \hat{\gamma}_0  \right) x_k\sigma_{k,1}(X) \epsilon^2 + \mathcal{O}(\epsilon^3).
% \end{align}

% Notice that $\sigma_{\xi_j,2}(X)$  does not depend on $j$, that is, it is the same for all investors. Moreover, $\sigma_{\xi_{j},2}(X)> 0$, as $\sigma_{k,1}(X)$ is inversely related to $\hat{\gamma}_k$. 

% Similarly, the diffusion for dividend yield $y$ can be written as
% \begin{align}\label{eq:sig-y2}
%     \sigma_{y}(X,\epsilon) &= \frac{y_{X,1}(X)}{y_{0}(X)} \sigma_{X,1}(X) \epsilon^2+\mathcal{O}(\epsilon^2)\nonumber\\
%     &= \left(1-\psi^{-1}\right)\frac{\gamma\sigma^2}{2 y_{0}(X)} \sum_{k=1}^J \left( \hat{\gamma}_k - \hat{\gamma}_0  \right) x_k \sigma_{k,1}(X)  \epsilon^2 + \mathcal{O}(\epsilon^3).
% \end{align}
% The expression above is negative if $\psi>1$, which implies that $\sigma_{R,2}(X) = -\sigma_{y,2}(X)$ is positive. In this case, a negative aggregate shock redistribute wealth to more risk averse investors, leading to a rise in the risk premium and a decline in the risk-free rate. If $\psi>1$, the risk premium effect dominates, so the price-dividend ratio, $ 1/y $, falls in response to the shock. The movement in the price-dividend ratio amplifies the initial effect of the drop in dividends. 

% The drift of $\xi_{j,t}$ is given by
% \begin{align}
%     \mu_{\xi_j}(X,\epsilon) &= \frac{\xi_{j,X}(X,\epsilon)}{\xi_j(X,\epsilon)} \mu_{X}(X,\epsilon)+ \frac{1}{2}  \sigma_{X}^\top(Z,\epsilon) \frac{\xi_{j,XX}(X,\epsilon)}{\xi_{j}(X,\epsilon)} \sigma_{X}(X,\epsilon)\nonumber\\
%     &=  \frac{\xi_{j,X,1}(X)}{\xi_{j,0}(X)} \mu_{X,1}(X)\epsilon^2+\mathcal{O}(\epsilon^3)\nonumber\\
%     &=  -(\psi-1)\left(1-\psi^{-1}\right)\frac{\gamma \sigma^2}{2 \xi_{j,0}(X)} \sum_{k=1}^J \left( \hat{\gamma}_k - \hat{\gamma}_0  \right) \mu_{X_k,1}(X) \epsilon^2 + \mathcal{O}(\epsilon^3).
% \end{align}

% The drift for $y_t$ is given by
% \begin{align}
%     \mu_{y}(X,\epsilon) &= \frac{y_{X}(X,\epsilon)}{y(X,\epsilon)} \mu_{X}(X,\epsilon)+ \frac{1}{2}  \sigma_{X}^\top(Z,\epsilon) \frac{y_{XX}(X,\epsilon)}{y(X,\epsilon)} \sigma_{X}(X,\epsilon)\nonumber\\
%     &=  \frac{y_{X,1}(X)}{y_{0}(X)} \mu_{X,1}(X)\epsilon^2+\mathcal{O}(\epsilon^3)\nonumber\\
%     &= \left(1-\psi^{-1}\right)\frac{\gamma \sigma^2}{2 y_{0}(X)} \sum_{k=1}^J \left( \hat{\gamma}_k - \hat{\gamma}_0  \right) \mu_{X_k,1}(X) \epsilon^2 + \mathcal{O}(\epsilon^3).
% \end{align}

% \noindent\textbf{Step 2: Risk exposures of investors.}\\
% We focus on the \textit{inner region}, that is, the case where all investors are sufficiently far from the constraint boundary (on either side). For a constrained investor, the second-order term is zero. For an unconstrained investor, the second-order term is given by
% \begin{align}
%     \sigma_{j,2}(X) &= \frac{\theta_2(X)}{\gamma} - \frac{\theta_1(X)}{\gamma}\hat{\gamma}_j + \frac{\theta_0(X)}{\gamma}\hat{\gamma}_j^2 + \frac{1-\gamma^{-1}}{\psi-1} \sigma_{\xi_j,2}(X),
% \end{align}
% where investor $j$ is unconstrained if the following condition holds:
% \begin{equation}
%      \sum_{k\in\mathcal{J}^u} \frac{x_k}{x_u} \hat{\gamma}_k -\left(\frac{\hat{\alpha}_p  x_0+\frac{\hat{\sigma}}{\|\sigma\|} x_c }{1-x_0-x_c}\right)- \hat{\gamma}_j  < \frac{\hat{\sigma}}{\|\sigma\|}.
% \end{equation}

% \noindent\textbf{Step 3: Aggregate risk aversion.}\\
% The aggregate risk aversion can be written as
% \begin{equation}
%     \gamma_u(X,\epsilon) = \gamma \left[ 1 + \mathbb{E}^u\left[\hat{\gamma}_j\right] \epsilon - \var^u\left[\hat{\gamma}_j\right] \epsilon^2 \right] + \mathcal{O}(\epsilon^3),
% \end{equation}
% where 
% \[\mathbb{E}^u\left[ \hat{\gamma}_j\right] \equiv\sum_{j\in\mathcal{J}^u} \frac{x_j}{x_u} \hat{\gamma}_j,\quad\text{and}\quad \var^u[\hat{\gamma}_j] \equiv \sum_{j\in\mathcal{J}^u} \frac{x_j}{x_u} \hat{\gamma}_j^2 - \left(\sum_{j\in\mathcal{J}^u} \frac{x_j}{x_u} \hat{\gamma}_j\right)^2.\]

% \noindent\textbf{Step 4: Market price of risk.}\\
% The market price of risk is given by
% \begin{small}
% 	\begin{equation}
% 	\theta(X,\epsilon) =  \frac{\gamma_u(X,\epsilon)}{x_{u}} \left[ (1-(1+\hat{\alpha}_p \epsilon) x_{0,t})(\sigma - \sigma_{y}(X,\epsilon)) - (\sigma + \hat{\sigma}\epsilon) x_{c,t}+ \sum_{j\in \mathcal{J}_t^u} x_{j,t} \frac{1-\gamma_j^{-1}}{1-\psi} \sigma_{c_j}(X,\epsilon)\right],
% 	\end{equation}	
% \end{small}

% The second-order term is then given by
% \begin{align}
%     \theta_2(X) &=  \frac{\gamma_{u,0}(X)}{x_{u}} \left[ -(1-x_{0,t})\sigma_{y,2}(X) + \sum_{j\in \mathcal{J}_t^u} x_{j,t} \frac{1-\gamma_j^{-1}}{1-\psi} \sigma_{\xi_j,2}(X)\right] + \nonumber\\
%     &+ \frac{\gamma_{u,1}(X)}{x_{u}} \left[ - \hat{\alpha}_p x_{0} \sigma -\hat{\sigma} x_{c}\right]+ \frac{\gamma_{u,2}(X)}{x_{u}} \left[ (1-x_{0})\sigma - \sigma x_{c} \right].
% \end{align}

% The expression above can be written as
% \begin{align}
%     \theta_2(X) =  - \frac{(\gamma-1) x_c + 1-x_0}{1-x_0-x_c} \sigma_{y,2}(X) - \gamma \sigma \mathbb{E}^u[\hat{\gamma}_j] \frac{ \hat{\alpha}_p x_0+ \frac{\hat{\sigma}}{\sigma}x_c }{1-x_0-x_c} - \gamma \sigma \var^u[\hat{\gamma}_j].
% \end{align}

% \noindent\textbf{Step 5: Interest rate.}\\
% The interest rate is given by
% \begin{align}
%      r(X,\epsilon)   &  =  \rho - \theta(X,\epsilon) \sigma_{R}(X,\epsilon) +\psi^{-1} (\mu - \mu_{y}(X,\epsilon)+\sigma_{R,t}^2(X,\epsilon)- \sigma \sigma_{R}(X,\epsilon))\nonumber\\
%      &+ \left(1-\psi^{-1}\right)  \sum_{j=0}^J x_{j}\frac{\gamma_j}{2} \sigma_{j}^2(X,\epsilon)  \nonumber\\
%     &+\psi^{-1}\sum_{j=0}^J x_{j,t} \left[\mu_{\xi_j}(X,\epsilon) + (1-\gamma_j)  \sigma_{\xi_j}(X,\epsilon) \sigma_{j}(X,\epsilon) + \frac{\psi - \gamma_j}{1-\psi} \frac{\sigma_{\xi_j,t}^2(X,\epsilon)}{2}\right].
% \end{align}

% The second-order term is given by
% \begin{align}
%     r_2(X) &= - \theta_2(X) \sigma + \theta_0(X) \sigma_{y,2}(X)-\psi^{-1} (\mu_{y,2}(X) + \sigma \sigma_{y,2}(X))\nonumber\\
%     &+ \left(1-\psi^{-1}\right) \sum_{j=0}^J x_{j} \gamma\sigma \left[ \sigma_{j,2}(X) + \frac{\sigma_{j,1}^2(X)}{2 \sigma} + \hat{\gamma}_j \sigma_{j,1} \right]\nonumber\\
%     &+\psi^{-1} \sum_{j=0}^J x_{j}\left[ \mu_{\xi_j,2}(X) + (1-\gamma) \sigma_{\xi_j,2}(X) \sigma \right].
% \end{align}

% \noindent\textbf{Step 6: Dividend yield.}\\
% The dividend yield, $y$, is given by
% \begin{align}
%     y(X,\epsilon) &= \psi \rho + (1-\psi) \left[ r(X,\epsilon) + \theta(X,\epsilon) \sigma_{R}(X,\epsilon) - \sum_{j=0}^J x_{j}\frac{\gamma_j}{2} \sigma_{j}^2(X,\epsilon) \right] \nonumber\\
%     &\qquad +\sum_{j=0}^J x_{j} \left[\mu_{\xi_j}(X,\epsilon) + (1-\gamma_j)  \sigma_{\xi_j}(X,\epsilon) \sigma_{j}(X,\epsilon) + \frac{\psi - \gamma_j}{1-\psi} \frac{\sigma_{\xi_j}^2(X,\epsilon)}{2}\right].
% \end{align}

% The second-order term in the expansion of $y(X,\epsilon)$ is given by 
% \begin{small}
% 	\begin{align}
% 	y_2(X) &=  (1-\psi) \left[ r_2(X) + \theta_2(X) \sigma + \theta_0(X) \sigma_{R,2}(X) - \sum_{j=0}^J x_{j}\frac{\gamma}{2} \left( \sigma_{j,1}^2(X) + 2 \sigma \sigma_{j,2}(X) + 2 \hat{\gamma}_j \sigma \sigma_{j,1}(X)  \right)\right] \nonumber\\
% 	&\qquad +\sum_{j=0}^J x_{j} \left[\mu_{\xi_j,2}(X) + (1-\gamma)  \sigma_{\xi_j,2}(X) \sigma\right].
% 	\end{align}
% \end{small}


% Using the expression for the interest rate, we obtain 
% \begin{align}
%     y_2(X) &=  \left(1-\psi^{-1}\right) \left[ \mu_{y,2}(X) + \sigma \sigma_{y,2}(X) + \sum_{j=0}^J x_{j}\gamma \sigma \left( \frac{\sigma_{j,1}^2(X)}{2 \sigma} + \sigma_{j,2}(X) + \hat{\gamma}_j \sigma_{j,1}(X)  \right)\right] \nonumber\\
%     &\qquad +\psi^{-1}\sum_{j=0}^J x_{j} \left[\mu_{\xi_j,2}(X) + (1-\gamma)  \sigma_{\xi_j,2}(X) \sigma\right].
% \end{align}

% Given that $\mu_{\xi_j,2} = (1-\psi) \mu_{y,2} $ and $\sigma_{\xi_j,2} = (1-\psi) \sigma_{y,2} $, we obtain
% \begin{align}
%     y_2(X) &=  \left(1-\psi^{-1}\right) \left[\gamma \sigma_{y,2}(X) \sigma+\sum_{j=0}^J x_{j}\gamma \sigma \left( \frac{\sigma_{j,1}^2(X)}{2 \sigma} + \sigma_{j,2}(X) + \hat{\gamma}_j \sigma_{j,1}(X)  \right)\right] \nonumber\\
%      &=\left(1-\psi^{-1}\right) \sum_{j=0}^J x_{j}\gamma \sigma \left( \frac{\sigma_{j,1}^2(X)}{2 \sigma}  + \hat{\gamma}_j \sigma_{j,1}(X)  \right),    
% \end{align}
%  where in the second equality, we use the fact that, from the market clearing condition for the risky asset, we have $\sum_{j=0}^J x_j\sigma_{j,2}(X)=\sigma_{R,2}(X)=-\sigma_{y,2}(X).$

% \noindent\textbf{Step 7: Risk premium.}\\
% Since the risk premium is given by $ \pi_t=\theta_t\sigma_{R,t} $, we can write
% \begin{equation*}
% \pi(X,\epsilon) = \theta_0(X)\sigma_{R,0}(X) + \theta_1(X)\sigma_{R,0}(X)\epsilon+
% \left(\theta_{0}(X)\sigma_{R,2}(X)+\theta_2(X)\sigma_{R,0}(X)\right)\epsilon^2+\mathcal{O}\left(\epsilon^3\right).
% \end{equation*}
% Given that $ \sigma_{R,0}=\sigma $ and $ \sigma_{R,2}(X)=-\sigma_{y,2}(X) $, we get
% \begin{equation}\label{eq:rp-second-order}
% 	\pi_2(X) = -\theta_{0}(X)\sigma_{y, 2}(X)+\theta_2(X)\sigma.
% \end{equation}

% \end{proof}

\section{Derivation of the Market Elasticity}

Let $p(X,\epsilon) \equiv 1/y(X,\epsilon)$ denote the price-dividend ratio.
The second-order expansion of $p(X,\epsilon)$ is given by
\begin{equation}\label{eq:price-dividend}
    p(X,\epsilon)
    =
    \frac{1}{y_0(X)}
    -
    \frac{y_1(X)}{y_0^2(X)} \epsilon
    +
    \left[
    \frac{y_1^2(X)}{y_0^3(X)}
    -
    \frac{y_{2}(X)}{y_0^2(X)}
    \right] \epsilon^2
    + \mathcal{O}(\epsilon^3).
\end{equation}

Let
\begin{equation}
    F(X,\epsilon)
    \equiv
    \omega_0 \frac{W_0 (1+\hat{\alpha}_p \epsilon) - W_0}{P}
    =
    \hat{\alpha}_p \epsilon x_0
\end{equation}
denote the flow into the risky asset relative to the benchmark economy.
Note that $F = \mathcal{O}(\epsilon)$.

In this section, we consider the case where the portfolio of passive investors is constant. 

\begin{proof}

Using \eqref{eq:price-dividend}, the derivative of the price-dividend ratio with respect to flows can be written as
\begin{equation}\label{eq:elasticity_first-order}
\frac{\partial p(X,\epsilon)}{\partial F(X,\epsilon)}
=
\frac{1}{\partial F / \partial \hat{\alpha}_p}
\frac{\partial p(X,\epsilon)}{\partial \hat{\alpha}_p}
=
-\frac{1}{y_0^2(X)}\frac{\partial y_1(X)}{\partial \hat{\alpha}_p}\frac{1}{x_0}
+ \mathcal{O}(\epsilon).
\end{equation}

From Proposition~\ref{prop:first_order}, $y_1(X)$ does not depend on $\hat{\alpha}_p$, so
\[
\frac{\partial y_1(X)}{\partial \hat{\alpha}_p} = 0.
\]
Therefore,
\[
\frac{\partial p(X,\epsilon)}{\partial F(X,\epsilon)} = 0 + \mathcal{O}(\epsilon),
\]
which implies
\[
\varepsilon_M^{-1}
=
\frac{1}{p(X,\epsilon)}
\frac{\partial p(X,\epsilon)}{\partial F(X,\epsilon)}
=
0 + \mathcal{O}(\epsilon).
\]

Thus, the aggregate market elasticity is infinite to first order.

To understand this result, note that from the pricing condition
\[
y_t = r_t + \pi_t - \mu_{P,t},
\]
we can write the first-order term as
\[
y_1(X) = r_1(X) + \pi_1(X) - \mu_{P,1}(X).
\]

Since $y_1(X)$ is constant, we must have
\[
\mu_{y,1}(X) = 0, \qquad \sigma_{y,1}(X) = 0,
\]
which implies $\sigma_{R,1}(X)=0$ and therefore $\mu_{P,1}(X)=0$.

Hence,
\[
\frac{\partial y_1(X)}{\partial \hat{\alpha}_p}
=
\frac{\partial r_1(X)}{\partial \hat{\alpha}_p}
+
\frac{\partial \pi_1(X)}{\partial \hat{\alpha}_p}.
\]

From Proposition~\ref{prop:first_order},
\begin{align*}
\frac{\partial r_1(X)}{\partial \hat{\alpha}_p}
&= -\sigma \frac{\partial \eta_1(X)}{\partial \hat{\alpha}_p}
= \frac{\gamma \|\sigma\|^2}{x_a} x_0, \\
\frac{\partial \pi_1(X)}{\partial \hat{\alpha}_p}
&= \sigma \frac{\partial \eta_1(X)}{\partial \hat{\alpha}_p}
= -\frac{\gamma \|\sigma\|^2}{x_a} x_0.
\end{align*}

These two effects exactly cancel.
Thus, at first order, portfolio flows raise the risk premium but lower the interest rate by the same amount, leaving prices unchanged.

\medskip

We now turn to the second-order effect.
From \eqref{eq:price-dividend}, we have
\begin{equation}
\frac{\partial p(X,\epsilon)}{\partial F(X,\epsilon)}
=
-\frac{1}{y_0^2(X)}
\frac{\partial y_2(X)}{\partial \hat{\alpha}_p}
\frac{\epsilon}{x_0}
+ \mathcal{O}(\epsilon^2).
\end{equation}

Because $F = \mathcal{O}(\epsilon)$, second-order changes in the dividend yield generate a first-order response in prices.

At second order, the cancellation between interest rates and risk premia breaks down due to endogenous reallocation of risk across investors.
In particular, from the expression for $y_2(X)$,
\begin{equation}
    y_{2}(X)
    =
    (1-\psi^{-1})\gamma \|\sigma\|^2
    \sum_{j=0}^2 x_j
    \left(
        \frac{1}{2} m_{j,1}^2
        +
        \hat{\gamma}_j m_{j,1}
    \right),
\end{equation}
where $\sigma_{j,1} = m_{j,1}\sigma$.

Thus, $\partial y_2 / \partial \hat{\alpha}_p \neq 0$, implying a non-zero price impact and a finite aggregate elasticity.

\end{proof}

% \subsection{Proof of Proposition~\ref{prop:elasticity_passive_demand}}\label{app:prop_elasticity_passive_demand}
% \begin{proof}
% 	Consider the case where there is no preference heterogeneity, so $\hat{\gamma}_j = 0$ for $j = 0,1,2$, and there are no leverage constraints, $\nu_{2,t} = 0$.
% 	In this case, $y_2(X)$ simplifies to
% 	\begin{align}
% 	y_2(X) &=  \left(1-\psi^{-1}\right) \gamma \|\sigma\|^2 \sum_{j=0}^2 x_{j} \frac{m_{j,1}^2(X)}{2},
% 	\end{align}
%     where
%     \begin{equation}
%         m_{0,1} = \hat{\alpha}_p, \qquad \qquad 
%         m_{j,1} =   -\frac{x_0}{x_a}\hat{\alpha}_{p}, \qquad \text{for $j = 1,2$}.
%     \end{equation}
	

% 	Using the expression for $ \sigma_{j,1} $ in Proposition~\ref{prop:first_order}, the (inverse) aggregate market elasticity is given by
% 	\begin{equation}
% 	\frac{1}{p(X,\epsilon)}\frac{\partial p(X,\epsilon)}{\partial F(X,\epsilon)} = - \frac{1-\psi^{-1}}{2 y_0(X)} \gamma \|\sigma\|^2 \left(2x_0 \hat{\alpha}_p  + x_{a}  \frac{2 \hat{\alpha}_p x_0^2}{x_{a}^2}\right)  \frac{\epsilon}{x_0} + \mathcal{O}(\epsilon^2), 
% 	\end{equation}
% 	where $ x_{a} \equiv 1-x_{0} $ denotes the wealth share of active investors.	This can be written as
% 	\begin{equation}
% 	\frac{1}{p(X,\epsilon)}\frac{\partial p(X,\epsilon)}{\partial F(X,\epsilon)} =  \left(1-\psi^{-1}\right)\frac{\gamma \|\sigma\|^2}{y_0(X)}  \frac{1-\overline{\alpha}_p}{x_{a}} + \mathcal{O}(\epsilon^2), 
% 	\end{equation}
% 	where we use $ \overline{\alpha}_p = 1+\hat{\alpha}_p \epsilon $ from Equation~\eqref{eq:apbar_index}.
% 	%\textcolor{red}{Notice that the aggregate elasticity in this case is decreasing in $ \overline{\alpha}_p $. The price-dividend ratio drops as passive demand exceeds one.}
% \end{proof}

% The impact of flows of the risk premium can be written as
\subsection{Proof of Proposition~\ref{prop:elasticity_passive_demand}}\label{app:prop_elasticity_passive_demand}
\begin{proof}
Consider the case with no preference heterogeneity and no leverage constraints.
That is, $\hat{\gamma}_j=0$ for $j=0,1,2$ and $\nu_{2,t}=0$.
In this case, the first-order exposure coefficients are
\begin{equation}
    m_{0,1}=\hat{\alpha}_p, \qquad
    m_{j,1}=-\frac{x_0}{x_a}\hat{\alpha}_p, \qquad j=1,2,
\end{equation}
where $x_a\equiv 1-x_0$ is the wealth share of active investors and $\sigma_{j,1}=m_{j,1}\sigma$.
The second-order correction to the dividend yield therefore simplifies to
\begin{align}
    y_{Y,2}(X)
    &= (1-\psi^{-1})\gamma\|\sigma\|^2
    \sum_{j=0}^2 x_j \frac{m_{j,1}^2}{2} \nonumber\\
    &= (1-\psi^{-1})\gamma\|\sigma\|^2
    \frac{1}{2}\left[x_0\hat{\alpha}_p^2+x_a\left(\frac{x_0}{x_a}\hat{\alpha}_p\right)^2\right] \nonumber\\
    &= (1-\psi^{-1})\gamma\|\sigma\|^2
    \frac{x_0}{2x_a}\hat{\alpha}_p^2.
\end{align}
Hence,
\begin{equation}
    \frac{\partial y_{Y,2}(X)}{\partial \hat{\alpha}_p}
    = (1-\psi^{-1})\gamma\|\sigma\|^2
    \frac{x_0}{x_a}\hat{\alpha}_p.
\end{equation}

Using the change of variables from the previous subsection,
\begin{equation}
    \frac{1}{p(X,\epsilon)}\frac{\partial p(X,\epsilon)}{\partial F(X,\epsilon)}
    = -\frac{1}{y_0(X)}
    \frac{\partial y_{Y,2}(X)}{\partial \hat{\alpha}_p}
    \frac{\epsilon}{x_0}
    +\mathcal{O}(\epsilon^2).
\end{equation}
Substituting the derivative above gives
\begin{equation}
    \frac{1}{p(X,\epsilon)}\frac{\partial p(X,\epsilon)}{\partial F(X,\epsilon)}
    = -(1-\psi^{-1})\frac{\gamma\|\sigma\|^2}{y_0(X)}
    \frac{\hat{\alpha}_p\epsilon}{x_a}
    +\mathcal{O}(\epsilon^2).
\end{equation}
Finally, since $\overline{\alpha}_p=1+\hat{\alpha}_p\epsilon$, we obtain
\begin{equation}
    \frac{1}{p(X,\epsilon)}\frac{\partial p(X,\epsilon)}{\partial F(X,\epsilon)}
    = (1-\psi^{-1})\frac{\gamma\|\sigma\|^2}{y_0(X)}
    \frac{1-\overline{\alpha}_p}{x_a}
    +\mathcal{O}(\epsilon^2).
\end{equation}
This proves the result.
\end{proof}


% \[\frac{\partial \pi(X,\epsilon)}{\partial F(X,\epsilon)}=\sigma\frac{\partial \theta_1(X)}{\partial\hat{\alpha}_p}\frac{\epsilon}{x_0}+
% \frac{\partial \pi_2(X)}{\partial\hat{\alpha}_p}\frac{\epsilon}{x_0}\]

\subsection{Proof of Proposition~\ref{prop_elasticity_pref_heterog}}\label{app:prop_elasticity_pref_heterog}
\begin{proof}
Consider the case in which investors have heterogeneous risk aversions but the leverage constraints are slack.
Then $\hat{\nu}_{j,1}=0$ for $j=1,2$.
Let
\begin{equation}
    \bar{\gamma}_a \equiv \sum_{j=1}^2 \frac{x_j}{x_a}\hat{\gamma}_j,
    \qquad x_a\equiv x_1+x_2=1-x_0.
\end{equation}
From the first-order solution, the scalar exposure coefficients defined by $\sigma_{j,1}=m_{j,1}\sigma$ are
\begin{equation}
    m_{0,1}=\hat{\alpha}_p,
    \qquad
    m_{j,1}=\bar{\gamma}_a-\frac{x_0}{x_a}\hat{\alpha}_p-\hat{\gamma}_j,
    \qquad j=1,2.
\end{equation}
The second-order correction to the dividend yield is
\begin{equation}
    y_{Y,2}
    = (1-\psi^{-1})\gamma\|\sigma\|^2
    \sum_{j=0}^2 x_j\left(\frac{1}{2}m_{j,1}^2+\hat{\gamma}_j m_{j,1}\right).
\end{equation}
We differentiate this expression with respect to $\hat{\alpha}_p$.
Since
\begin{equation}
    \frac{\partial m_{0,1}}{\partial \hat{\alpha}_p}=1,
    \qquad
    \frac{\partial m_{j,1}}{\partial \hat{\alpha}_p}=-\frac{x_0}{x_a},
    \qquad j=1,2,
\end{equation}
we obtain
\begin{align}
    \frac{\partial y_{Y,2}}{\partial \hat{\alpha}_p}
    &= (1-\psi^{-1})\gamma\|\sigma\|^2
    \sum_{j=0}^2 x_j(m_{j,1}+\hat{\gamma}_j)
    \frac{\partial m_{j,1}}{\partial \hat{\alpha}_p} \nonumber\\
    &= (1-\psi^{-1})\gamma\|\sigma\|^2 x_0
    \left[\hat{\alpha}_p+\hat{\gamma}_0
    -\sum_{j=1}^2 \frac{x_j}{x_a}(m_{j,1}+\hat{\gamma}_j)\right].
\end{align}
For active investors,
\begin{equation}
    m_{j,1}+\hat{\gamma}_j
    = \bar{\gamma}_a-\frac{x_0}{x_a}\hat{\alpha}_p.
\end{equation}
Therefore,
\begin{equation}
    \frac{\partial y_{Y,2}}{\partial \hat{\alpha}_p}
    = (1-\psi^{-1})\gamma\|\sigma\|^2 x_0
    \left[\frac{\hat{\alpha}_p}{x_a}+\hat{\gamma}_0-\bar{\gamma}_a\right].
\end{equation}

Using the change of variables from the market-elasticity derivation,
\begin{equation}
    \frac{1}{p(X,\epsilon)}\frac{\partial p(X,\epsilon)}{\partial F(X,\epsilon)}
    = -\frac{1}{y_0(X)}
    \frac{\partial y_{Y,2}(X)}{\partial \hat{\alpha}_p}
    \frac{\epsilon}{x_0}
    +\mathcal{O}(\epsilon^2).
\end{equation}
Substituting the derivative above gives
\begin{equation}
    \frac{1}{p(X,\epsilon)}\frac{\partial p(X,\epsilon)}{\partial F(X,\epsilon)}
    = -(1-\psi^{-1})\frac{\gamma\|\sigma\|^2}{y_0(X)}
    \left[\frac{\hat{\alpha}_p\epsilon}{x_a}+(\hat{\gamma}_0-\bar{\gamma}_a)\epsilon\right]
    +\mathcal{O}(\epsilon^2).
\end{equation}
Finally, using $\overline{\alpha}_p=1+\hat{\alpha}_p\epsilon$ and $\gamma_j=\gamma(1+\hat{\gamma}_j\epsilon)$, we can write the expression in primitives as
\begin{equation}
    \frac{1}{p(X,\epsilon)}\frac{\partial p(X,\epsilon)}{\partial F(X,\epsilon)}
    = (1-\psi^{-1})\frac{\gamma\|\sigma\|^2}{y_0(X)}
    \left[
        \frac{1-\overline{\alpha}_p}{x_a}
        -\frac{\gamma_0-\mathbb{E}_a[\gamma_j]}{\gamma}
    \right]
    +\mathcal{O}(\epsilon^2),
\end{equation}
where
\begin{equation}
    \mathbb{E}_a[\gamma_j]\equiv \sum_{j=1}^2 \frac{x_j}{x_a}\gamma_j.
\end{equation}
Equivalently, this can be written in the form of Proposition~\ref{prop_elasticity_pref_heterog} by setting $x_c=0$, $x_u=x_a$, and
\begin{equation}
    \overline{\alpha}_{p}^{\mathrm{opt}}
    =1-x_a\frac{\gamma_0-\mathbb{E}_a[\gamma_j]}{\gamma} = 1 + \left( \sum_{j=0}^2 x_k \hat{\gamma}_k  - \hat{\gamma}_0 \right) \epsilon,
\end{equation}
where the expression above coincides with first-order correction for the passive investor's optimal portfolio share, that is, the portfolio they would choose if they could freely adjust their portfolio share.

Then
\begin{equation}
    \frac{1}{p(X,\epsilon)}\frac{\partial p(X,\epsilon)}{\partial F(X,\epsilon)}
    = (1-\psi^{-1})\frac{\gamma\|\sigma\|^2}{y_0(X)}
    \frac{\overline{\alpha}_{p}^{\mathrm{opt}}-\overline{\alpha}_p}{x_a}
    +\mathcal{O}(\epsilon^2).
\end{equation}
This proves the version of the result without leverage constraints.

\medskip

We next derive the corresponding expression when investor $2$ is constrained and investor $1$ is unconstrained.
In this case, the first-order exposure coefficients are
\begin{equation}
    m_{0,1}=\hat{\alpha}_p, \qquad
    m_{2,1}=\hat{\sigma}, \qquad
    m_{1,1}=-\frac{x_0\hat{\alpha}_p+x_2\hat{\sigma}}{x_1},
\end{equation}
where the expression for $m_{1,1}$ follows from market clearing,
\begin{equation}
    x_0m_{0,1}+x_1m_{1,1}+x_2m_{2,1}=0.
\end{equation}
Using again
\begin{equation}
    y_{Y,2}
    = (1-\psi^{-1})\gamma\|\sigma\|^2
    \sum_{j=0}^2 x_j\left(\frac{1}{2}m_{j,1}^2+\hat{\gamma}_j m_{j,1}\right),
\end{equation}
we differentiate with respect to $\hat{\alpha}_p$.
Since
\begin{equation}
    \frac{\partial m_{0,1}}{\partial \hat{\alpha}_p}=1, \qquad
    \frac{\partial m_{1,1}}{\partial \hat{\alpha}_p}=-\frac{x_0}{x_1}, \qquad
    \frac{\partial m_{2,1}}{\partial \hat{\alpha}_p}=0,
\end{equation}
we obtain
\begin{align}
    \frac{\partial y_{Y,2}}{\partial \hat{\alpha}_p}
    &= (1-\psi^{-1})\gamma\|\sigma\|^2
    \left[x_0(m_{0,1}+\hat{\gamma}_0)
    -x_0(m_{1,1}+\hat{\gamma}_1)\right] \nonumber\\
    &= (1-\psi^{-1})\gamma\|\sigma\|^2 x_0
    \left[
        \frac{x_0+x_1}{x_1}\hat{\alpha}_p
        +\frac{x_2}{x_1}\hat{\sigma}
        +\hat{\gamma}_0-\hat{\gamma}_1
    \right].
\end{align}
Using the change of variables from the market-elasticity derivation,
\begin{equation}
    \frac{1}{p(X,\epsilon)}\frac{\partial p(X,\epsilon)}{\partial F(X,\epsilon)}
    = -\frac{1}{y_0(X)}
    \frac{\partial y_{Y,2}(X)}{\partial \hat{\alpha}_p}
    \frac{\epsilon}{x_0}
    +\mathcal{O}(\epsilon^2),
\end{equation}
we get
\begin{align}
    \frac{1}{p(X,\epsilon)}\frac{\partial p(X,\epsilon)}{\partial F(X,\epsilon)}
    &= -(1-\psi^{-1})\frac{\gamma\|\sigma\|^2}{y_0(X)}
    \left[
        \frac{x_0+x_1}{x_1}\hat{\alpha}_p\epsilon
        +\frac{x_2}{x_1}\hat{\sigma}\epsilon
        +(\hat{\gamma}_0-\hat{\gamma}_1)\epsilon
    \right]
    +\mathcal{O}(\epsilon^2).
\end{align}
Finally, using $\overline{\alpha}_p=1+\hat{\alpha}_p\epsilon$, $\overline{\sigma}=\|\sigma\|(1+\hat{\sigma}\epsilon)$, and $\gamma_j=\gamma(1+\hat{\gamma}_j\epsilon)$, this becomes
\begin{align}
    \frac{1}{p(X,\epsilon)}\frac{\partial p(X,\epsilon)}{\partial F(X,\epsilon)}
    &= (1-\psi^{-1})\frac{\gamma\|\sigma\|^2}{y_0(X)}
    \left[
        \frac{x_0+x_1}{x_1}(1-\overline{\alpha}_p)
        -\frac{x_2}{x_1}\left(\frac{\overline{\sigma}}{\|\sigma\|}-1\right)
        -\frac{\gamma_0-\gamma_1}{\gamma}
    \right]
    +\mathcal{O}(\epsilon^2).
\end{align}
Define $x_c\equiv x_2$, $x_u\equiv x_1=x_a-x_c$, and $\mathbb{E}^u[\gamma_j]\equiv\gamma_1$.
Also define
\begin{equation}
    \overline{\alpha}_{p}^{\mathrm{opt}}
    = 1
    - \frac{x_u}{1-x_c}\frac{\gamma_0-\mathbb{E}^{u}[\gamma_j]}{\gamma}
    - \frac{x_c}{1-x_c}\left(\frac{\overline{\sigma}}{\|\sigma\|}-1\right) = 1 + \left( \sum_{j=0}^1 \frac{x_k}{1-x_c} \hat{\gamma}_k  - \hat{\gamma}_0 - \frac{x_c}{1-x_c} \hat{\sigma} \right) \epsilon,
\end{equation}
where the expression above coincides with first-order correction for the passive investor's optimal portfolio share when the leverage constraint of investor $2$ binds.

Since $1-x_c=x_0+x_1$, the previous display is equivalent to
\begin{equation}
    \frac{1}{p(X,\epsilon)}\frac{\partial p(X,\epsilon)}{\partial F(X,\epsilon)}
    = (1-\psi^{-1})\frac{\gamma\|\sigma\|^2}{y_0(X)}
    \frac{\overline{\alpha}_{p}^{\mathrm{opt}}-\overline{\alpha}_{p}}{x_a-x_c}
    (1-x_c)
    +\mathcal{O}(\epsilon^2).
\end{equation}
This proves the version of the result in the regime in which the leverage constraint of investor $2$ binds.


\end{proof}
% \begin{proof}
% 	Consider the case where there is no preference heterogeneity and active investors do not face leverage constraints.
% 	In this case, $y_2(X)$ simplifies to
% 	\begin{align}
% 	y_2(X) &=  \left(1-\psi^{-1}\right) \gamma \sigma^2 \sum_{j=0}^J x_{j} \frac{\sigma_{j,1}^2(X)}{2 \sigma^2},
% 	\end{align}
% 	using the fact that $\sigma_{y,2} = \sigma_{j,2} = 0$ when $\hat{\gamma}_j = 0$ for $j = 0,1,\ldots,J$ when investors have the same preferences.
	
% 	Using the expression for $ \sigma_{j,1} $ in Proposition~\ref{prop:first_order}, the (inverse) aggregate market elasticity is given by
% 	\begin{equation}
% 	\frac{1}{p(X,\epsilon)}\frac{\partial p(X,\epsilon)}{\partial F(X,\epsilon)} = - \frac{1-\psi^{-1}}{2 y_0(X)} \gamma \sigma^2 \left(2x_0 \hat{\alpha}_p  + x_{a}  \frac{2 \hat{\alpha}_p x_0^2}{x_{a}^2}\right)  \frac{\epsilon}{x_0} + \mathcal{O}(\epsilon^2), 
% 	\end{equation}
% 	where $ x_{a} \equiv 1-x_{0} $ denotes the wealth share of active investors.	This can be written as
% 	\begin{equation}
% 	\frac{1}{p(X,\epsilon)}\frac{\partial p(X,\epsilon)}{\partial F(X,\epsilon)} =  \left(1-\psi^{-1}\right)\frac{\gamma \sigma^2}{y_0(X)}  \frac{1-\overline{\alpha}_p}{x_{a}} + \mathcal{O}(\epsilon^2), 
% 	\end{equation}
% 	where we use $ \overline{\alpha}_p = 1+\hat{\alpha}_p \epsilon $ from Equation~\eqref{eq:apbar_index}.
% 	%\textcolor{red}{Notice that the aggregate elasticity in this case is decreasing in $ \overline{\alpha}_p $. The price-dividend ratio drops as passive demand exceeds one.}
% \end{proof}
% From Equation~\eqref{eq:rp-second-order}, the impact of flows of the risk premium can be written as
% \[\frac{\partial \pi(X,\epsilon)}{\partial F(X,\epsilon)}=\sigma\frac{\partial \theta_1(X)}{\partial\hat{\alpha}_p}\frac{\epsilon}{x_0}+
% \frac{\partial \pi_2(X)}{\partial\hat{\alpha}_p}\frac{\epsilon}{x_0}\]

% \subsection{Proof of Proposition~\ref{prop_elasticity_pref_heterog}}\label{app:prop_elasticity_pref_heterog}
% \begin{proof}
% 	Consider the case in which active investors have heterogeneous risk aversions, but face no leverage constraints.
% 	In this case, the expression for $y_{2}(X)$ in Equation~\eqref{eq:y-second-order} can be written as
% 	\begin{small}
% 	\begin{align}
%             y_2(X) &=  \left(1-\psi^{-1}\right) \gamma \sigma^2 \left[  \sum_{j=1}^J x_{j}\left( \frac{\sigma_{j,1}^2(X)}{2 \sigma^2} +\hat{\gamma}_j\frac{ \sigma_{j,1}(X)}{\sigma}  \right)\right] + \left(1-\psi^{-1}\right) \gamma \sigma^2 x_0 \left(\frac{\hat{\alpha}_p^2}{2} + \hat{\gamma}_0 \hat{\alpha}_p  \right)
% 		% y_2(X) &=  \left(1-\psi^{-1}\right) \gamma \sigma^2 \left[  \sum_{j=1}^J x_{j}\left( \frac{\sigma_{j,1}^2(X)}{2 \sigma^2} - \mathbb{E}^u[\hat{\gamma}_j] \frac{ \hat{\alpha}_p x_0}{1-x_0} - \var^u[\hat{\gamma}_j]-\frac{\theta_1(X)}{\gamma\sigma}\hat{\gamma}_j + \hat{\gamma}_j^2+ \hat{\gamma}_j\frac{ \sigma_{j,1}(X)}{\sigma}  \right)\right] \nonumber\\
% 		% &+\left(1-\psi^{-1}\right) \gamma \sigma^2 x_0 \left(\frac{\hat{\alpha}_p^2}{2} + \hat{\gamma}_0 \hat{\alpha}_p  \right)
% 		\end{align}
% 	\end{small}
	
% 	Note that with unconstrained active investors, the effect of endogenous volatility and hedging demand exactly cancel out.
% 	We first compute the derivative of the term involving $\sigma_{j,1}^2$:
% 	\begin{align}
% 		\sum_{j=0}^J x_j   \frac{\partial \sigma_{j,1}^2}{\partial \hat{\alpha}_p} &= \sum_{j=0}^J 2x_j  \sigma_{j,1} \frac{\partial \sigma_{j,1}}{\partial \hat{\alpha}_p} \nonumber\\
% 		&= 2 x_0  \sigma^2 \hat{\alpha}_p  + 2 \sigma^2 \sum_{j=1}^J x_j \left( \sum_{k\in\mathcal{J}^u} \frac{x_k}{x_u} \hat{\gamma}_k -\left(\frac{\hat{\alpha}_p  x_0}{1-x_0}\right)- \hat{\gamma}_j\right) \left( -\frac{x_0}{1-x_0}\right) \nonumber \\
% 		&= 2 \sigma^2 \left[ x_0 \hat{\alpha}_p + \frac{\hat{\alpha}_p x_0^2}{1-x_0}+ \left((1-x_0) \sum_{k=1}^J \frac{x_k}{x_u} \hat{\gamma}_k - \sum_{j=1}^J x_j\hat{\gamma}_j\right) \left( -\frac{x_0}{1-x_0}\right)\right] \nonumber \\
% 		&= 2 \sigma^2 \frac{\hat{\alpha}_p x_0}{1-x_0}.
% %		\sum_{j=0}^J x_j \gamma  \frac{\partial \sigma_{j,1}^2}{\partial \hat{\alpha}_p} &= \sum_{j=0}^J x_j \gamma 2 \sigma_{j,1} \frac{\partial \sigma_{j,1}}{\partial \hat{\alpha}_p} \nonumber\\
% %		&= 2 x_0 \gamma \sigma^2 \hat{\alpha}_p  + 2 \gamma \sigma^2 \sum_{j=1}^J x_j \left( \sum_{k\in\mathcal{J}^u} \frac{x_k}{x_u} \hat{\gamma}_k -\left(\frac{\hat{\alpha}_p  x_0}{1-x_0}\right)- \hat{\gamma}_j\right) \left( -\frac{x_0}{1-x_0}\right) \nonumber \\
% %		&= 2\gamma \sigma^2 \left[ x_0 \hat{\alpha}_p + \frac{\hat{\alpha}_p x_0^2}{1-x_0}+ \left((1-x_0) \sum_{k=1}^J \frac{x_k}{x_u} \hat{\gamma}_k - \sum_{j=1}^J x_j\hat{\gamma}_j\right) \left( -\frac{x_0}{1-x_0}\right)\right] \nonumber \\
% %		&= 2\gamma \sigma^2 \frac{\hat{\alpha}_p x_0}{1-x_0}.
% 	\end{align}
	
% 	% The derivatives of the terms involving $ \E^u\left [\hat{\gamma}_j\right ] $, $ \theta_1(X) $, and $ \sigma_{j,1}(X) $ with respect to $ \hat{\alpha}_p $ are given by:
% 	% \begin{align}
% 	% \sum_{j=1}^J x_{j} \left[ -\mathbb{E}^u[\hat{\gamma}_k] \frac{x_0}{1-x_0} + \frac{x_0\hat{\gamma}_j}{1-x_0}      \right]&= -\mathbb{E}^u[\hat{\gamma}_k] x_0 + \sum_{j=1}^J\frac{x_j}{1-x_0} \hat{\gamma}_j x_0= 0\\
% 	% \frac{1}{\sigma}  \sum_{j=0}^J x_{j} \hat{\gamma}_j \frac{\partial \sigma_{j,1}(X)}{\partial \hat{\alpha}_p} &=  \sum_{j=1}^J x_j \hat{\gamma}_j \left( -\frac{x_0}{1-x_0} \right)+x_0 \hat{\gamma}_0 = x_0 \left( \hat{\gamma}_0 - \mathbb{E}^u[\hat{\gamma}_k] \right)
% 	% \end{align}
% The derivatives of the term involving $ \sigma_{j,1}(X) $ with respect to $ \hat{\alpha}_p $ are given by:
% \begin{align}
%     % \sum_{j=1}^J x_{j} \left[ -\mathbb{E}^u[\hat{\gamma}_k] \frac{x_0}{1-x_0} + \frac{x_0\hat{\gamma}_j}{1-x_0}      \right]&= -\mathbb{E}^u[\hat{\gamma}_k] x_0 + \sum_{j=1}^J\frac{x_j}{1-x_0} \hat{\gamma}_j x_0= 0\\
%     \frac{1}{\sigma}  \sum_{j=0}^J x_{j} \hat{\gamma}_j \frac{\partial \sigma_{j,1}(X)}{\partial \hat{\alpha}_p} &=  \sum_{j=1}^J x_j \hat{\gamma}_j \left( -\frac{x_0}{1-x_0} \right)+x_0 \hat{\gamma}_0 = x_0 \left( \hat{\gamma}_0 - \mathbb{E}^u[\hat{\gamma}_k] \right)
% \end{align}
 
% 	Thus, the (inverse) aggregate elasticity is then given by:
% 	\begin{align}
% 	\frac{1}{p(X,\epsilon)}\frac{\partial p(X,\epsilon)}{\partial F(X,\epsilon)} =  \left(1-\psi^{-1}\right)\frac{\gamma \sigma^2}{y_0(X)} \left[ \frac{1-\overline{\alpha}_p}{x_a} - \frac{\gamma_0-\mathbb{E}^u\left[\gamma_j \right]}{\gamma}\right]+ \mathcal{O}(\epsilon^2),
% 	\end{align}
% 	where $x_a\equiv1-x_0$ is the wealth share of active investors, and we use $ \gamma_{j}=\gamma\left(1+\hat{\gamma}_j\epsilon\right) $ from Equation~\eqref{eq:gamma-perturb}.
% \end{proof}

\begin{comment}
\subsection{Proof of Proposition~\ref{prop:elasticity_leverage_constraint}}\label{app:prop_elasticity_leverage_constraint}
% \begin{proof}
% 	Consider the case in which active investors have heterogeneous risk aversions and also face leverage constraints. We can then write $y_2(X)$ as follows
% 	\begin{align}
% 	y_2(X) &=  \left(1-\psi^{-1}\right) \gamma \sigma^2 \sum_{j=0}^J x_{j} \left[\frac{\sigma_{y,2}(X)}{\sigma}+ \frac{\sigma_{j,1}^2(X)}{2 \sigma^2} +\frac{ \sigma_{j,2}(X)}{\sigma} + \hat{\gamma}_j \frac{\sigma_{j,1}(X)}{\sigma}  \right] \nonumber\\
% 	&=\left(1-\psi^{-1}\right) \gamma \sigma^2 \sum_{j=1}^J x_{j} \left[\frac{\sigma_{j,1}^2(X)}{2 \sigma^2}-\frac{ x_c}{x_u}\frac{\sigma_{y,2}(X)}{\sigma}-\mathbb{E}^u[\hat{\gamma}_k] \frac{ \hat{\alpha}_p x_0+ \frac{\hat{\sigma}}{\sigma}x_c }{1-x_0-x_c} -\var^u[\hat{\gamma}_k]  \right.\nonumber\\
% 	&-\left. \hat{\gamma}_j \left(\sum_{k\in\mathcal{J}^u} \frac{x_k}{x_a} \hat{\gamma}_k -\frac{\hat{\alpha}_p  x_0 + \frac{\hat{\sigma}}{\sigma} x_c}{1-x_0-x_c}-\hat{\gamma}_j\right)+\hat{\gamma}_j \min\left\{ \sum_{k\in\mathcal{J}^u} \frac{x_k}{x_u} \hat{\gamma}_k -\frac{\hat{\alpha}_p  x_0+\frac{\hat{\sigma}}{\sigma} x_c }{1-x_0-x_c}- \hat{\gamma}_j  ,\frac{\hat{\sigma}}{\sigma}  \right\}  \right] \nonumber \\
% 	&+\left(1-\psi^{-1}\right) \gamma \sigma^2 x_0 \left[ \frac{\hat{\alpha}_p^2}{2} + \hat{\gamma}_0 \hat{\alpha}_p \right],
% 	\end{align}
% 	using the fact that $\sigma_{0,2} = -\sigma_{y,2}$ in the last equality.
% 	The derivative of $y_2(X)$ with respect to $\hat{\alpha}_p$ is given by 
% 	\begin{align}
% 	\frac{\partial y_2(X)}{\partial \hat{\alpha}_p}   &=\left(1-\psi^{-1}\right) \gamma \sigma^2 \left(\sum_{j=1}^J x_{j} \left[\frac{\sigma_{j,1}(X)}{ \sigma^2} \frac{\partial \sigma_{j,1}}{\partial \hat{\alpha}_p}-\frac{ x_c}{x_u \sigma}\frac{\partial \sigma_{y,2}(X)}{\partial \hat{\alpha}_p}\right]-\mathbb{E}^u[\hat{\gamma}_k] \frac{  x_0 (1-x_0)}{1-x_0-x_c}\right. \nonumber\\
% 	&\left. +\mathbb{E}^c[\hat{\gamma}_k] \frac{  x_0 x_c }{1-x_0-x_c} + x_0 (\hat{\alpha}_p+ \hat{\gamma}_0) \right), \nonumber 
% 	\end{align}
% 	where
% 	\begin{align}
% 	\frac{1}{\sigma^2}\sum_{j=0}^J x_j \frac{\partial \sigma_{j,1}^2}{\partial \hat{\alpha}_p} &=\frac{1}{\sigma^2} \sum_{j=0}^J x_j2 \sigma_{j,1} \frac{\partial \sigma_{j,1}}{\partial \hat{\alpha}_p} \nonumber\\
% 	&= 2 x_0  \hat{\alpha}_p  + 2 \sum_{j\in\mathcal{J}^u} x_j \left( \sum_{k\in\mathcal{J}^u} \frac{x_k}{x_u} \hat{\gamma}_k -\left(\frac{\hat{\alpha}_p  x_0+\frac{\hat{\sigma}}{\sigma}x_c}{1-x_0-x_c}\right)- \hat{\gamma}_j\right) \left( -\frac{x_0}{1-x_0-x_c}\right) \nonumber \\
% 	&= 2 \left[ x_0 \hat{\alpha}_p + \frac{\hat{\alpha}_p x_0+\frac{\hat{\sigma}}{\sigma}x_c}{1-x_0-x_c}x_0+ \left(x_u \sum_{k\in\mathcal{J}^u} \frac{x_k}{x_u} \hat{\gamma}_k - \sum_{j\in\mathcal{J}^u} x_j\hat{\gamma}_j\right) \left( -\frac{x_0}{1-x_0-x_c}\right)\right] \nonumber \\
% 	&= 2 x_0 \left[\frac{\hat{\alpha}_p}{x_u} + \left(\frac{\hat{\sigma}}{\sigma} - \hat{\alpha}_p \right)\frac{x_c}{x_u}\right].
% 	\end{align}
	
% 	The second derivative of $y(X,\epsilon)$ can then be written as
% 	\begin{align}
% 	\frac{\partial y_2(X)}{\partial \hat{\alpha}_p}   &=\left(1-\psi^{-1}\right) \gamma \sigma^2 \left[\frac{\hat{\alpha}_p}{x_u}- (\mathbb{E}^u[\hat{\gamma}_k]-\hat{\gamma}_0) + \left(\frac{\hat{\sigma}}{\sigma} - \hat{\alpha}_p \right)\frac{x_c}{x_u}\nonumber\right.\\ 
% 	&\left. - \frac{x_c}{\sigma x_0} \frac{\partial \sigma_{y,2}(X)}{\partial \hat{\alpha}_p} -\frac{  x_c \left(\mathbb{E}^u[\hat{\gamma}_k]- \mathbb{E}^c[\hat{\gamma}_k]\right) }{1-x_0-x_c}  \right] x_0, \nonumber 
% 	\end{align}
% 	where 
% 	\begin{align}
% 	\frac{1}{\sigma}\frac{\partial \sigma_{y,2}(X)}{\partial \hat{\alpha}_p}  &= \left(1-\psi^{-1}\right)\frac{\gamma\sigma^2}{2 y_{0}(X)} \sum_{k=1}^J \left( \hat{\gamma}_k - \hat{\gamma}_0  \right) x_k\frac{1}{\sigma}\frac{\partial \sigma_{k,1}(X)}{\partial \hat{\alpha}_p}\nonumber\\
% 	&= \left(1-\psi^{-1}\right)\frac{\gamma\sigma^2}{2 y_{0}(X)} \sum_{j\in\mathcal{J}^u} \left( \hat{\gamma}_k - \hat{\gamma}_0  \right) x_k \left( -\frac{x_0}{x_u} \right) \nonumber\\
% 	&= \left(1-\psi^{-1}\right)\frac{\gamma\sigma^2}{2 y_{0}(X)}\left(\hat{\gamma}_0- \mathbb{E}^u[\hat{\gamma}_k] \right) x_0
% 	\end{align}
	
% 	Thus, the (inverse) aggregate market elasticity is then given by
% 	\begin{small}
% 		\begin{align}
% 		%\frac{\partial \ln p(X,\epsilon)}{\partial F(X,\epsilon)} 
% 		\frac{1}{p(X, \epsilon)} \frac{\partial p(X, \epsilon)}{\partial F(X, \epsilon)}
% 		&=  \left(1-\psi^{-1}\right)\frac{\gamma \sigma^2}{y_0(X)} \left[ \frac{1-\overline{\alpha}_p}{x_u} - \frac{\gamma_0-\mathbb{E}^u\left[\gamma_j \right]}{\gamma}\right]\nonumber\\
% 		&+\left(1-\psi^{-1}\right)\frac{\gamma \sigma^2}{y_0(X)} \left[\frac{ \mathbb{E}^u[\gamma_k]- \mathbb{E}^c[\gamma_k] }{\gamma x_u}+\frac{\gamma\sigma \left(1-\psi^{-1}\right)}{2 y_{0}(X)}\frac{\gamma_0- \mathbb{E}^u[\gamma_k]}{\gamma} + \frac{\overline{\alpha}_p - \frac{\overline{\sigma}}{\sigma}}{x_u}   \right]x_c\nonumber\\
% 		&+ \mathcal{O}(\epsilon^2), 
% 		\end{align}
% 	\end{small}
% where we use $ \overline{\sigma}=\sigma+\hat{\sigma}\epsilon $ from Equation~\eqref{eq:sig_bar-perturb}.
% \end{proof}

% \paragraph{\textcolor{red}{NEW VERSION: Fixing typos.}}


\begin{proof}
	Consider the case in which active investors have heterogeneous risk aversions and also face leverage constraints.
	We can then write $y_2(X)$ as follows
	\begin{align}
	y_2(X) &=  \left(1-\psi^{-1}\right) \gamma \sigma^2 \sum_{j=0}^J x_{j} \left[ \frac{\sigma_{j,1}^2(X)}{2 \sigma^2} + \hat{\gamma}_j \frac{\sigma_{j,1}(X)}{\sigma}  \right] \nonumber\\
	&=\left(1-\psi^{-1}\right) \gamma \sigma^2 \sum_{j=1}^J x_{j} \left[\frac{\sigma_{j,1}^2(X)}{2 \sigma^2}+ \hat{\gamma}_j \min\left\{ \sum_{k\in\mathcal{J}^u} \frac{x_k}{x_u} \hat{\gamma}_k -\frac{\hat{\alpha}_p  x_0+\frac{\hat{\sigma}}{\sigma} x_c }{1-x_0-x_c}- \hat{\gamma}_j  ,\frac{\hat{\sigma}}{\sigma}  \right\}  \right] \nonumber \\
	&+\left(1-\psi^{-1}\right) \gamma \sigma^2 x_0 \left[ \frac{\hat{\alpha}_p^2}{2} + \hat{\gamma}_0 \hat{\alpha}_p \right],
	\end{align}
	The derivative of $y_2(X)$ with respect to $\hat{\alpha}_p$ is given by
	\begin{align}
	\frac{\partial y_2(X)}{\partial \hat{\alpha}_p}   &=\left(1-\psi^{-1}\right) \gamma \sigma^2 \left(\sum_{j=1}^J x_{j} \frac{\sigma_{j,1}(X)}{ \sigma^2} \frac{\partial \sigma_{j,1}}{\partial \hat{\alpha}_p}-\mathbb{E}^u[\hat{\gamma}_k] x_0 + x_0 (\hat{\alpha}_p+ \hat{\gamma}_0) \right), \nonumber 
	\end{align}
	where
	\begin{align}
	\frac{1}{\sigma^2}\sum_{j=0}^J x_j \frac{\partial \sigma_{j,1}^2}{\partial \hat{\alpha}_p} &=\frac{1}{\sigma^2} \sum_{j=0}^J x_j2 \sigma_{j,1} \frac{\partial \sigma_{j,1}}{\partial \hat{\alpha}_p} \nonumber\\
	&= 2 x_0  \hat{\alpha}_p  + 2 \sum_{j\in\mathcal{J}^u} x_j \left( \sum_{k\in\mathcal{J}^u} \frac{x_k}{x_u} \hat{\gamma}_k -\left(\frac{\hat{\alpha}_p  x_0+\frac{\hat{\sigma}}{\sigma}x_c}{1-x_0-x_c}\right)- \hat{\gamma}_j\right) \left( -\frac{x_0}{1-x_0-x_c}\right) \nonumber \\
	&= 2 \left[ x_0 \hat{\alpha}_p + \frac{\hat{\alpha}_p x_0+\frac{\hat{\sigma}}{\sigma}x_c}{1-x_0-x_c}x_0+ \left(x_u \sum_{k\in\mathcal{J}^u} \frac{x_k}{x_u} \hat{\gamma}_k - \sum_{j\in\mathcal{J}^u} x_j\hat{\gamma}_j\right) \left( -\frac{x_0}{1-x_0-x_c}\right)\right] \nonumber \\
	&= 2 x_0 \left[\frac{\hat{\alpha}_p}{x_u} + \left(\frac{\hat{\sigma}}{\sigma} - \hat{\alpha}_p \right)\frac{x_c}{x_u}\right].
	\end{align}
	
	The second derivative of $y(X,\epsilon)$ can then be written as
	\begin{align}
	\frac{\partial y_2(X)}{\partial \hat{\alpha}_p}   &=\left(1-\psi^{-1}\right) \gamma \sigma^2 \left[\frac{\hat{\alpha}_p}{x_u}+ \left(\frac{\hat{\sigma}}{\sigma} - \hat{\alpha}_p \right)\frac{x_c}{x_u}- (\mathbb{E}^u[\hat{\gamma}_k]-\hat{\gamma}_0) \right] x_0. \nonumber 
	\end{align}
	
	Thus, the (inverse) aggregate market elasticity is then given by
	\begin{small}
		\begin{align}
		%\frac{\partial \ln p(X,\epsilon)}{\partial F(X,\epsilon)} 
		\frac{1}{p(X, \epsilon)} \frac{\partial p(X, \epsilon)}{\partial F(X, \epsilon)}
	&=  \left(1-\psi^{-1}\right)\frac{\gamma \sigma^2}{y_0(X)} \left[ \frac{(1-\overline{\alpha}_p)(1-x_c)-\left( \frac{\overline{\sigma}}{\sigma} -1\right)x_c}{x_u} - \frac{\gamma_0-\mathbb{E}^u\left[\gamma_j \right]}{\gamma}\right]+ \mathcal{O}(\epsilon^2),
		\end{align}
	\end{small}
where we use $ \overline{\sigma}=\sigma+\hat{\sigma}\epsilon $ from Equation~\eqref{eq:sig_bar-perturb}.
\end{proof}

%\begin{proof}
%We present the proof under three cases below.
%\paragraph{(i) Passive demand} 
%First, 
%
%\paragraph{(ii) Preference heterogeneity} 
%We next 
%
%\paragraph{(iii) Leverage constraints} 
%Finally, 
%\end{proof}


% \section{Things to do}
% \begin{itemize}
% 	\item Law of motion of wealth for the second-order (for completeness)
% 	\item Check the elasticity formulas (het.+leverage constraints)
% 	\item Obtain derivatives for $r$, $\theta$, and $\sigma_{y,2}$
% 	\item Code
% 	\item Calibration
% 	\item Section 3
% 	\item Elasticity decomposition
% 	\begin{itemize}
% 		\item From Equation~\eqref{eq:pricing_cond}, we have:
% 		\begin{equation*}
% 		y_2(X) = r_2(X) + \pi_2(X) - \mu_{P,2}(X)
% 		\end{equation*}
% 		\item Since $ \varepsilon_M^{-1}=-\dfrac{1}{y_{0}(X)} \dfrac{\partial y_{2}(X)}{\partial \hat{\alpha}_{p}} \dfrac{\epsilon}{x_{0}}+\mathcal{O}\left(\epsilon^{2}\right)$, we can decompose the market elasticity into the terms coming from the interest rate, risk premium, and drift of the price.
% 	\end{itemize}      
% \end{itemize}
\end{comment}

\section{Market Elasticity with Time-Varying Passive Portfolio Shares}

\noindent\textbf{Step 8 (time-varying $\alpha_p$): Aggregate market elasticity (no heterogeneity or leverage).}\\
Consider the case without preference heterogeneity or leverage constraints.
Then $\hat{\gamma}_j=0$ for $j=0,1,2$ and $\hat{\nu}_2=0$.
In the fixed-$\hat{\alpha}_p$ economy, Step 6 gives
\begin{equation}
    y_{Y,2}^{0}(x,\hat{\alpha}_p)
    =d_Y^{0}\hat{\alpha}_p^2,
    \qquad
    d_Y^{0}=\frac{1}{2}(1-\psi^{-1})\gamma\|\sigma\|^2\frac{x_0}{x_a},
\end{equation}
so $a_Y^{0}=b_Y^{0}=0$.

With time-varying passive demand, Step 7 implies that the total second-order dividend-yield correction is
\begin{equation}
    y_{Y,2}^{T}(x,\hat{\alpha}_p)
    \equiv y_{Y,2}^{0}(x,\hat{\alpha}_p)+\Delta y_{Y,2}(x,\hat{\alpha}_p)
    =a_Y^T+b_Y^T\hat{\alpha}_p+d_Y^T\hat{\alpha}_p^2,
\end{equation}
where, using $b_Y^{0}=0$,
\begin{align}
    d_Y^T
    &=\frac{1}{1+\frac{2\kappa_p}{y_{Y,0}}}d_Y^{0},\\
    b_Y^T
    &=\frac{2d_Y^T M_p}{y_{Y,0}+\kappa_p},\\
    a_Y^T
    &=a_Y
    =\frac{1}{y_{Y,0}}\left(b_Y^T M_p+d_Y^T V_p\right),
\end{align}
with
\begin{equation}
    M_p\equiv \kappa_p\hat{\alpha}+(1-\gamma)\sigma_p\sigma^\top,
    \qquad
    V_p\equiv \sigma_p\sigma_p^\top.
\end{equation}
The constant term $a_Y^T$ does not affect the elasticity.
Differentiating with respect to $\hat{\alpha}_p$ gives
\begin{equation}\label{eq:dy2-total-alpha-timevarying-noheterog}
    \frac{\partial y_{Y,2}^{T}}{\partial\hat{\alpha}_p}
    =b_Y^T+2d_Y^T\hat{\alpha}_p
    =2d_Y^T\left(\hat{\alpha}_p+\frac{M_p}{y_{Y,0}+\kappa_p}\right).
\end{equation}

Using the price-dividend-ratio expansion,
\begin{equation}
    \frac{1}{p(X,\epsilon)}\frac{\partial p(X,\epsilon)}{\partial F(X,\epsilon)}
    =-\frac{1}{y_{Y,0}}
    \frac{\partial y_{Y,2}^{T}}{\partial\hat{\alpha}_p}
    \frac{\epsilon}{x_0}
    +\mathcal{O}(\epsilon^2).
\end{equation}
Substituting \eqref{eq:dy2-total-alpha-timevarying-noheterog} and the expression for $d_Y^0$ yields
\begin{align}\label{eq:elasticity-timevarying-noheterog-hatalpha}
    \frac{1}{p(X,\epsilon)}\frac{\partial p(X,\epsilon)}{\partial F(X,\epsilon)}
    &=-\frac{1}{y_{Y,0}}
    \frac{2d_Y^{0}}{1+\frac{2\kappa_p}{y_{Y,0}}}
    \left(\hat{\alpha}_p+\frac{M_p}{y_{Y,0}+\kappa_p}\right)
    \frac{\epsilon}{x_0}
    +\mathcal{O}(\epsilon^2) \\
    &=-\left(1-\psi^{-1}\right)
    \frac{\gamma\|\sigma\|^2}{y_{Y,0}}
    \frac{1}{x_a}
    \frac{1}{1+\frac{2\kappa_p}{y_{Y,0}}}
    \left(\hat{\alpha}_p+\frac{M_p}{y_{Y,0}+\kappa_p}\right)
    \epsilon
    +\mathcal{O}(\epsilon^2).
\end{align}
Equivalently, using $\overline{\alpha}_p=1+\hat{\alpha}_p\epsilon$, this can be written as
\begin{equation}\label{eq:elasticity-timevarying-noheterog-final}
    \frac{1}{p(X,\epsilon)}\frac{\partial p(X,\epsilon)}{\partial F(X,\epsilon)}
    =\left(1-\psi^{-1}\right)
    \frac{\gamma\|\sigma\|^2}{y_{Y,0}}
    \frac{1}{x_a}
    \frac{1}{1+\frac{2\kappa_p}{y_{Y,0}}}
    \left[1-\overline{\alpha}_p
    -\frac{\epsilon M_p}{y_{Y,0}+\kappa_p}\right]
    +\mathcal{O}(\epsilon^2).
\end{equation}
The first term is the fixed-passive-share elasticity attenuated by mean reversion in passive demand.
The additional term involving $M_p$ captures the fact that predictable rebalancing and passive-demand risk affect prices even when the current passive share equals its benchmark value.
When $\kappa_p=0$ and $\sigma_p=0$, $M_p=0$ and \eqref{eq:elasticity-timevarying-noheterog-final} collapses to the fixed-$\hat{\alpha}_p$ result.

\noindent\textbf{Step 9 (time-varying $\alpha_p$): Aggregate market elasticity (preference heterogeneity).}\\
Consider next the case with preference heterogeneity but no leverage constraints.
Then $\hat{\nu}_2=0$, and the first-order exposure coefficients are
\begin{equation}
    m_{0,1}=\hat{\alpha}_p,
    \qquad
    m_{j,1}=\bar{\gamma}_a-\frac{x_0}{x_a}\hat{\alpha}_p-\hat{\gamma}_j,
    \qquad j=1,2,
\end{equation}
where
\begin{equation}
    \bar{\gamma}_a\equiv \sum_{j=1}^2\frac{x_j}{x_a}\hat{\gamma}_j,
    \qquad
    x_a\equiv 1-x_0.
\end{equation}
In the fixed-$\hat{\alpha}_p$ economy, the second-order dividend-yield correction is
\begin{equation}
    y_{Y,2}^{0}(x,\hat{\alpha}_p)
    =a_Y^{0}(x)+b_Y^{0}(x)\hat{\alpha}_p+d_Y^{0}(x)\hat{\alpha}_p^2,
\end{equation}
with
\begin{equation}
    b_Y^{0}(x)
    =(1-\psi^{-1})\gamma\|\sigma\|^2x_0(\hat{\gamma}_0-\bar{\gamma}_a),
    \qquad
    d_Y^{0}(x)
    =\frac{1}{2}(1-\psi^{-1})\gamma\|\sigma\|^2\frac{x_0}{x_a}.
\end{equation}
The constant term $a_Y^0$ does not affect the elasticity.

With time-varying passive demand, Step 7 implies that the total second-order dividend-yield correction satisfies
\begin{equation}
    y_{Y,2}^{T}(x,\hat{\alpha}_p)
    \equiv y_{Y,2}^{0}(x,\hat{\alpha}_p)+\Delta y_{Y,2}(x,\hat{\alpha}_p)
    =a_Y^T+b_Y^T\hat{\alpha}_p+d_Y^T\hat{\alpha}_p^2,
\end{equation}
where
\begin{equation}
    d_Y^T=\frac{1}{1+\frac{2\kappa_p}{y_{Y,0}}}d_Y^{0},
    \qquad
    b_Y^T=\frac{y_{Y,0}b_Y^{0}+2d_Y^T M_p}{y_{Y,0}+\kappa_p},
\end{equation}
and
\begin{equation}
    M_p\equiv \kappa_p\hat{\alpha}+(1-\gamma)\sigma_p\sigma^\top.
\end{equation}
Therefore,
\begin{align}\label{eq:dy2-total-alpha-timevarying-heterog}
    \frac{\partial y_{Y,2}^{T}}{\partial\hat{\alpha}_p}
    &=b_Y^T+2d_Y^T\hat{\alpha}_p \nonumber\\
    &=\frac{y_{Y,0}b_Y^{0}}{y_{Y,0}+\kappa_p}
    +2d_Y^T\left(\hat{\alpha}_p+\frac{M_p}{y_{Y,0}+\kappa_p}\right).
\end{align}
Using the price-dividend-ratio expansion,
\begin{equation}
    \frac{1}{p(X,\epsilon)}\frac{\partial p(X,\epsilon)}{\partial F(X,\epsilon)}
    =-\frac{1}{y_{Y,0}}
    \frac{\partial y_{Y,2}^{T}}{\partial\hat{\alpha}_p}
    \frac{\epsilon}{x_0}
    +\mathcal{O}(\epsilon^2).
\end{equation}
Substituting \eqref{eq:dy2-total-alpha-timevarying-heterog} and the expressions for $b_Y^0$ and $d_Y^0$ gives
\begin{align}\label{eq:elasticity-timevarying-heterog-hatalpha}
    \frac{1}{p(X,\epsilon)}\frac{\partial p(X,\epsilon)}{\partial F(X,\epsilon)}
    &=-\left(1-\psi^{-1}\right)\frac{\gamma\|\sigma\|^2}{y_{Y,0}}
    \left[
        \frac{y_{Y,0}}{y_{Y,0}+\kappa_p}(\hat{\gamma}_0-\bar{\gamma}_a)
        +\frac{1}{x_a}\frac{1}{1+\frac{2\kappa_p}{y_{Y,0}}}
        \left(\hat{\alpha}_p+\frac{M_p}{y_{Y,0}+\kappa_p}\right)
    \right]\epsilon
    +\mathcal{O}(\epsilon^2).
\end{align}
Equivalently, using $\overline{\alpha}_p=1+\hat{\alpha}_p\epsilon$ and
\begin{equation}
    \frac{\gamma_0-\mathbb{E}_a[\gamma_j]}{\gamma}
    =(\hat{\gamma}_0-\bar{\gamma}_a)\epsilon,
\end{equation}
we can write
\begin{align}\label{eq:elasticity-timevarying-heterog-final}
    \frac{1}{p(X,\epsilon)}\frac{\partial p(X,\epsilon)}{\partial F(X,\epsilon)}
    &=(1-\psi^{-1})\frac{\gamma\|\sigma\|^2}{y_{Y,0}}
    \left[
        \frac{1}{x_a}\frac{1}{1+\frac{2\kappa_p}{y_{Y,0}}}
        \left(1-\overline{\alpha}_p-\frac{\epsilon M_p}{y_{Y,0}+\kappa_p}\right)
        -\frac{y_{Y,0}}{y_{Y,0}+\kappa_p}
        \frac{\gamma_0-\mathbb{E}_a[\gamma_j]}{\gamma}
    \right]
    +\mathcal{O}(\epsilon^2).
\end{align}
Thus, time variation in passive demand attenuates both components of the fixed-passive-share elasticity, but with different attenuation factors.
The demand-misalignment term is scaled by $(1+2\kappa_p/y_{Y,0})^{-1}$, while the preference-heterogeneity term is scaled by $y_{Y,0}/(y_{Y,0}+\kappa_p)$.
The term involving $M_p$ is new and captures predictable rebalancing and passive-demand risk.
When $\kappa_p=0$ and $\sigma_p=0$, \eqref{eq:elasticity-timevarying-heterog-final} collapses to the fixed-$\hat{\alpha}_p$ result with preference heterogeneity.


\section{Useful Formula}
The following are useful for computing the derivatives above:
\begin{align}
\gamma \sigma \frac{\partial \sigma_{y,2}(X)}{\partial \hat{\alpha}_p} &= \left(1-\psi^{-1}\right)\frac{\gamma\sigma^2}{2 y_{0}(X)} \sum_{k=1}^J \left( \hat{\gamma}_k - \hat{\gamma}_0  \right) x_k \left(- \frac{\gamma \sigma^2 x_0}{1-x_0} \right).
\end{align}
Note that we can write $\sigma_{j,2}(X)$ as follows
\begin{align}
\frac{\sigma_{j,2}(X)}{\sigma} &= \frac{\theta_2(X)}{\gamma \sigma} - \frac{\theta_1(X)}{\gamma\sigma}\hat{\gamma}_j + \frac{\theta_0(X)}{\gamma\sigma}\hat{\gamma}_j^2 - (1-\gamma^{-1}) \frac{\sigma_{y,2}(X)}{\sigma}\nonumber\\
&= -\left( 1 + \frac{ x_c}{x_u}\right)\frac{\sigma_{y,2}(X)}{\sigma}-\mathbb{E}^u[\hat{\gamma}_k] \frac{ \hat{\alpha}_p x_0+ \frac{\hat{\sigma}}{\sigma}x_c }{1-x_0-x_c} -\var^u[\hat{\gamma}_k] +\nonumber\\
& - \hat{\gamma}_j \left[\sum_{k\in\mathcal{J}^u} \frac{x_k}{x_u} \hat{\gamma}_j -\frac{\hat{\alpha}_p  x_0}{1-x_0}\right]+\hat{\gamma}_j^2.
\end{align}

\begin{align}
\sigma_{j,1}(X) &= \sigma\min\left\{ \sum_{k\in\mathcal{J}^u} \frac{x_k}{x_u} \hat{\gamma}_k -\left(\frac{\hat{\alpha}_p  x_0+\frac{\hat{\sigma}}{\|\sigma\|} x_c }{1-x_0-x_c}\right)- \hat{\gamma}_j  ,\frac{\hat{\sigma}}{\sigma}  \right\}\\
\sigma_{j,2}(X) &= \frac{\theta_2(X)}{\gamma} - \frac{\theta_1(X)}{\gamma}\hat{\gamma}_j + \frac{\theta_0(X)}{\gamma}\hat{\gamma}_j^2 - (1-\gamma^{-1})   \sigma_{y,2}(X)\\
\sigma_{y,2}(X)  &= \left (1-\psi^{-1}\right )\frac{\gamma\sigma^2}{2 y_{0}(X)} \sum_{k=1}^J \left( \hat{\gamma}_k - \hat{\gamma}_0  \right) x_k \sigma_{k,1}(X)\\
\theta_1(X) &= \gamma \sigma \left[\sum_{j\in\mathcal{J}^u} \frac{x_j}{x_u} \hat{\gamma}_j -\frac{\hat{\alpha}_p  x_0}{1-x_0}\right]\\
\theta_2(X) &=  - \frac{(\gamma-1) x_c + 1-x_0}{1-x_0-x_c} \sigma_{y,2}(X) - \gamma \sigma \mathbb{E}^u[\hat{\gamma}_j] \frac{ \hat{\alpha}_p x_0+ \frac{\hat{\sigma}}{\sigma}x_c }{1-x_0-x_c} - \gamma \sigma \var^u[\hat{\gamma}_j],
\end{align}
where
\begin{align}
\frac{\partial \sigma_{y,2}}{\partial \hat{\alpha}_p} &= \left (1-\psi^{-1}\right )\frac{\gamma\sigma^2}{2 y_{0}(X)} \sum_{k=1}^J \left( \hat{\gamma}_k - \hat{\gamma}_0  \right) x_k \left(- \frac{\sigma x_0}{1-x_0} \right)\nonumber\\
&=\left (1-\psi^{-1}\right )\frac{\gamma\sigma^2}{2 y_{0}(X)}\sigma\left(\hat{\gamma}_0-\E^u[\hat{\gamma}_j]\right)x_0\\
\sum_{j=0}^J x_{j}\gamma \sigma \hat{\gamma}_j \frac{\partial \sigma_{j,1}(X)}{\partial \hat{\alpha}_p} &= \gamma \sigma^2 \sum_{j=0}^J x_j \hat{\gamma}_j \left( -\frac{x_0}{1-x_0} \right)
\end{align}

\section{Derivation of the perturbed solution}

In this section, we compute the first-order and second-order correction of the equilibrium objects.
It turns out that the system of equations determining the perturbed solution is block-recursive, so we are able to solve for the equilibrium objects one by one, provided we proceed in the appropriate order.

In contrast to the case considered in the text, we allow for portfolio-flow shocks.
In particular, we assume that the portfolio share of the passive investor is given by $\alpha_{0,t} = 1 + \epsilon (\overline{\alpha}_{p,t} -1)$, where $\overline{\alpha}_{p,t}$ follows the process
\begin{equation}
    d \overline{\alpha}_{p,t} = \theta_p (\overline{\alpha} - \overline{\alpha}_{p,t}) \epsilon dt + \sigma_{p} \sqrt{\overline{\alpha}_{p,t}} \epsilon d Z_t.
\end{equation}
Notice that $\alpha_{0,t} = 1$ and $\overline{\alpha}_{p,t}$ is constant when $\epsilon = 0$.
Finally, we assume that the mortality parameter is given by $\kappa = \hat{\kappa} \epsilon$.

\subsection{First-order correction}

\paragraph{Diffusion and drift terms.} The diffusion term for the price-dividend ratio is given by
\begin{equation}
    \sigma_{p,t} = \frac{p_{x}}{p} \sigma_{x} + \frac{p_{\overline{\alpha}_p}}{p} \sigma_{p} \sqrt{\overline{\alpha}_{p,t}} \epsilon = \mathcal{O}(\epsilon^2).
\end{equation}
Notice that $p_{x_j} = \mathcal{O}(\epsilon)$ and $\sigma_{x_j} = \mathcal{O}(\epsilon)$, as $p$ and $x_j$ are constant when $\epsilon = 0$, so the zeroth-order terms for $p_{x_j}$, $p_{\overline{\alpha}_p}$, and $\sigma_{x_j}$ are equal to zero.
This implies that the first-order correction for $\sigma_{p,t}$ is equal to zero.
A similar argument shows that $\sigma_{c_j,t} = \mathcal{O}(\epsilon^2)$.

The drift of $p$ is given by
\begin{equation}
    \mu_{p,t} = \frac{p_{x}}{p} \mu_{x} + \frac{p_{\overline{\alpha}_p}}{p} \theta_{p} (\overline{\alpha} -\overline{\alpha}_{p,t}) \epsilon + \frac{1}{2} \sum_{k=1}^d\left[  \sigma_{x,k}^\top\frac{ p_{xx}}{p} \sigma_{x,k} + 2 \sigma_{p,k} \sqrt{\overline{\alpha}_{p,t}} \epsilon \frac{ p_{x\alpha_p}}{p} \sigma_{x,k} + \frac{p_{\overline{\alpha}_p\overline{\alpha}_p}}{p} \sigma_{p,k}^2 \overline{\alpha}_{p,t} \epsilon^2 \right],
\end{equation}
where $\sigma_{x,k}$ is the $k$-th column of the $J\times d$ matrix $\sigma_x$.
As $p_{x}$ and $p_{\overline{\alpha}_p}$ are first-order in $\epsilon$, and the same goes for $\mu_x$ and $\sigma_x$, then $\mu_{p,t} = \mathcal{O}(\epsilon^2)$.
A similar argument shows that $\mu_{c_j,t} = \mathcal{O}(\epsilon^2)$.
Notice these facts imply that $\varsigma_j = \mathcal{O}(\epsilon^2)$ and $\xi_{j,t} = \mathcal{O}(\epsilon^2)$.

\paragraph{Risk premium.} The risk premium is given by
\begin{equation}
    \pi_1(X) =  \gamma ||\sigma||^2 \left[  \sum_{j\in \mathcal{J}^u} \frac{x_{j}}{x_u} \hat{\gamma}_j -  \frac{x_0 \hat{\alpha}_{p} + x_c \frac{\hat{\sigma}}{||\sigma||}}{1-x_0 - x_c}\right],
\end{equation}
using the fact that $||\sigma_{R,t}|| = ||\sigma|| + \mathcal{O}(\epsilon^2)$, and $\hat{\alpha}_{p,t} \equiv \overline{\alpha}_{p,t}-1$.

\paragraph{Portfolio share.} The portfolio share of an unconstrained investor is given by
\begin{equation}
    \alpha_{j}(X,\epsilon) = 1 + \left[ \frac{\pi_1(X)}{\gamma ||\sigma||^2} - \hat{\gamma}_j \right] \epsilon + \mathcal{O}(\epsilon^2),
\end{equation}
the portfolio share of a constrained investor is given by
\begin{equation}
    \alpha_j(X,\epsilon) = 1 + \frac{\hat{\sigma}}{||\sigma||} \epsilon + \mathcal{O}(\epsilon^2),
\end{equation}
and the portfolio share of the passive investor is given by $\alpha_{0}(X,\epsilon) = 1 + \hat{\alpha}_{p,t} \epsilon$.
Notice that $\sum_{j=0}^J x_j \alpha_{j,1}(X) = 0$, consistent with market clearing.

\paragraph{Interest rate.} The first-order correction for the interest rate is given by
\begin{equation}
    r_1(X) = \left(1-\psi^{-1}\right) \gamma ||\sigma||^2\sum_{j=0}^J x_j  \left[ \frac{\hat{\gamma}_j}{2} + \alpha_{j,1}(X)\right]  - \pi_1(X),
\end{equation}
using the fact that $\xi_t = \mathcal{O}(\epsilon^2)$, $\mu_{p,t} = \mathcal{O}(\epsilon^2)$, and $\sigma_{p,t} = \mathcal{O}(\epsilon^2)$.
Given the market clearing for the risky asset, we can write:
\begin{equation}
    r_1(X) = \left(1-\psi^{-1}\right) \gamma ||\sigma||^2\sum_{j=0}^J x_j  \frac{\hat{\gamma}_j}{2}  - \pi_1(X),
\end{equation}
so $r_1(X) + \pi_1(X)$ is independent of $\overline{\alpha}_{p}$.

\paragraph{Price-dividend ratio.} From the pricing condition, we obtain
\begin{equation}
    - \frac{1}{p_0(X)^2} p_1(X) = r_1(X) + \pi_1(X).
\end{equation}
Rearranging the expression above, and using the expression for the interest rate, we obtain
\begin{equation}
    p_1(X) = - p_0(X)^2 \left(1-\psi^{-1}\right) \gamma ||\sigma||^2\sum_{j=0}^J x_j  \frac{\hat{\gamma}_j}{2},
\end{equation}
which is independent of $\overline{\alpha}_{p,t}$.

\paragraph{Consumption-wealth ratio.} The consumption-wealth ratio is given by
\begin{equation}
    c_{j,1}(X) = (1-\psi) \left[ r_1(X) + \pi_1(X) + \pi_0(X) \alpha_{j,1}(X)  - \frac{1}{2} \gamma \||\sigma\||^2 (\hat{\gamma}_j + 2 \alpha_{j,1}(X))\right].
\end{equation}
Using the expression for $r_1(X)$, we can write the expression above as follows:
\begin{equation}
    c_{j,1}(X) = (1-\psi) \left[ (1-\psi^ {-1}) \sum_{i=0}^J x_i  \frac{\hat{\gamma}_i}{2}  - \frac{\hat{\gamma}_j }{2} \right] \gamma \||\sigma\||^2.
\end{equation}

\paragraph{Wealth dynamics.} The diffusion term of $x_j$ is given by
\begin{equation}
    \sigma_{x_j}(X) = x_j \alpha_{j,1}(X) \epsilon \sigma + \mathcal{O}(\epsilon^2).
\end{equation}

The drift of $x_j$ is given by
\begin{equation}
    \mu_{x_j}(X) = x_j\left[r_1(X) + \pi_1(X) + \pi_0(X) \alpha_{j,1}(X) - c_{j,1}(X)- \alpha_{j,1}(X)||\sigma||^2 + \hat{\kappa} \frac{\omega_j - x_j}{x_j} \right] \epsilon + \mathcal{O}(\epsilon^2).
\end{equation}

We can write the first-order correction of $\mu_{x_j}$ as follows:
\begin{equation}
    \mu_{x_j,1}(X) = x_j\left[(\psi-1) \frac{\gamma \||\sigma\||^ 2 }{2} \left(\sum_{i=0}^J x_i \hat{\gamma}_i -\hat{\gamma}_j \right) + (\gamma-1)\||\sigma\||^ 2 \alpha_{j,1}(X) \right] + \hat{\kappa} (\omega_j - x_j).
\end{equation}

\subsection{Second-order correction}

\paragraph{Diffusion and drift terms.} The diffusion term for the price-dividend ratio is given by
\begin{equation}
    \sigma_{p,2}(X) = \frac{p_{x,1}(X)}{p_0(X)} \sigma_{x,1}(X) + \frac{p_{\overline{\alpha}_p,1}(X)}{p_0(X)} \sigma_{p} \sqrt{\overline{\alpha}_{p}}.
\end{equation}
We can write the expression above as follows:
\begin{equation}
    \sigma_{p,2}(X) = - p_0(X) (1-\psi^{-1})\frac{\gamma \|\sigma\||^ 2}{2} \sum_{j=1}^J \left( \hat{\gamma}_j - \hat{\gamma}_0 \right) x_j \alpha_{j,1}(X) \sigma,
\end{equation}
where we used the fact that $p_{\overline{\alpha}_p,1}(X) = 0$.

Similarly, the diffusion for $c_j$ is given by
\begin{equation}
    \sigma_{c_j,2}(X) = -(\psi-1) \frac{\gamma \||\sigma\||^2}{2 c_{j,0}(X)} (1-\psi^{-1}) \sum_{i=1}^J\left( \hat{\gamma}_i - \hat{\gamma}_0 \right) x_i \alpha_{i,1}(X) \sigma.
\end{equation}
The second-order correction for the hedging demand is then given by $\varsigma_{j,2}(X) = \frac{1-\gamma^{-1}}{\psi-1} \frac{\sigma_{c_j,2} \sigma^\top }{\||\sigma\||^2}$.

The second-order correction of the drift of $p$ and $c_{j}$ are given by
\begin{equation}
    \mu_{p,2}(X) = \frac{p_{x,1}(X)}{p_0(X)} \mu_{x,1}(X), \qquad \qquad 
    \mu_{c_j,2}(X) = \frac{c_{j,x,1}(X)}{c_{j,0}(X)} \mu_{x,1}(X),
\end{equation}
which can be written as
\begin{align}
    \mu_{p,2}(X) &= -p_0(X) (1-\psi^{-1}) \frac{\gamma\|\sigma\||^2}{2} \sum_{j=1}^J (\hat{\gamma}_j-\hat{\gamma}_0) \mu_{x_j,1}(X)\\
    \mu_{c_i,2}(X) &= (1-\psi)(1-\psi^{-1}) \frac{\gamma\|\sigma\||^2}{2 c_{i,0}(X)} \sum_{j=1}^J (\hat{\gamma}_j-\hat{\gamma}_0) \mu_{x_j,1}(X).
\end{align}

The second-order correction for $\xi_{j,t}$ is then given by $\xi_{j,2}(X) = \mu_{c_j,2}(X) + (1-\gamma) \sigma_{c_j,2} \sigma^\top$.

\paragraph{Risk premium.} The risk premium is given by
\begin{equation}
    \pi_2(X) = \gamma \|\sigma\||^2\left[ \frac{\gamma_{u,2}(X)}{\gamma} - \frac{\gamma_{u,1}(X)}{\gamma} \frac{x_0 \hat{\alpha}_{p} + x_c \frac{\hat{\sigma}}{\|\sigma\||}}{1-x_0-x_c}  + 
    2\sum_{k=1}^d \frac{\sigma_k \sigma_{p,2,k}}{||\sigma||^ 2}
    -\frac{x_c}{1-x_0 -x_c} \overline{\alpha}_{c,2}(X) - \varsigma_{2}(X) \right].
\end{equation}

Notice that we can write the aggregate risk aversion as follows:
\begin{equation}
    \gamma_u(X) = \gamma \left[ 1 + \mathbb{E}^u[\hat{\gamma}_j] \epsilon - \delta^u[\hat{\gamma}_j] \epsilon^ 2 \right] + \mathcal{O}(\epsilon^3),
\end{equation}
where $\mathbb{E}^u[\hat{\gamma}_j] \equiv \sum_{j\in\mathcal{J}^u} \frac{x_j}{x_u} \hat{\gamma}_j$ and $\delta^u[\hat{\gamma}_j] \equiv \sum_{j\in\mathcal{J}^u} \frac{x_j}{x_u} \hat{\gamma}_j^ 2 - \left(\sum_{j\in\mathcal{J}^u} \frac{x_j}{x_u} \hat{\gamma}_j \right)^ 2$, so $\gamma_{u,2}(X) / \gamma = -\delta^u[\hat{\gamma}_j]$.

Combining the previous two expressions, we obtain
\begin{align}
    \frac{\pi_2(X)}{\gamma \|\sigma\||^2} &=  -\delta^u[\hat{\gamma}_j] - \mathbb{E}^u[\hat{\gamma}_j] \frac{x_0 \hat{\alpha}_{p} + x_c \frac{\hat{\sigma}}{\|\sigma\||}}{1-x_0-x_c}  + 
    2\sum_{k=1}^d \frac{\sigma_k \sigma_{p,2,k}}{||\sigma||^ 2}
    +\frac{x_c}{1-x_0 -x_c} \sum_{k=1}^d  \frac{\sigma_k \sigma_{p,2,k}}{\|\sigma\||^2} - \varsigma_{2}(X) \nonumber\\
     &=  -\delta^u[\hat{\gamma}_j] - \mathbb{E}^u[\hat{\gamma}_j] \frac{x_0 \hat{\alpha}_{p} + x_c \frac{\hat{\sigma}}{\|\sigma\||}}{1-x_0-x_c}  
    +\left(1+\gamma^{-1}+\frac{x_c}{1-x_0 -x_c}\right) \sum_{k=1}^d  \frac{\sigma_k \sigma_{p,2,k}}{\|\sigma\||^2},
\end{align}
where we used the fact that $\varsigma_{2}(X) = \frac{1-\gamma^{-1}}{\psi-1} \frac{\sigma_{c_j,2}\sigma^\top}{\|\sigma\||^2}$, $\sigma_{p,2} = \frac{\sigma_{c_j,2}}{\psi-1}$, and $\overline{\alpha}_{c,2}(X) = - \sum_{k=1}^d  \frac{\sigma_k \sigma_{p,2,k}}{\|\sigma\||^2}$.

\paragraph{Portfolio share.} The portfolio share of an unconstrained investor is given by
\begin{equation}
    \alpha_{j,2}(X) = \frac{\pi_2(X)}{\gamma \||\sigma\||^2} - \frac{\pi_1(X)}{\gamma \||\sigma\||^2} \hat{\gamma}_j + \hat{\gamma}_j^2 - 2 \sum_{k=1}^d \frac{\sigma_k \sigma_{p,2,k}(X)}{\|\sigma\|^2} + \varsigma_{j,2}(X),
\end{equation}
the portfolio share of a constrained investor is $\alpha_{j,2}(X) = - \sum_{k=1}^d  \frac{\sigma_k \sigma_{p,2,k}}{\|\sigma\||^2}$, and the portfolio share of the passive investor satisfies $\alpha_{0,2} = 0$.

We can write the expression above as follows:
\begin{equation}
    \alpha_{j,2}(X) = \frac{\pi_2(X)}{\gamma \||\sigma\||^2} - \frac{\pi_1(X)}{\gamma \||\sigma\||^2} \hat{\gamma}_j + \hat{\gamma}_j^2 - (1+\gamma^{-1}) \sum_{k=1}^d \frac{\sigma_k \sigma_{p,2,k}(X)}{\|\sigma\|^2}.
\end{equation}
Notice that $\sum_{j=0}^J x_j \alpha_{j,2}(X) = 0 $, consistent with market clearing.


\paragraph{Interest rate.} The interest rate is given by
\begin{align}
    r_2(X) &= \psi^{-1} (\mu_{p,2}(X) + \sigma \sigma_{p,2}(X)^\top )  + (1-\psi^{-1}).\frac{\gamma \|\sigma\||^ 2}{2} \sum_{j=0}^J x_j \left[ 2\hat{\gamma}_j \alpha_{j,1}(X) + \alpha_{j,1}^2(X) + 2 \alpha_{j,2}(X)\right]\nonumber\\
    &+(1-\psi^{-1}) \gamma \|\sigma\|^ 2 \sum_{k=1}^d \frac{\sigma_k \sigma_{p,2,k}(X)}{\|\sigma\|^2} -\pi_2(X) + \psi^{-1} \xi_{2}(X),
\end{align}
where $\xi_2(X) = (\psi-1) \left[ \mu_{p,2}(X) + (1-\gamma) \sigma_{p,2} \sigma^\top \right]$

We can write the expression for the portfolio of the unconstrained investor as follows:
\begin{align}
    r_2(X) &= \mu_{p,2}(X) + \sigma \sigma_{p,2}(X)^\top   + (1-\psi^{-1}) \gamma \|\sigma\||^ 2 \sum_{j=0}^J x_j \left[ \hat{\gamma}_j \alpha_{j,1}(X) + \frac{\alpha_{j,1}^2(X)}{2} \right] -\pi_2(X),
\end{align}

\paragraph{Price-dividend ratio.} The price-dividend ratio is given by
\begin{equation}
\frac{p_1^2(X)}{p_0^3(X)} - \frac{p_2(X)}{p_0^2(X)} = r_2(X) + \pi_2(X) - \mu_{p,2}(X) - \sigma \sigma_{p,2}^\top.
\end{equation}
Rearranging the expression above, and using the expression for $r_2(X)$, we obtain
\begin{equation}
    \frac{p_2(X)}{p_0(X)} = - p_0(X)  (1-\psi^{-1}) \gamma \|\sigma\||^ 2 \sum_{j=0}^J x_j \left[ \hat{\gamma}_j \alpha_{j,1}(X) + \frac{\alpha_{j,1}^2(X)}{2} \right]  + \left( \frac{p_1(X)}{p_0(X)} \right)^2.
\end{equation}

\paragraph{Consumption-wealth ratio.} The second-order correction for the consumption-wealth ratio is given by
\begin{align}
    c_{j,2}(X) &= (1-\psi) \left[ r_2(X) + \pi_2(X) + \pi_1(X) \alpha_{j,1}(X) + \pi_0(X) \alpha_{j,2}(X) \right] + \xi_{j,2}(X)\\
    &-(1-\psi) \gamma \|\sigma\||^2 \left[ \alpha_{j,1}(X) \hat{\gamma}_j + \alpha_{j,2}(X) + \frac{\alpha_{j,1}^2(X)}{2} + \sum_{k=1}^d \frac{\sigma_k \sigma_{p,k,2}(X)}{\|\sigma\||^2} \right].
\end{align}

\paragraph{Wealth dynamics.} The diffusion term of $x_j$ is given by
\begin{equation}
    \sigma_{x_j,2}(X) = x_j \alpha_{j,2} \sigma.
\end{equation}

The drift of $x_j$ is given by
\begin{equation}
    \mu_{x_j,2} = x_j \left[ r_2(X) + \pi_2(X) + \pi_1(X) \alpha_{j,1}(X) + \pi_0(X) \alpha_{j,2}(X) - c_{j,2}(X)- \mu_{p,2}(X)- \sigma \sigma_{p,2}(X)^\top - \alpha_{j,2}(X) \|\sigma\|^2 \right].
\end{equation}

\paragraph{Aggregate market elasticity.} The derivative of $p$ with respect to $\overline{\alpha}_p$ is given by
\begin{equation}
    \frac{1}{p(X,\epsilon)}\frac{\partial p(X,\epsilon)}{\partial \overline{\alpha}_p} = \frac{1}{p_0(X)} \frac{\partial p_2(X)}{\partial \overline{\alpha}_p} \epsilon^2 + \mathcal{O}(\epsilon^3).
\end{equation}

The market elasticity satisfies the condition
\begin{equation}
    \frac{1}{p_0(X)} \frac{\partial p_2(X)}{\partial \overline{\alpha}_p} = - p_0(X)  (1-\psi^{-1}) \gamma \|\sigma\||^ 2 \left[ x_0(\hat{\gamma}_0 + \hat{\alpha}_{p}) + \sum_{j\in\mathcal{J}^u} x_j \left( \hat{\gamma}_j + \alpha_{j,1}(X) \right)  \left(- \frac{x_0}{x_u} \right) \right] \epsilon^2.
\end{equation}
From the market clearing for the risky asset, we have $x_0 \hat{\alpha}_p + \sum_{j \in \mathcal{J}^u} x_j \alpha_{j,1}(X) + x_c \frac{\hat{\sigma}}{\|\sigma\|} = 0$, so we can write the expression above as follows:
\begin{equation}
    \frac{1}{p_0(X)} \frac{\partial p_2(X)}{\partial \overline{\alpha}_p} = p_0(X)  (1-\psi^{-1}) \gamma \|\sigma\||^ 2 \left[ \hat{\gamma}_u(X) - \hat{\gamma}_0 + \frac{1-x_c}{1-x_0 -x_c}  (1-\overline{\alpha}_p) -  \frac{ x_c \frac{\hat{\sigma}}{\|\sigma\||}}{1-x_0 - x_c}  \right] x_0 \epsilon^2.
\end{equation}


\subsection{Third-order approximation}

\subsubsection{Passive demand}

Suppose there is no preference heterogeneity and no leverage constraint.
Without loss of generality, set $J = 1$.
In this case, the price-dividend ratio is given by
\begin{align}
    p(X,\epsilon) &= p^* -(p^*)^2 (1-\psi^{-1}) \frac{\gamma \|\sigma\|^2}{2} \frac{x_0 \hat{\alpha}_p^2}{x_1} \epsilon^2 + \mathcal{O}(\epsilon^3)\\
    \pi(X,\epsilon) &= \pi_0(X) - \gamma \|\sigma\|^2 \frac{x_0\hat{\alpha}_p}{x_1} \epsilon  + \mathcal{O}(\epsilon^3)\\
    r(X,\epsilon) &= r_0(X)+\gamma \|\sigma\|^2 \frac{x_0 \hat{\alpha}_p}{x_1} \epsilon +(1-\psi^{-1}) \frac{\gamma \|\sigma\|^2}{2} \frac{x_0 \hat{\alpha}_p^2}{x_1} \epsilon^2 + \mathcal{O}(\epsilon^3)\\
    \alpha_{0}(X,\epsilon) &= 1 + \hat{\alpha}_p\epsilon + \mathcal{O}(\epsilon^3)\\
    \alpha_{1}(X,\epsilon) &= 1 - \frac{x_0}{x_1} \hat{\alpha}_p\epsilon + \mathcal{O}(\epsilon^3)\\
    c_{j}(X,\epsilon) &= c_{j,0}(X) + (1-\psi^{-1}) \frac{\gamma \|\sigma\|^2}{2} \frac{x_0 \hat{\alpha}_p^2}{x_1} \epsilon^2 + \mathcal{O}(\epsilon^3)\\
    \sigma_{x_1}(X,\epsilon) &= - \frac{x_0}{x_1} \hat{\alpha}_p \sigma \epsilon  + \mathcal{O}(\epsilon^3)\\
    \mu_{x_1}(X,\epsilon) &= \left[ (1-\gamma) \|\sigma\|^2 (1-x_1) \hat{\alpha}_p + \hat{\kappa} (\omega_j - x_1)\right]  \epsilon + \mathcal{O}(\epsilon^3),
\end{align}
where $\mu_p, \sigma_p, \mu_{c_j},  $ and $\sigma_{c_j}$ are all equal to zero up to second order, and $x_0 = 1-x_1$.

The law of motion of $x_{1,t}$ can be written as
\begin{equation}
    \frac{d x_{1,t}}{x_{1,t}} = \left[ x_{0,t} (c_{0,t}-c_{1,t})+x_{0,t}(\alpha_{0,t}-\alpha_{1,t})\left( \|\sigma_{R,t}\|^2-\pi_t \right) 
 + \kappa\frac{\omega_1 - x_{1,t}}{x_{1,t}}\right] dt + (\alpha_{1,t}-1) \sigma_{R,t} d Z_t.
\end{equation}

% Notice that $\alpha_{1,t} = 1 + \frac{1-x_{1,t}}{x_{1,t}} (1-\alpha_{0,t})$. Define $z = \frac{1-x_1}{x_1}$, so the diffusion term for $z_t$ is given by
% \begin{equation}
%     \sigma_{z,t} = -\frac{1}{x_{1,t}^2} (1-x_{1,t}) (1-\alpha_{0,t}) \sigma_{R,t} = - \frac{z_{1,t}}{1+z_{1,t}} (1-\alpha_{0,t}) \sigma_{R,t}.
% \end{equation}






\paragraph{Diffusion and drift terms.} The derivatives of $p(X,\epsilon)$ with respect to $x_1$ and $\overline{\alpha}_p$ are given by
\begin{equation}
    \frac{p_{x_1}(X,\epsilon)}{p^*} = (1-\psi^{-1})\frac{\gamma\|\sigma\|^2}{2 y^*} \frac{1}{x_1^2} \hat{\alpha}_p^2 \epsilon^2 + \mathcal{O}(\epsilon^3), \qquad  
    \frac{p_{\overline{\alpha}_p}(X,\epsilon)}{p^*} = -(1-\psi^{-1}) \frac{\gamma\|\sigma\|^2}{y^*} \frac{x_0}{x_1} \hat{\alpha}_p \epsilon^2 + \mathcal{O}(\epsilon^3)
\end{equation}

The diffusion term for $p(X,\epsilon)$ is then given by
\begin{equation}
    \sigma_{p,3}(X) =  -(1-\psi^{-1}) \frac{\gamma\|\sigma\|^2}{y^*} \frac{x_0}{x_1} \left[   \frac{\hat{\alpha}_p^3}{2x_1^2} \sigma + \hat{\alpha}_p \sigma_p \sqrt{\overline{\alpha}_p} \right].
\end{equation}
and $\sigma_{c_j,3}(X) = \sigma_{p,3}(X)$.
Notice that excess volatility depends on the market elasticity times the volatility of portfolio flows.

The drift of $p$ is given by
\begin{equation}
    \mu_{p,3}(X) = (1-\psi^{-1})\frac{\gamma\|\sigma\|^2}{2 y^*} \frac{1}{x_1^2} \hat{\alpha}_p^2 \mu_{x_1,1} -(1-\psi^{-1}) \frac{\gamma\|\sigma\|^2}{y^*} \frac{x_0}{x_1} \hat{\alpha}_p \theta_p (\overline{\alpha}-\overline{\alpha}_{p,t}),
\end{equation}
and $\mu_{c_j,3}(X) = \mu_{p,3}(X)$.

\paragraph{Risk premium.} The risk premium is given by
\begin{equation}
\pi(X) = \gamma \|\sigma\|^2 \left[ 1 - \frac{x_0(\overline{\alpha}_p -1) }{x_1}  -  \frac{1-\gamma^{-1}}{\psi-1} \frac{\gamma\|\sigma\|^2}{y^*} \frac{x_0}{x_1} \left[   \frac{(1-\overline{\alpha}_p)^3}{2x_1^2} \sigma + \hat{\alpha}_p \sigma_p \sqrt{\overline{\alpha}_p} \right] \right]
\end{equation}

\section{Higher-order perturbations}

Suppose we have the $(n-1)$-th order perturbation of $c_{j}(X,\epsilon) = \sum_{k=0}^{n-1}c_{j,k}(X) \epsilon^k$ and the law of motion of $X$.
Let $lc_{j}(X,\epsilon) = \sum_{k=0}^{n-1} lc_{j,k}(X) \epsilon^k$ denote the expansion of log $c_j(X,\epsilon)$.
Then we can compute $\sigma_{j}(X)$ up to order $n$:
\begin{equation}
    \sigma_{c_j,n}(X) = \sum_{k=1}^{n-1} lc_{j,k,X}(X) \sigma_{X,n-k}(X),
\end{equation}
which is independent of the $n$-th order term in $c_{j}(X,\epsilon)$ and $\sigma_X(X,\epsilon)$, as $lc_{j,0,X}(X) = 0$.
Similarly, we can compute $\sigma_{y,n}(X)$.
A similar argument gives $\mu_{p,n}(X)$ and $\mu_{c_j,n}(X)$.
We can then compute $\pi_n(X)$ and $\alpha_{j,n}(X)$.
The $n$-th term of the consumption-wealth ratio satisfies the condition
\begin{equation}
c_{j,t} = \psi \rho + (1-\psi) \left[ \pi_t (\alpha_{j,t}-1) + \mu + \mu_{p,t} + \sigma \sigma_{p,t}^\top - \frac{\gamma_j}{2} \|\sigma_{R,t}\|^2 \alpha_{j,t}^2 + \sum_{j=0}^J x_j c_{j,t} \right] + \xi_{j,t}
\end{equation}

We can rewrite the system above in matrix form as follows:
\begin{align}
    [ I - (1-\psi) \textbf{1}_{J+1} x_t^\top] c_t = \zeta_t,
\end{align}
where $c_t = [c_{0,t},\ldots, c_{J,t}]^\top$, $x_t=[x_{0,t},\ldots ,x_{J,t}]^\top$, $\textbf{1}_{J+1}$ is a $(J+1)$-th dimensional vector of ones, and
$\zeta_{j,t} \equiv \psi \rho + (1-\psi) \left[ \pi_t (\alpha_{j,t}-1) + \mu + \mu_{p,t} + \sigma \sigma_{p,t}^\top - \frac{\gamma_j}{2} \|\sigma_{R,t}\|^2 \alpha_{j,t}^2 \right] + \xi_{j,t}$.
Applying the Sherman-Morrison formula, we obtain
\begin{equation}
    c_t = [I - (1-\psi^{-1}) \textbf{1}_{J+1} x_t^\top] \zeta_t,
\end{equation}
or $c_{j,t} = \zeta_{j,t} - (1-\psi^{-1}) x_t^\top \zeta_t$.
Notice that $\zeta_t$ can be computed at order $n$ based on the coefficients of order $n-1$ and their derivatives.

\paragraph{Computing the derivatives.} The derivation above shows that, given the order $n-1$ expansion of $\zeta_j(X,\epsilon)$ and its derivatives, we can compute the expansion of order $n$.
Suppose the expansion of $\zeta_{j}(X,\epsilon)$ is given by
\begin{equation}
    \zeta_j(X,\epsilon) = \sum_{k=0}^{n-1} \zeta_{j,k}(X) \epsilon^k,
\end{equation}
where $\zeta_{j,k}(X)$ takes the form:
\begin{equation}
    \zeta_{j,k}(X) = A_{j,k} + B_{j,k}^\top (X-\overline{X}) + \frac{1}{2} (X- \overline{X})^\top C_{j,k} (X - \overline{X}),
\end{equation}
where $\overline{X}$ is a reference point, $A_{j,k}$ is a scalar, $B_{j,k}$ is a vector, and $C_{j,k}$ is a matrix.
Notice that .
$\zeta_{j,k}(\overline{X}) = A_{j,k}$, $\zeta_{j,k,X}(\overline{X}) = B_{j,k}$ and $\zeta_{j,k,XX} = C_{j,k}$.
Given this expansion, we can compute $c_j(X,\epsilon) = \sum_{k=0}^{n-1} c_{j,k}(X) \epsilon^k$.

\subsection{Inner region}

Consider the case of no preference heterogeneity and no leverage constraints.
Consider the following change of variables: $x_1 = \epsilon \tilde{x}_1$.
Define $\tilde{c}_j( \tilde{x}_1, \overline{\alpha}_p ) = c_j\left(x_1, \overline{\alpha}_p\right) $, so $c_{j,x_1}= \frac{1}{\epsilon} \tilde{c}_{j,\tilde{x}_1}  $ and $ c_{j,x_1 x_1}= \frac{1}{\epsilon^2} \tilde{c}_{j,\tilde{x}_1 \tilde{x}_1}$.
This implies the following is true:
\begin{equation}
    \sigma_{c_j} =  \frac{1}{\epsilon}\frac{\tilde{c}_{j,\tilde{x}_1}}{\tilde{c}_j} \sigma_{x_1},
\end{equation}
where $\sigma_{x_1} = - \frac{1-\epsilon \tilde{x}_1}{\tilde{x}_1} \hat{\alpha}_p \sigma_R$ .
Similarly, we can write $\sigma_{y}$
\begin{equation}
    \sigma_y = \frac{1}{\epsilon} \frac{\tilde{y}_{x_1}}{y} \sigma_{x_1}.
\end{equation}

The drift of $y$ is given by
\begin{equation}
    \mu_{y} = \frac{1}{\epsilon} \frac{\tilde{y}_{\tilde{x}_1}}{y} \mu_{\tilde{x}_1} + \frac{1}{2 \epsilon^2}  \frac{\tilde{y}_{\tilde{x}_1 \tilde{x}_1}}{y} \sigma_{\tilde{x}_1}^2.
\end{equation}

The term of order 0 is the same as before.

\section{Volatility}

Consider the case without preference heterogeneity or leverage constraints.
The price-dividend ratio is given by:
\begin{equation}
    \hat{p}_t = - (1-\psi^{-1})\frac{\gamma\|\sigma\|^2}{2 y^*}\,
                    \frac{(1-x_{a,t}) (1-\overline{\alpha}_{p,t})^2}{x_{a,t}}.
\end{equation}
where, given $\overline{\alpha}_{a,t} \equiv \frac{1-(1-x_{a,t})  \overline{\alpha}_{p,t}}{x_{a,t}}$, we have
\begin{equation}
    \sigma_{x_a,t} = x_{a,t} (\overline{\alpha}_{a,t}-1) \sigma, \qquad \qquad 
    \sigma_{\overline{\alpha}_p,t} = \sigma_{\overline{\alpha}_p}
\end{equation}

The diffusion of $\hat{p}_t$ is given
\begin{equation}
    \sigma_{p,t} = \hat{p}_{x_a,t} \sigma_{x_a,t} + \hat{p}_{\overline{\alpha}_p,t} \sigma_{\overline{\alpha}_p}.
\end{equation}

\begin{equation}
    \hat c_t \;\equiv\; c_t - c^*
    \;=\; \chi_{0,t} \;+\; \chi_p\,\hat p_t \;+\; \mathcal{O}(\epsilon^3),
\end{equation}
with coefficients
\[
    \chi_p = (\psi-1)\,\frac{1}{p^*},
    \qquad
    \chi_{0,t} = (\psi-1)\,\frac{\gamma\|\sigma\|^2}{2}\,
                    \frac{f_t^2}{x_{0,t}(1-x_{0,t})}.
\]


\section{Flow of Funds aggregation}\label{app:fof_aggregation}

This appendix describes the procedure used in Section~\ref{sec:calibration} to aggregate Flow of Funds (FoF) sectors into the three ultimate sectors of the model: passive, constrained active, and unconstrained active.
We follow the aggregation in \citet{koijen2024investors}, drop sectors with no or minimal equity holdings, and route pass-through vehicles to their ultimate end investors.
ETFs (L.124) and closed-end funds (L.123) are aggregated onto the household balance sheet (L.101); mutual-fund holdings (L.122) are split across end investors using the FoF's breakdown, with the household portion kept separate.
Defined-contribution pensions (L.118c, L.119c, L.120c) are routed to households.
Hedge funds are separated out using Table~B.101.f for domestic funds and the Enhanced Financial Accounts for foreign funds, and the corresponding holdings are subtracted from the household (L.101) and rest-of-the-world (L.133) balance sheets.
The remaining insurance and pension sectors (L.115, L.116, L.118b, L.119b, L.120b) are aggregated into a single long-horizon sector.
The three ultimate sectors correspond, respectively, to \emph{passive} (households and non-profits net of hedge funds, government, monetary authority, and money-market funds), \emph{constrained} (hedge funds, broker-dealers, and other short-horizon intermediaries), and \emph{unconstrained} (insurance, pensions, and rest-of-the-world net of hedge funds).

\section{Numerical methodology}\label{app:numerical}
This appendix describes the numerical procedure used to solve the dynamic model in Section \ref{sec:model}.
We follow a two-step strategy where we first solve a `simplified' model with an aggregate claim in positive net supply and a derivative security in zero net supply.
In the second step (valuation step), we solve the valuation equation for the dividend and risky debt claims.
This ordering is justified since, once the equilibrium is computed in the first step, the market prices of risk and interest rates are known functions of the simplified model, and the individual valuation ratios can be recovered from the respective pricing  equations.

\subsection{Simplified economy}
The state variables in the simplified economy are
\[
X=(x_c,x_u,\alpha_p)\in\mathcal S,
\qquad x_p=1-x_c-x_u,
\]

where $x_{c}$ and $x_{u}$ are the wealth shares of the constrained and unconstrained active investors, and $\alpha_{p}$ is the exogenous risky-share demand of the passive sector.
The aggregate claim is in positive net supply and the derivative is in zero net supply.
Denote the active portfolio exposures by $(\alpha_{j,Y},\theta_j)$, where $\alpha_{j,Y}$ is the exposure to the aggregate claim and $\alpha_{j,F}$ is the exposure to the derivative.
The passive investor has exposure $(\alpha_{p,t},0)$.
We approximate
\[
{\xi}_c(X),\quad {\xi}_u(X), \quad {\xi}_p(X),\quad
{\alpha}_{c,Y}(X), \quad {\alpha}_{c,F}(X),\quad
{\alpha}_{u,Y}(X), \quad {\alpha}_{u,F}(X)
\]
with neural networks parameterized by $\Theta$:
\begin{align}
    \hat\xi_j &: \mathcal S\to\mathbb R_+, \qquad j\in\{c,u,p\}, \\
    \hat\alpha_{j,Y}&: \mathcal S\to\mathbb R, \qquad j\in\{c,u\}, \\
    \hat\alpha_{j,F} &: \mathcal S\to\mathbb R, \qquad j\in\{c,u\}.
\end{align}
and derive for all other equilibrium quantities as functions of the state variables and the approximated objects.
Note that the passive investor does not trade the derivative, so $\alpha_{p,F}(X)\equiv 0$, while the passive exposure to the aggregate claim is the state variable itself, $\alpha_{p,Y}(X)=\alpha_p$.
Thus the passive portfolio is imposed directly rather than approximated.
We use fully connected feed-forward neural networks with 8 hidden layers and 10 neurons in each layer.
Each neuron in the hidden and final layer is represented by a softplus activation function for the consumption-wealth ratios $\hat\xi_j$'s, and a linear function for the others.
The state space is sampled globally using Latin-hypercube points on the simplex
\[
\{(x_c,x_u,\alpha_p): x_c>0,\ x_u>0, \alpha_p >0, \ x_c+x_u<1\}
\]
Given a candidate collection of network outputs at $X$, the remaining equilibrium objects are recovered analytically.

First, goods-market clearing implies the aggregate claim dividend yield
\begin{equation}
\hat y(X)=x_c\hat\xi_c(X)+x_u\hat\xi_u(X)+x_p\hat\xi_p(X).
\end{equation}

Second, the volatility of the dividend yield is recovered from the It\^o consistency condition.
Third, once $\sigma_y(X)$ is known, the excess-return volatilities of the aggregate claim and the derivative, denoted $(\sigma_{R,Y}(X),\sigma_{R,F}(X))$, are recovered.
The unconstrained active investor's first-order conditions then pin down the two risk premia $\pi_Y(X), \pi_F(X)$.
Fourth, conditional on $(\hat y,\pi_Y,\pi_F,\sigma_y)$, the risk-free rate $r(X)$, the state drift $\mu_X(X)$, the state diffusion matrix $\Sigma_X(X)$, the consumption-wealth ratio drifts $\mu_{c_j}(X)$, and the market price of risk vector $\eta(X)$ are all recovered from the equilibrium identities and It\^o's lemma.
These calculations are explicit once the neural-network outputs and their first and second derivatives are known.
The problem reduces to minimizing the following loss functions

\begin{enumerate}
    \item \textbf{HJB residuals.} For each investor type $j\in\{c,u,p\}$, we construct the residual from the investor's Hamilton--Jacobi--Bellman equation after replacing all equilibrium objects by either neural-network outputs or the derived quantities defined above.
    We denote these residuals by
    \[
    \mathcal L^{\mathrm{HJB}}_j(X;\Theta).
    \]

    \item \textbf{Portfolio FOCs.} The constrained investor's first-order conditions for the aggregate claim and derivative lead to two residuals denoted as:
    \[
    \mathcal L^{\mathrm{FOC}}_{c,Y}(X;\Theta),
    \qquad
    \mathcal L^{\mathrm{FOC}}_{c,F}(X;\Theta).
    \]
    Note that the unconstrained investor's first-order conditions are used to recover $(\pi_Y,\pi_F)$ exactly and therefore do not appear separately in the loss.

    \item \textbf{Market-clearing residuals.} We impose clearing in the aggregate-claim and derivative markets as residuals:
    \[
    \mathcal L^{\mathrm{MC}}_Y(X;\Theta),
    \qquad
    \mathcal L^{\mathrm{MC}}_F(X;\Theta).
    \]
\end{enumerate}
The resulting loss is
\begin{equation}
    \label{eq:baseline_pointwise_loss_appendix}
    \mathcal L(X;\Theta)
    =
    \sum_{j\in\{c,u,p\}}\bigl(\mathcal L^{\mathrm{HJB}}_j(X;\Theta)\bigr)^2
    +\bigl(\mathcal L^{\mathrm{FOC}}_{c,Y}(X;\Theta)\bigr)^2
    +\bigl(\mathcal L^{\mathrm{FOC}}_{c,F}(X;\Theta)\bigr)^2
    +\bigl(\mathcal L^{\mathrm{MC}}_{Y}(X;\Theta)\bigr)^2
    +\bigl(\mathcal L^{\mathrm{MC}}_{F}(X;\Theta)\bigr)^2.
\end{equation}
For a training batch $\{X^n\}_{n=1}^N$, the training objective is the empirical mean
\begin{equation}
\label{eq:baseline_batch_loss_appendix}
\widehat{\mathcal L}(\Theta)
=\frac{1}{N}\sum_{n=1}^N \mathcal L(X^n;\Theta).
\end{equation}
The full algorithm is given in Algorithm (\ref{alg:first_step}).



\begin{algorithm}[!ht]
\setstretch{1.08}
\begin{algorithmic}[1]
\STATE Initialize NNs
$\{\hat\xi_c,\hat\xi_u,\hat\xi_p,\hat\alpha_{c,Y},\hat\alpha_{c,F},\hat\alpha_{u,Y},\hat\alpha_{u,F}\}$
with parameters $\Theta$.
\FOR{$t=1,2,\dots,T_{\max}$}
 \STATE Sample a batch of aggregate states $\{X^n=(x_c^n,x_u^n,\alpha_p^n)\}_{n=1}^N$.
 \STATE Evaluate the seven network outputs at each sampled state.
 \STATE Recover the derived equilibrium objects state-by-state:
 \begin{enumerate}[label*=(\alph*)]
 \item compute $\hat y(X^n)$ from goods-market clearing;
 \item compute $\sigma_y(X^n)$ using It\^o's lemma;
 \item recover $(\sigma_{R,Y}(X^n),\sigma_{R,F}(X^n))$;
 \item compute the hedging terms from the gradients of $(\hat\xi_c,\hat\xi_u,\hat\xi_p)$;
 \item recover $(\pi_Y(X^n),\pi_F(X^n))$ from the unconstrained investor's first-order conditions;
 \item get $(r(X^n),\mu_X(X^n),\Sigma_X(X^n),\mu_{\xi_j}(X^n))$ from equilibrium identities and It\^o's lemma.
\end{enumerate}
 \STATE Construct the residual blocks
 $\mathcal L^{\mathrm{HJB}}_j$, $\mathcal L^{\mathrm{FOC}}_{c,Y}$, $\mathcal L^{\mathrm{FOC}}_{c,F}$, $\mathcal L^{\mathrm{MC}}_Y$, and $\mathcal L^{\mathrm{MC}}_F$.
 \STATE Form the batch loss $\widehat{\mathcal L}(\Theta)$ in \eqref{eq:baseline_batch_loss_appendix}.
 \STATE Update $\Theta$ using ADAM, and update learning-rate scheduler using the current training loss.
\ENDFOR
\end{algorithmic}
\caption{First-stage training.}
\label{alg:first_step}
\end{algorithm}



\subsection{Valuation ratios}
After solving the simplified economy, we price the dividend and risky-debt claims as a function of the state variables $(X,s)$, where $s\in[0,1]$ denote the dividend share of the aggregate claim.
We solve for the valuation ratio of the dividend claim, denoted by $p_D(X,s)$, and the valuation ratio of risky-debt claim, denoted by $p_B(X,s)$.
We approximate these functions by two additional neural networks parameterized by $\Theta_D$ and $\Theta_B$, respectively:
\begin{align}
\hat p_D &: \mathcal S\times[0,1]\to\mathbb R_+, \\
\hat p_B&: \mathcal S\times[0,1]\to\mathbb R_+,
\end{align}
At this stage the first-stage simplified equilibrium is held fixed.
For each state $(X,s)$, the first-stage solution supplies the coefficients entering the claim-pricing PDEs: the risk-free rate $r(X)$, the market price of risk $\eta(X)$, the state drift $\mu_X(X)$, the state diffusion matrix $\Sigma_X(X)$, and the aggregate-claim return volatility.
The remaining coefficients associated with the $s$-state are known analytically from the decomposition of the aggregate claim into dividend and risky-debt components.
Given a candidate valuation ratio $\hat p_k(X,s)$ for claim $k\in\{D,B\}$, we compute its first and second derivatives with respect to $(X,s)$ using automatic differentiation, leading to the residual
\[
\mathcal L^{\mathrm{val}}_k(X,s;\Theta_k),
\qquad k\in\{D,B\}.
\]
The training objective is simply the empirical mean squared PDE residual,
\begin{equation}
\widehat{\mathcal L}^{\mathrm{val}}_k(\Theta_k)
=\frac{1}{N}\sum_{n=1}^N
\bigl(\mathcal L^{\mathrm{val}}_k(X^n,s^n;\Theta_k)\bigr)^2,
\qquad k\in\{D,B\}.
\end{equation}

The full algorithm is given in Algorithm (\ref{alg:valuation_step}).


\begin{algorithm}[h!]
 \setstretch{1.08}
 \begin{algorithmic}[1]
 \STATE Load the trained first-stage networks.
 \STATE Initialize the valuation-ratio network $\hat p_k$ for claim $k\in\{D,B\}$ with parameters $\Theta_k$.
\FOR{$t=1,2,\dots,T_{\max}^{\mathrm{val}}$}
 \STATE Sample a batch of states $\{(X^n,s^n)\}_{n=1}^N$.
 \STATE Evaluate $\hat p_k(X^n,s^n)$ and compute its first and second derivatives with respect to $(X,s)$.
 \STATE Recover the drift and volatility of the valuation ratio from It\^o's lemma.
 \STATE Use the fixed baseline solution to evaluate the coefficients entering the pricing PDE.
 \STATE Construct the claim-pricing residual $\mathcal L_k^{\mathrm{val}}(X^n,s^n;\Theta_k)$.
 \STATE Compute batch loss $\widehat{\mathcal L}^{\mathrm{val}}_k(\Theta_k)$.
 \STATE Update $\Theta_k$ using ADAM step, update the learning-rate scheduler using current training loss.
 \ENDFOR
\end{algorithmic}
 \caption{Valuation-ratio step.}
 \label{alg:valuation_step}
\end{algorithm}

We follow \cite{Gopalakrishna2026} and use vectorized automatic differentiation to compute the gradients that speeds up the computation.
The total time to run the baseline model on a single-threaded computer with 128GB RAM and NVIDIA Ada 4500 GPU is 8 hours.
The HJB loss in the first-stage training problem converges to a mean-squared loss of $3.7 \times 10^{-8}$.
The corresponding valuation-ratio training problems converge to a normalized mean-squared training loss of $1.4 \times 10^{-5}$.
These losses are computed from the batch objectives in \eqref{eq:baseline_batch_loss_appendix} and in the valuation-ratio objective, respectively.


% \section{The unit EIS case}

% Suppose $\psi = 1$ and $\gamma_j = 1$ for $j = 1,\ldots, J$. In this case, the consumption-wealth ratio satisfies $c_{j,t} = \rho$ and the dividend-yield is given by $y_t = \rho$. This implies that $\sigma_{R,t} = \sigma$. The risk premium is given by
% \begin{equation}
%     \pi_t = \left[1+\frac{x_{0,t} (1- \overline{\alpha}_{p,t}) }{1-x_{0,t}}\right] \|\sigma\|^2.
% \end{equation}

% The interest rate is given by 
% \begin{equation}
%     r_t = \rho + \mu -\left[1+\frac{x_{0,t} (1- \overline{\alpha}_{p,t}) }{1-x_{0,t}}\right] \|\sigma\|^2.
% \end{equation}

% The portfolio share of active investors is $\alpha_{j,t} = \frac{1-x_{0,t} \overline{\alpha}_{p,t}}{1-x_{0,t}}$.
% The drift and diffusion of the state variable are given by
% \begin{align}
%     \mu_{x,t} &= \left[ r_t + \pi_t \alpha_{j,t} - \rho - \mu + (1-\alpha_{j,t}) \|\sigma\|^2 \right] x_t + \kappa (\omega_j - x_t)\\
%     \sigma_{x,t} &= (\alpha_{j,t}-1) \sigma 
% \end{align}

% \newpage
% The HJB equation can be written as\small
% \begin{equation}
%     c_{j,t} = \psi \rho + (1-\psi) \left[ r_t + \pi_t \alpha_{j,t} - \frac{\gamma_j}{2} \| \sigma \|^2 \alpha_{j,t}^2 \right] + \frac{c_{j,x}}{c_j} \mu_x \epsilon + \frac{1}{2} \frac{c_{j,xx}}{c_j} \sigma_x^2 \epsilon^2 + (1-\gamma_j) \frac{c_{j,x}}{c_j} \sigma_x \epsilon \sigma^\top \alpha_{j,t} + \frac{\psi - \gamma_j}{1-\psi} \frac{1}{2} \left(\frac{c_{j,x}}{c_j} \sigma_x \epsilon\right)^2
% \end{equation}\normalsize

% The zeroth-order term is given by
% \begin{equation}
%     c_{j,t}^{(0)} = \psi \rho + (1-\psi) \left[ r_t + \pi_t \alpha_{j,t} - \frac{\gamma_j}{2} \| \sigma \|^2 \alpha_{j,t}^2 \right].
% \end{equation}

% The first-order term is given by
% \begin{equation}
%     c_{j,t}^{(1)} = \frac{c_{j,x}^{(0)}}{c_j^{(0)}} \mu_x  + (1-\gamma_j) \frac{c_{j,x}^{(0)}}{c_j^{(0)}} \sigma_x \sigma^\top \alpha_{j,t}.
% \end{equation}

% % The second-order term is given by
% % \begin{equation}
% %     c_{j,t}^{(2)} = \hat{c}_{j,x}^{(1)} \mu_x  + \frac{1}{2}\hat{c}_{j,xx}^{(0)}\sigma_x^2 + (1-\gamma_j) \hat{c}_j^{(1)} \sigma_x  \sigma^\top \alpha_{j,t} + \frac{\psi - \gamma_j}{1-\psi} \frac{1}{2} \left(\hat{c}_{j,x}^{(0)} \sigma_x \right)^2
% % \end{equation}

% The term of order $k\geq 2$ is given by
% \begin{equation}
%     c_{j,t}^{(k)} = \hat{c}_{j,x}^{(k-1)} \mu_x  + \frac{1}{2}\hat{c}_{j,xx}^{(k-2)}\sigma_x^2 + (1-\gamma_j) \hat{c}_j^{(k-1)} \sigma_x  \sigma^\top \alpha_{j,t} + \frac{\psi - \gamma_j}{1-\psi} \frac{1}{2} \left(\hat{c}_{j,x}^{(k-2)} \sigma_x \right)^2.
% \end{equation}




